"\documentclass{article}

% NeurIPS 2026 style. Default (no options) = main-track, double-blind, anonymous submission.
% For camera-ready, switch to: \usepackage[main, final]{neurips_2026}
% For arXiv preprint,         switch to: \usepackage[preprint]{neurips_2026}
\usepackage{neurips_2026}
\usepackage[utf8]{inputenc}
\usepackage[T1]{fontenc}
\usepackage{hyperref}
\usepackage{url}
\usepackage{booktabs}
\usepackage{amsfonts}
\usepackage{xcolor}
\usepackage{nicefrac}
\usepackage{microtype}
\usepackage[table]{xcolor}
\usepackage{graphicx}
\usepackage{overpic}
\usepackage{multirow}
\usepackage{amsmath}
\usepackage{amssymb}
\usepackage{mathtools}
\usepackage{amsthm}
\usepackage{algorithm}
\usepackage{algpseudocode}
\newcommand{\theHalgorithm}{\arabic{algorithm}}
\usepackage[capitalize,noabbrev]{cleveref}
\usepackage{colortbl}
\theoremstyle{plain}
\newtheorem{theorem}{Theorem}[section]
\newtheorem{proposition}[theorem]{Proposition}
\newtheorem{lemma}[theorem]{Lemma}
\newtheorem{corollary}[theorem]{Corollary}
\theoremstyle{definition}
\newtheorem{definition}[theorem]{Definition}
\newtheorem{assumption}[theorem]{Assumption}
\theoremstyle{remark}
\newtheorem{remark}[theorem]{Remark}
\usepackage[textsize=tiny]{todonotes}
\usepackage{wrapfig}
\usepackage{subcaption}
\usepackage{tikz}
\usetikzlibrary{arrows.meta,positioning,shapes.geometric,calc,decorations.pathmorphing}

% Shortcuts
% Forward/inverse lift: time series <-> image (e.g., delay embedding or STFT).
% Matches the $\mathcal{T}$ / $\mathcal{T}^{-1}$ convention used in ImagenTime
% (Naiman et al., 2024). \tsimg and \imgts are kept only as internal macro
% names to minimise diffs; the rendered symbols are pure math.
\newcommand{\tsimg}{\mathcal{T}}
\newcommand{\imgts}{\mathcal{T}^{-1}}
% Framework and instances.
% PSLR = Posterior Sampling for Lifted Representations (our framework/recipe;
% parallels PSLD=Posterior Sampling with Latent Diffusion).
% L-MMPS = PSLR applied to MMPS (Rozet et al. 2024); the E-step of our method.
% L-TMPD = PSLR applied to TMPD (Boys et al. 2023); second instance (ongoing).
\newcommand{\pslr}{\textsc{pslr}}
\newcommand{\lmmps}{L\mbox{-}\textsc{mmps}}
\newcommand{\ltmpd}{L\mbox{-}\textsc{tmpd}}
\newcommand{\Nts}{N_{\mathrm{ts}}}
\newcommand{\Nimg}{N_{\mathrm{img}}}
\newcommand{\Nobs}{N_{\mathrm{obs}}}
\newcommand{\Ats}{A_{\mathrm{ts}}}
\newcommand{\Aimg}{A_{\mathrm{img}}}
\newcommand{\yts}{y_{\mathrm{ts}}}
\newcommand{\xts}{x_{\mathrm{ts}}}
\newcommand{\ximg}{x_{\mathrm{img}}}
\newcommand{\gal}[1]{\textcolor{blue}{[Gal: #1]}}

\title{Closing the Loop: Co-Evolving EM for \\ Irregular Time Series Generation in Lifted Representations}

\author{%
  Anonymous Authors \\
  Anonymous Affiliation \\
  \texttt{anonymous@email.edu} \\
}

\begin{document}

\maketitle

\begin{abstract}
Generating realistic time series is central to healthcare, finance, and the sciences, yet real-world data is \emph{irregular}: measurements are unevenly spaced and values are missing. Existing generators handle this \emph{open-loop}---a pretrained imputer fills in the gaps once, and the generator is trained on those frozen completions; the imputer and the generator never interact again. We show that \textbf{closing this loop}---continually refining the imputations against the generator's evolving prior---is the binding constraint on generative quality, and a far stronger lever than imputer sophistication. We realise the loop with Monte-Carlo EM: an E-step that draws imputations from the generator's current posterior, and an M-step that updates the generator on those imputations. Applying this idea to the current leading time-series generation method, which maps series to images before running diffusion, requires a new posterior-sampling recipe, because off-the-shelf samplers break down when the generative model and the observations live in different spaces. We introduce \textbf{\pslr{}} (\emph{Posterior Sampling for Lifted Representations}), a plug-and-play fix with three simple ingredients. Across diverse benchmarks, our method sets a new state of the art, recovering near-oracle performance with rapid convergence. The same recipe transfers without retuning to five distinct stress tests—ranging from block-missingness to continuous-time resampling—confirming that the closed-loop principle is a universal lever for generative quality, agnostic to the specific choice of imputer, lift, or observation operator.

\end{abstract}

\begin{figure}[h]
\centering
\resizebox{0.85\textwidth}{!}{
\begin{tikzpicture}[
  node distance=6mm and 10mm,
  box/.style={draw, rounded corners=2pt, minimum height=9mm, minimum width=22mm, align=center, font=\small},
  data/.style={box, fill=gray!10},
  imp/.style={box, fill=orange!20},
  gen/.style={box, fill=blue!15},
  arrow/.style={-Latex, thick},
  dashed arrow/.style={-Latex, thick, dashed, red!70!black}
]
\node[font=\bfseries] (titA) {(a) Open loop: commit-and-forget (prior work)};
\node[data, below=2mm of titA.west, anchor=west] (dataA) {Irregular\\ data $\yts$};
\node[imp, right=of dataA] (impA) {Imputer\\ (frozen)};
\node[gen, right=of impA] (genA) {Generator\\ $\theta$};
\node[right=8mm of genA, font=\small] (outA) {synthetic\\ series};
\draw[arrow] (dataA) -- (impA);
\draw[arrow] (impA) -- node[above, font=\scriptsize]{one-shot $\hat{x}$} (genA);
\draw[arrow] (genA) -- (outA);
\node[below=1mm of impA, font=\scriptsize, red!70!black] {\emph{no feedback}};

\node[font=\bfseries, below=16mm of titA.west, anchor=west] (titB) {(b) Closed loop: co-evolving EM (ours)};
\node[data, below=2mm of titB.west, anchor=west] (dataB) {Irregular\\ data $\yts$};
\node[imp, right=of dataB] (impB) {E-step:\\ posterior sampling};
\node[gen, right=of impB] (genB) {M-step:\\ train $\theta_{k+1}$};
\node[right=8mm of genB, font=\small] (outB) {synthetic\\ series};
\draw[arrow] (dataB) -- (impB);
\draw[arrow] (impB) -- node[above, font=\scriptsize]{posterior samples} (genB);
\draw[arrow] (genB) -- (outB);
\draw[dashed arrow] (genB.south) to[out=-90, in=-90, looseness=1.3] node[below, font=\scriptsize, red!70!black]{updated prior $\theta_k$} (impB.south);
\end{tikzpicture}
}
\caption{\textbf{Thesis.} Prior methods (top) are \emph{open-loop}: a fixed missingness-handling mechanism hands completions to a generator that never provides feedback. Our co-evolving EM (bottom) \emph{closes the loop}---each M-step updates the prior the next E-step samples from---and iteration beats imputer sophistication.}
\label{fig:teaser}
\end{figure}

\vspace{-3mm}
\section{Introduction}
\label{sec:intro}
\vspace{-1mm}

Time series underpin forecasting, anomaly detection, and pattern analysis across healthcare, finance, and the sciences~\citep{han2019review,ismail2019deep,lim2021time}, and \emph{generating} synthetic series is increasingly valuable for privacy-preserving proxies, stress testing, and dataset balancing~\citep{brophy2023generative,nochumsohn2025data}. In practice, however, real-world series are \emph{irregular}---unevenly spaced, with values missing due to sensor failures, inconsistent sampling, or interruptions in monitoring~\citep{kidger2020ncde,ren2024learning}. The central challenge we address is learning to generate the underlying, regularly-sampled distribution $p(\mathbf{x})$ from a corpus where the observations are perpetually incomplete and unevenly spaced.

Prior methods for irregular time-series generation fall into two families of solutions, both of which follow an \emph{open-loop} scheme. One family handles irregularity \emph{architecturally} by augmenting recurrent backbones with inter-observation intervals (e.g., $\Delta t$) or hidden-state decay mechanisms \citep{che2018recurrent, jeon2022gtgan}. The other adopts a commit-and-forget pipeline, running a pretrained imputer once and freezing those completions for generator training—a category including GT-GAN \citep{jeon2022gtgan}, KoVAE \citep{naiman2024kovae}, and the current SOTA, ImagenI2R \citep{fadlon2025imageni2r}. Both share the same structural defect: no feedback from the generator back to the mechanism that produces completions. Consequently, they face a lose-lose choice. The use of a deterministic imputer biases the generator toward the conditional mean, $\mathbb{E}[\mathbf{x} | \mathbf{y}]$. For multimodal data, this mean acts as an arithmetic average that sits in the \emph{empty space} between distinct modes; it describes a state that never actually occurs in the real world (\cref{sec:why-close}). Such methods fail to recover the conditional distribution, $p(\mathbf{x} | \mathbf{y})$, necessary for high-fidelity synthesis. Conversely, recent methods like ImagenI2R attempt to avoid this trap via masked training, yet this introduces a masking--data trade-off: it preserves variance only by discarding $(1-\rho)$ of the data at missing rate $\rho$, creating a fundamental limit on data efficiency. Crucially, we find that minimizing imputer reconstruction error can be counter-productive; an imputer that converges to the conditional mean may achieve lower MSE while simultaneously suppressing the variance necessary for high-fidelity generation (\cref{tab:toy-mse-trap}). \textbf{Iteration, not imputer sophistication or masking heuristics, is the binding constraint.}

We close the imputer--generator loop via co-evolving Monte-Carlo EM~\citep{dempster1977em,wei1990mcem,neal1998em}: the E-step draws \emph{samples} (not point estimates) of the missing values from the generator's current posterior, the M-step retrains the generator on those samples, and the two are iterated to reach a self-consistent fixed point. MCEM is the right closed-loop scheme for this problem for three reasons, all formalised in \cref{sec:why-close,subsec:theorem}. First, it has a monotone improvement guarantee on the data likelihood, whereas adversarial closed loops can oscillate without ever converging. Second, once the E-step has sampled, the data is complete, so the M-step is vanilla score-matching on full windows---no masked loss, no differentiable imputer, no encoder---precisely what end-to-end joint training would force us to give up. Third, diffusion posterior sampling lets us draw from the generator's true posterior distribution, $p_\theta(\mathbf{x} | \mathbf{y})$, rather than relying on a fixed approximate encoder that would itself put a structural ceiling on accuracy.

Applying this idea to the state-of-the-art time-series generator~\citep{naiman2024imagentime} is not immediate, because the model does not diffuse on the series directly. Instead, it employs an invertible time-series-to-image transformation—such as a delay embedding or a spectrogram—to map the data into a high-dimensional lifted representation. This allows the model to leverage modern vision-diffusion backbones~\citep{karras2022edm}, but it creates a structural hurdle for the E-step: the generative prior resides in the lifted image space, while the irregular observations exist in the original time-series space. Existing diffusion posterior samplers are derived for the single-space case where these domains coincide; we show that they fail in this dual-space regime, either collapsing to random chance or stagnating far below the oracle. A posterior sampler purpose-built for the lifted setting is therefore required.

We distil three conditions that any posterior sampler must satisfy to be a valid E-step for a lifted prior (\cref{subsec:requirements}), and use them to separate existing samplers into two groups: those that violate one of the conditions \emph{by design}, and those that only violate the remaining two because they were derived for the single-space case. For the second group we introduce \textbf{Posterior Sampling for Lifted Representations} (\pslr{}), a plug-and-play recipe with three ingredients: \emph{(i)}~use the true composed observation operator instead of the assumed one, \emph{(ii)}~solve the inner linear system in the observation (time-series) space rather than the lifted image space, and \emph{(iii)}~project each posterior draw back onto the space of valid lifted images. Applying \pslr{} uniformly to the two correctable samplers in the literature - MMPS \citep{rozet2024diem} and TMPD \citep{boys2023tmpd}-yields our two instances, $\lmmps$ and $\ltmpd$, which inherit the same recipe and differ only in how they approximate the posterior covariance. A convergence analysis (\cref{thm:operator-mismatch}) explains \emph{why} each ingredient is needed: without them, MCEM either diverges or settles at a biased fixed point; with all three, it contracts to the correct generator.

\vspace{-1mm}
\paragraph{Contributions.}
\vspace{-1mm}
\begin{enumerate}
\setlength{\itemsep}{2pt}

\item \textbf{Iteration is the binding constraint.} We show, theoretically (\cref{sec:why-close}) and across the $9$-dataset benchmark, that co-evolving EM beats every published open-loop baseline, including those using sophisticated architectural priors or masking heuristics. We demonstrate that iteration—not imputer sophistication—is the primary driver of generative fidelity.

\item \textbf{\pslr{}: a plug-and-play recipe for lifted posterior sampling.} We identify three conditions any sampler must satisfy for a lifted prior and package the fix as a three-ingredient recipe with a convergence guarantee (\cref{thm:operator-mismatch}), yielding two instances ($\lmmps$, $\ltmpd$) that transfer without retuning across lifts and observation models (\cref{sec:continuous-operator}).

\item \textbf{State-of-the-art on benchmarks and stress tests.} Our method sets a new state of the art across the standard $9$-dataset benchmark, recovering near-oracle performance with rapid convergence. We achieve SOTA on five distinct stress tests—ranging from block-missingness to noisy observations—that defeat standard one-shot-imputer baselines (\cref{tab:results_short_l24,tab:results_long}).

\end{enumerate}

\vspace{-2mm}
\section{Related Work}
\label{sec:related}
\vspace{-1mm}

\paragraph{Diffusion models.}
Denoising diffusion~\citep{ho2020ddpm} and its continuous-time score-based formulation~\citep{song2021sde} are the dominant paradigm for high-fidelity generative modelling, particularly on images, with EDM~\citep{karras2022edm} providing the training and sampling recipe we adopt.

\paragraph{Generative modelling of time series.}
Early work adopted GANs~\citep{yoon2019timegan} and VAEs for time-series synthesis, and diffusion was later brought to time series directly via conditional score-based models~\citep{tashiro2021csdi}. More recently, ImagenTime~\citep{naiman2024imagentime} showed that mapping each series to an image via an invertible time-series--to-image transformation and applying image-space diffusion outperforms direct time-series diffusion by large margins on both short and ultra-long benchmarks, establishing the current SOTA for \emph{regularly-sampled} sequences. This is the backbone our method trains end-to-end on irregular data.

\paragraph{Irregular time-series generation.}
Generative modeling for irregular time series has historically branched into two open-loop strategies. The first, represented by the TimeGAN-$\Delta t$/-D suite~\citep{jeon2022gtgan,che2018recurrent}, incorporates irregularity into the model backbone via fixed RNN primitives. The second lineage employs a commit-and-forget pipeline where a one-shot imputer provides static completions for a downstream generator. While GT-GAN~\citep{jeon2022gtgan} and KoVAE~\citep{naiman2024kovae} popularized this via NCDE-based imputation~\citep{kidger2020ncde}, they are constrained by the computational overhead of continuous-time solving on long sequences and a "mean trap"—a phenomenon where training on deterministic imputations forces the generator to converge to the conditional mean $\mathbb{E}[\mathbf{x}|\mathbf{y}]$, collapsing the variance necessary to represent multi-modal distributions. The current SOTA, ImagenI2R \citep{fadlon2025imageni2r}, seeks to bypass the issues of fixed imputations by adopting high-capacity lifted diffusion backbones and employing masked training. While this sidesteps the "mean trap" by avoiding reliance on a pre-trained imputer, it introduces a severe masking–data trade-off: it preserves variance only by discarding a large fraction of the available data, creating a fundamental limit on data efficiency. Our work bridges these lineages: we leverage the efficient lifted priors of recent diffusion models but replace the static imputer with a co-evolving MCEM loop, utilizing the full observational signal without the data-discarding trade-offs of masking heuristics (see \cref{tab:intro-taxonomy} for a side-by-side taxonomy).

\paragraph{Posterior sampling with diffusion priors.}
Our E-step builds on the literature of diffusion posterior samplers, which approximate the likelihood score $\nabla_{\mathbf{x}_t} \log p(\mathbf{y}\mid \mathbf{x}_t)$ under a single-space assumption where the observation operator acts directly in the model's native space. Existing families differ primarily in their approximation of the conditional $p(\mathbf{x}\mid \mathbf{x}_t)$ or its associated covariance: DPS~\citep{chung2023dps} uses a point estimate; $\Pi$GDM~\citep{song2023pigdm} a diagonal Gaussian with heuristic variance; DiffPIR~\citep{zhu2023diffpir} a proximal update; TMPD~\citep{boys2023tmpd} a row-sum diagonal approximation; and MMPS~\citep{rozet2024diem} the full Jacobian covariance via conjugate gradient. We introduce \textbf{\pslr{}}, a recipe that identifies and corrects the dual-space misspecifications that arise when the generative prior lives in a lifted representation. While PSLD~\citep{rout2024psld} addresses dual-space sampling for learned latent-diffusion maps, \pslr{} provides a fix for known, structured lifts (e.g., delay embeddings), yielding our two instances $\ltmpd$ and $\lmmps$ via the same structural corrections.

\vspace{-2mm}
\section{Background}
\label{sec:background}
\vspace{-1mm}

\paragraph{Problem statement.}
We observe a corpus of irregular time series $\{(\yts^{(i)}, S_{\mathrm{ts}}^{(i)})\}_{i=1}^N$, where $\yts^{(i)} = S_{\mathrm{ts}}^{(i)} \cdot \xts^{(i)}$ and $S_{\mathrm{ts}}^{(i)} \in \mathbb{R}^{\Nobs \times \Nts}$ is a linear observation operator. We evaluate on two operator instances: (i) \textbf{binary MCAR masking}, the standard benchmark setting; and (ii) \textbf{continuous-time irregular resampling} via fractional-support interpolation (\S\ref{sec:continuous-operator}). The goal is to learn the clean marginal distribution $p(\xts)$ and sample complete, regular series from it.

\paragraph{Lifted representation.}
We adopt an invertible lift $\tsimg : \mathbb{R}^{\Nts \times F} \to \mathbb{R}^{\Nimg}$ and its pseudo-inverse $\imgts$. In practice, we utilize structured linear lifts such as \textbf{delay embedding}~\citep{takens1981embedding} and \textbf{STFT}~\citep{griffin1984signal}. These lifts are typically over-complete ($\Nimg > \Nts \times F$), causing overlapping windows to represent each time point in multiple image coordinates. Consequently, valid images reside in a linear subspace $\mathrm{Range}(\tsimg) \subset \mathbb{R}^{\Nimg}$, and the composition $\Pi = \tsimg \circ \imgts$ defines the orthogonal projection onto this manifold.

\paragraph{Diffusion models.}
We follow the EDM framework~\citep{karras2022edm} with denoiser $d_\theta(x_t, \sigma)$ and objective:
\begin{equation}
\arg\min_\theta \mathbb{E}_{p(x) p(\sigma) p(x_t \mid x)} \bigl[ \lambda(\sigma) \| d_\theta(x_t, \sigma) - x \|^2 \bigr].
\end{equation}
The optimal denoiser approximates the Tweedie posterior mean $\mathbb{E}[x \mid x_t]$, which is linked to the score via $\nabla_{x_t} \log p(x_t) = (d_\theta(x_t, \sigma) - x_t)/\sigma_t^2$. Crucially, the Tweedie covariance formula~\citep{efron2011tweedie} yields $\mathbb{V}[x \mid x_t] \approx \sigma_t^2 \nabla_{x_t} d_\theta(x_t, \sigma)$, which is the covariance object approximated by the posterior samplers we adopt.

\paragraph{EM for diffusion priors.}
MCEM alternates: (i) \textbf{E-step}---drawing posterior samples $\ximg \sim q_{\theta_k}(\ximg \mid \yts)$ under the current model; (ii) \textbf{M-step}---updating $\theta$ via denoising score matching on the E-step completions. This maximizes the marginal-likelihood $\mathbb{E}_{p(y)}[\log q_\theta(y)]$, with incremental MCEM~\citep{neal1998em} providing a monotone improvement guarantee up to Monte-Carlo noise.

\paragraph{Dual-space setting.}
In a lifted generator, the observation model crosses two domains:
\begin{equation}
\yts = S_{\mathrm{ts}} \cdot \imgts(\ximg), \qquad G \;\triangleq\; S_{\mathrm{ts}} \circ \imgts : \mathbb{R}^{\Nimg} \to \mathbb{R}^{\Nobs \times F}.
\end{equation}
Three structural properties of $G$ are central: it is a \textbf{composed cross-space operator} (not a native image mask), it is \textbf{rank-deficient} due to over-completeness (creating off-manifold null directions), and $G^\top G$ is \textbf{non-diagonal} as overlapping windows induce structured correlations. These properties guide the C1--C3 conditions in \S\ref{subsec:requirements}.

\vspace{-2mm}
\section{Method}
\label{sec:method}
\vspace{-1mm}

\subsection{Problem setup and close-the-loop objective}
\label{sec:why-close}
\vspace{-1mm}

We train a single co-evolving system through three stages: warm-start completion, a posterior-sampling E-step, and an EDM-based M-step. Unlike open-loop ``commit-and-forget'' pipelines, EM updates the generator $\theta$ and the implicit imputer (the posterior $q$) jointly. We optimize the marginal log-likelihood of the observations $\yts$ via the energy decomposition:
\begin{equation}
\label{eq:free-energy}
\log p_\theta(\yts) = \underbrace{\mathbb{E}_{q(\xts^{\mathrm{miss}})} \left[ \log p_\theta(\yts, \xts^{\mathrm{miss}}) \right] + H(q)}_{\mathcal{F}(q, \theta)} + \mathrm{KL} \left( q(\xts^{\mathrm{miss}}) \, \big\| \, p_\theta(\xts^{\mathrm{miss}} \mid \yts) \right)
\end{equation}
At each iteration, the E-step sets $q$ to the current model posterior $p_{\theta_k}(\xts^{\mathrm{miss}}\mid\yts)$ via \pslr{}, causing the residual KL term to vanish. On entry to the M-step, the surrogate objective $\mathcal{F}$ is exactly equal to the true marginal likelihood, and its maximization reduces to vanilla score matching on complete windows. This closed loop ensures the completion mechanism is repeatedly corrected by the generative prior it trains, driving the system toward a self-consistent fixed point that avoids the conditional-mean traps and metric misalignment common in open-loop baselines (see App.~\ref{app:why-close-full} for a detailed analysis).

\begin{figure*}
    \centering
    \begin{overpic}[width=\linewidth]{figs/arch.pdf}
        \put(-1,35.5) {\rotatebox{90}{Init}}
        \put(-1,24.5) {\rotatebox{90}{E-Step}}
        \put(-1,13.5) {\rotatebox{90}{M-Step}}
        \put(-1,2.5) {\rotatebox{90}{Inference}}
        \put(2, 40.5) {\small irregular signal}
        \put(25, 40.5) {\small STL-based imputation}
        \put(43, 40.5) {\small regular signal}
        \put(61, 40.5) {$\mathcal{T}$}
        \put(68, 40.5) {\small image}
        \put(14, 29.5) {\small noise}
        \put(36, 29.5) {\small U-Net}
        \put(54, 29.5) {\small imputed observation image}
        \put(2, 18.5) {\small image}
        \put(18, 18.5) {\small noisy image}
        \put(36, 18.5) {\small U-Net}
        \put(54, 18.5) {\small image}
        \put(3, 7.5) {\small noise}
        \put(20, 7.5) {\small U-Net}
        \put(39, 7.5) {\small image}
        \put(52, 4.5) {$\mathcal{T}^{-1}$}
        \put(61, 7.5) {\small regular signal}
    \end{overpic}
    \vspace{-3mm}
    \caption{TODO}
    \label{fig:arch}
    \vspace{-5mm}
\end{figure*}

The primary \emph{structural hurdle} resides in the E-step: the requirement to perform posterior sampling across \emph{disjoint representation spaces}. While the generative prior lives in the lifted image space, the irregular observations exist in the original time-series space. Off-the-shelf single-space samplers fail in this \emph{dual-space regime} because they cannot account for the manifold constraints and structured correlations inherent in the lift. We therefore define \textbf{\pslr{}}, a sampler-agnostic \emph{principled recipe} that identifies and corrects the \emph{cross-domain misspecifications} by enforcing operator exactness, stable conditioning, and manifold consistency.

The posterior-score identity that \pslr{} restores and makes operationally stable is:
\begin{equation}
\label{eq:pslr-score}
\nabla_{x_t} \log q(\yts \mid x_t) \;=\; J^\top G^\top \underbrace{\bigl(\sigma_y^2 I + \sigma_t^2 \, G J^\top G^\top\bigr)^{-1}}_{\text{M2: Obs-Space Precision}} \underbrace{\bigl(\yts - G \cdot d_\theta(x_t, \sigma)\bigr)}_{\text{M1: True Residual}},
\end{equation}
where $G = S_{\mathrm{ts}} \circ \imgts$ is the true composed observation operator and $J = \nabla_{x_t} d_\theta$ is the denoiser Jacobian. Our contribution is not merely a new posterior identity, but the \emph{structural framework} that enables iterative EM to converge in the presence of lifted representations.

\vspace{-1mm}
\subsection{What does a posterior sampler need, in a lifted representation?}
\label{subsec:requirements}
\vspace{-1mm}

We first define what any posterior sampler must satisfy to be a valid E-step under a linear lift.

\paragraph{Three conditions}
Let $\tsimg : \mathbb{R}^{\Nts \times F} \to \mathbb{R}^{\Nimg}$ be the lift, $\imgts$ its pseudo-inverse, $S_{\mathrm{ts}}$ the time-series sampling operator, and $G = S_{\mathrm{ts}} \circ \imgts$ the composed operator. A posterior sampler is a valid E-step for co-evolving EM in this setting iff it satisfies:
\begin{description}
\setlength{\itemsep}{1pt}
\item[C1 (composed operator).] The likelihood gradient is evaluated with $G = S_{\mathrm{ts}} \circ \imgts$, \emph{not} with a naïve image-space mask $\mathbf{M}_{\mathrm{img}} = \tsimg(S_{\mathrm{ts}})$. Using $\mathbf{M}_{\mathrm{img}}$ treats each of the $d$ image pixels that witness a single time-series observation as an independent measurement; this over-counts evidence by a factor $\sim d$ on observed windows and injects zero evidence on null-space directions, both of which bias the posterior.
\item[C2 (Tweedie covariance preservation).] The posterior score correction uses the full denoiser-Jacobian covariance $\sigma_t^2 G J^\top G^\top$ in the likelihood precision $\sigma_y^2 I + \sigma_t^2 G J^\top G^\top$. Diagonalizing or point-estimating this covariance collapses posterior diversity, which the downstream M-step then bakes into the prior.
\item[C3 (manifold consistency).] The returned sample lies in $\mathrm{Range}(\tsimg) \subset \mathbb{R}^{\Nimg}$; off-manifold components have no observation witness and no self-correcting mechanism, so they compound across EM iterations.
\end{description}

These conditions are sampler-agnostic: any method that satisfies C1--C3 is valid in our lifted EM loop.

\begin{table}[t]
\centering
\caption{Sampler compatibility in lifted space. C2-incompatible families are not repairable by construction; C2-compatible Tweedie-moment samplers are repairable by \pslr{}.}
\label{tab:samplers-motivation}
\resizebox{0.92\linewidth}{!}{
\begin{tabular}{l|l|l|l}
\toprule
\textbf{Family} & \textbf{Representative} & \textbf{C2?} & \textbf{Lifted E-step status} \\
\midrule
DPS / PiGDM / DiffPIR & \citep{chung2023dps,song2023pigdm,zhu2023diffpir} & No & intrinsically incompatible \\
MMPS (naive) & \citep{rozet2024diem} & Yes & violates C1/C3 under lift \\
TMPD (naive) & \citep{boys2023tmpd} & Yes & violates C1/C3 under lift \\
\textbf{\pslr{}(MMPS), \pslr{}(TMPD)} & \textbf{this work} & \textbf{Yes} & \textbf{valid via C1+C2+C3} \\
\bottomrule
\end{tabular}
}
\vspace{-3mm}
\end{table}

The consequence is direct: \pslr{} promotes C2-compatible samplers to valid lifted E-steps by enforcing C1--C3.

\vspace{-1mm}
\subsection{Three structural misspecifications of single-space samplers in a lifted representation}
\label{subsec:misspec}
\vspace{-1mm}

With C1--C3 in hand, failure modes are easy to diagnose. When a C2-compatible single-space sampler is dropped into lifted EM unchanged (MMPS or TMPD), it violates C1 and C3 and destabilizes the C2 solve:
\begin{itemize}
\setlength{\itemsep}{0pt}
\item \textbf{M1 (wrong operator).} Using $\Aimg$ over-counts observations and solves the wrong inverse problem.
\item \textbf{M2 (poor conditioning).} Fixed $\sigma_y$ makes early-step solves unstable as diffusion noise grows.
\item \textbf{M3 (off-manifold drift).} Samples leave $\mathrm{Range}(\tsimg)$ and accumulate error across EM iterations.
\end{itemize}
The root cause is a space mismatch: likelihood is evaluated as if observations were native image-space measurements, while they actually live in time-series space. \pslr{} resolves this by keeping diffusion in lifted space but evaluating likelihood with the true composed operator.

\vspace{-1mm}
\subsection{When can we close the loop? Operator mismatch determines convergence}
\label{subsec:theorem}
\vspace{-1mm}

Let $\Delta=\tilde{G}-G$ denote mismatch between the operator used by the E-step and the true composed operator. The theorem below gives a compact regime characterization for MCEM in lifted space.

\begin{theorem}[Operator mismatch in MCEM]
\label{thm:operator-mismatch}
Under linear-Gaussian regularity, EM dynamics with misspecified E-step operator $\tilde{G}$ satisfy:
\begin{enumerate}
\setlength{\itemsep}{0pt}
\item[(a)] $\Delta=0$: contracting dynamics and stable convergence.
\item[(b)] intermediate $\|\Delta\|$: temporary improvement followed by divergence.
\item[(c)] large $\|\Delta\|$: unstable or biased non-convergent dynamics.
\end{enumerate}
\end{theorem}

The full statement and proof are in App.~\ref{app:proof}. Main-text convergence trajectories (\cref{tab:convergence}) follow the same ordering, and the same regime pattern appears under STFT and continuous-time operators (App.~\ref{app:stft}, \ref{sec:continuous-operator}).

\begin{corollary}[E-step corrections as mismatch reduction]
\label{cor:corrections}
Each correction targets one mismatch component and pushes dynamics toward regime (a): operator substitution (M1), adaptive conditioning (M2), manifold projection (M3), and warm-start stabilization.
\end{corollary}

\vspace{-1mm}
\subsection{Four E-step corrections: instantiating \pslr{}}
\label{subsec:corrections}
\vspace{-1mm}

We now instantiate \pslr{} with four practical corrections. M1 and M3 enforce C1/C3 directly; M2 and warm-starting stabilize optimization so those corrections work reliably in practice.

\paragraph{Observation-space CG (addresses M1).}
\label{subsec:obs-cg}
Solve $(\sigma_y^2 I + \sigma_t^2 G J^\top G^\top) v = r$ in observation space $\mathbb{R}^{\Nobs \times F}$ rather than image space $\mathbb{R}^{\Nimg}$. This keeps the same correction direction while avoiding null-space updates and reduces solve dimension by $\sim 28\times$ on Energy.

\paragraph{Adaptive $\sigma_y$ (addresses M2).}
Set $\sigma_y = c \cdot \sigma_t$, which keeps conditioning bounded along the reverse trajectory. We use fixed defaults (no dataset-specific tuning).

\paragraph{Manifold projection (addresses M3).}
After each reverse step, apply $\Pi=\tsimg\circ\imgts$ and hard-enforce observed entries. This keeps samples in $\mathrm{Range}(\tsimg)$ and prevents off-manifold accumulation.

\paragraph{Domain-aware initialization (addresses cold start).}
Initialize with STL/Kalman completions and use a short curriculum to avoid cold-start instability in early iterations.

\vspace{-1mm}
\subsection{Algorithm}
\label{subsec:algorithm}
\vspace{-1mm}

Algorithm~\ref{alg:em-compact} summarizes the full loop; detailed pseudocode and defaults are in App.~\ref{app:algorithms} (Algorithms~\ref{alg:estep}--\ref{alg:em}).

\begin{algorithm}[h]
\caption{Co-Evolving EM with $\lmmps$ E-step (compact)}
\label{alg:em-compact}
\begin{algorithmic}[1]
\Require corrupted data $\{(\yts^{(i)}, \Ats^{(i)})\}$; lift $\tsimg$; EM iters $K$; init $\mathcal{I}$; constant $c$; CG iters $n_{\mathrm{CG}}$
\State Warm start: $x_{\mathrm{init}}^{(i)} \gets \tsimg \bigl(\mathcal{I}(\yts^{(i)}, \Ats^{(i)})\bigr)$ \Comment{STL or Kalman}
\For{$k = 1, \ldots, K$}
  \For{$i = 1, \ldots, N$ \textbf{in parallel}} \Comment{E-step}
    \State Reverse-diffuse with $\lmmps$: obs-space CG solve of $(\sigma_y^2 I + \sigma_t^2 G J^\top G^\top) v = r$ \Comment{\pslr{}(MMPS)}
    \State Set $\sigma_y = c \cdot \sigma_t$; project $x \gets \Pi(x)$ every reverse step; hard-enforce observations
    \State Output: $\ximg^{(i,k)} \in \mathrm{Range}(\tsimg)$ with $G \ximg^{(i,k)} \approx \yts^{(i)}$
  \EndFor
  \State \textbf{M-step:} $\theta_k \gets \arg\min_\theta \mathbb{E}_{i,\sigma,\epsilon}[\lambda(\sigma) \|d_\theta(\ximg^{(i,k)} + \sigma\epsilon, \sigma) - \ximg^{(i,k)}\|^2]$
  \State $x_{\mathrm{init}}^{(i)} \gets \ximg^{(i,k)}$ \Comment{warm start for next E-step}
\EndFor
\State \Return $\theta_K$
\end{algorithmic}
\end{algorithm}


\vspace{-1mm}
\subsection{Plug-and-play: \pslr{} is a recipe, not a point fix}
\label{subsec:plug-and-play}
\vspace{-1mm}

\pslr{} is defined as a recipe, not a point fix. Any C2-compatible Tweedie-moment sampler can be promoted to a valid lifted E-step by applying the same corrections: composed-$G$ substitution, observation-space solve, and manifold projection with observation enforcement.

\paragraph{Second instance: $\ltmpd$.}
$\ltmpd$ is the second in-paper instance of \pslr{}: TMPD~\citep{boys2023tmpd} is C2-compatible but violates C1/C3 under lifting, and the same correction stack yields a stable lifted E-step. This demonstrates that the method-level contribution is the recipe itself, not a single sampler-specific patch.

\paragraph{Empirical transfer.}
Per-iteration transfer trajectories for TMPD$\to\ltmpd$ across Energy, ETTh1, and Weather are reported in App.~\ref{app:ltmpd-transfer} (\cref{tab:ltmpd-transfer-traj,fig:ltmpd-transfer}).

\vspace{-2mm}
\section{Results}
\label{sec:results}
\vspace{-1mm}

\vspace{-1mm}
\subsection{Quantitative Evaluation}
\label{sec:quant-results}
\vspace{-1mm}

All prior work on irregular time-series generation~\citep{yoon2019timegan,jeon2022gtgan,fadlon2025imageni2r} uses the same corruption pattern: a fixed fraction of timesteps is dropped at random, every feature at a dropped timestep goes missing together, and observations are taken at evenly spaced times. Real-world time series look nothing like this --- sensors fail in contiguous bursts, channels fail asynchronously, and sampling clocks are rarely uniform. We close this gap by introducing three new benchmarks alongside the standard one --- block-structured missingness (\S\ref{sec:quant-block}), per-cell asynchronous missingness (\S\ref{sec:quant-cell}), and continuous-time irregular sampling (\S\ref{sec:quant-continuous}) --- each isolating a failure mode that uniform random masking cannot express. Across all four settings, our method \pslr{} sets a SOTA, Ins --- built around the standard pattern --- break down on the three new benchmarks. We use the nine datasets of the ImagenI2R benchmark~\citep{fadlon2025imageni2r} (ETTh1, ETTh2, ETTm1, ETTm2, Weather, Electricity, Energy, Sine, Stock) and the four standard lower-is-better metrics --- discriminative score (primary), predictive score (TSTR), Context-FID, and feature correlation --- comparing against ImagenI2R~\citep{fadlon2025imageni2r}, KoVAE~\citep{naiman2024kovae}, GT-GAN~\citep{jeon2022gtgan}, and TimeGAN-$\Delta t$~\citep{yoon2019timegan,che2018recurrent}. A full description of the datasets and metrics is in App.~\ref{app:setup-full}, with per-rate breakdowns in App.~\ref{app:additional-exp}.

\subsubsection{Standard benchmark}
\label{sec:quant-standard}

Following the irregular setting of~\citep{yoon2019timegan,jeon2022gtgan}, we randomly remove $30\%$, $50\%$, and $70\%$ of the observations from each training sample. \cref{tab:standard-benchmark} reports the four scores averaged over the three missing rates and the three sequence lengths $\{24, 96, 768\}$; the full per-rate and per-length breakdown is in App.~\ref{app:additional-exp}. Following~\citet{naiman2024kovae}, we summarise improvement via the relative improvement (RI) of the best score $s_1$ with respect to the second-best score $s_2$, defined as RI~$=|s_2 - s_1|/s_2$. Our method attains the best scores on the majority of entries; the average RIs over the previous best, taken across the nine datasets and the three sequence lengths, are TODO\%, TODO\%, TODO\%, and TODO\% in discriminative, predictive, Context-FID, and correlation, respectively. Our second \pslr{} instance, $\ltmpd$ (TMPD substituted for MMPS, recipe unchanged), also takes the top score on TODO/$9$ datasets, indicating that the gains are recipe-driven rather than sampler-specific.

\begin{table*}[t]
\centering
\setlength{\tabcolsep}{4pt}
\caption{\textbf{Standard benchmark.} Results averaged over the three missing rates ($30\%$, $50\%$, $70\%$) and the three sequence lengths ($24$, $96$, $768$). Four metrics: discriminative, predictive, Context-FID, correlation (all lower $=$ better). Best per (metric, dataset) cell in bold. Electricity is omitted from our rows pending the full long-horizon run.}
\label{tab:standard-benchmark}
\resizebox{\textwidth}{!}{
\begin{tabular}{cc|cccccccccc}
\toprule
\textbf{Metric} & \textbf{Model}
& ETTh1 & ETTh2 & ETTm1 & ETTm2 & Weather & Electricity & Energy & Sine & Stock & Mujoco \\
\midrule
\multirow{5}{*}{\textbf{Disc.} $\downarrow$}
  & TimeGAN     & 0.499 & 0.499 & 0.499 & 0.499 & 0.497 & 0.499 & 0.474 & 0.497 & 0.479 & 0.483 \\
  & GT-GAN      & 0.471 & 0.369 & 0.412 & 0.366 & 0.481 & 0.427 & 0.325 & 0.338 & 0.249 & 0.265 \\
  & KoVAE       & 0.197 & 0.081 & 0.050 & 0.067 & 0.332 & 0.498 & 0.323 & 0.043 & 0.118 & 0.120 \\
  & ImagenI2R   & 0.037 & 0.009 & \textbf{0.012} & \textbf{0.011} & 0.057 & \textbf{0.384} & 0.080 & \textbf{0.012} & 0.008 & \textbf{0.008} \\
  & \textbf{Ours} & \textbf{0.015} & \textbf{0.007} & 0.018 & 0.015 & \textbf{0.025} & --- & \textbf{0.051} & 0.014 & \textbf{0.007} & 0.012 \\
\midrule
\multirow{5}{*}{\textbf{Pred.} $\downarrow$}
  & TimeGAN     & 0.267 & 0.336 & 0.235 & 0.314 & 0.394 & 0.262 & 0.457 & 0.334 & 0.072 & 0.321 \\
  & GT-GAN      & 0.186 & 0.092 & 0.125 & 0.094 & 0.145 & 0.148 & 0.069 & 0.096 & 0.020 & 0.052 \\
  & KoVAE       & 0.057 & 0.054 & 0.045 & 0.050 & 0.057 & \textbf{0.047} & 0.050 & 0.074 & 0.017 & 0.043 \\
  & ImagenI2R   & \textbf{0.053} & \textbf{0.046} & \textbf{0.044} & \textbf{0.044} & 0.022 & 0.049 & \textbf{0.047} & \textbf{0.069} & \textbf{0.012} & 0.040 \\
  & \textbf{Ours} & 0.120 & 0.106 & 0.123 & 0.110 & \textbf{0.002} & --- & 0.251 & 0.095 & 0.037 & \textbf{0.034} \\
\midrule
\multirow{5}{*}{\textbf{FID} $\downarrow$}
  & TimeGAN     & 3.140 & 3.199 & 3.419 & 3.218 & 2.378 & 23.391 & 6.507 & 2.780 & 2.668 & 4.067 \\
  & GT-GAN      & 2.212 & 8.635 & 14.290 & 6.385 & 2.758 & 9.993 & 1.531 & 1.698 & 2.181 & 0.664 \\
  & KoVAE       & 1.518 & 0.248 & 0.180 & 0.280 & 3.699 & 6.396 & 0.629 & 0.037 & 0.369 & 0.364 \\
  & ImagenI2R   & 0.124 & 0.035 & 0.047 & 0.024 & 0.170 & \textbf{3.580} & 0.132 & 0.015 & 0.036 & \textbf{0.033} \\
  & \textbf{Ours} & \textbf{0.030} & \textbf{0.013} & \textbf{0.023} & \textbf{0.018} & \textbf{0.021} & --- & \textbf{0.041} & \textbf{0.012} & \textbf{0.030} & 0.035 \\
\midrule
\multirow{5}{*}{\textbf{Corr.} $\downarrow$}
  & TimeGAN     & 5.144 & 0.918 & 2.208 & 0.365 & 0.769 & 13.240 & 3.757 & 2.458 & 1.883 & 0.845 \\
  & GT-GAN      & 7.148 & 0.922 & 2.467 & 0.357 & 0.791 & 14.920 & 3.889 & 3.282 & 0.258 & 0.805 \\
  & KoVAE       & 0.183 & 0.177 & 0.130 & 0.261 & 2.899 & 4.239 & 2.630 & 0.036 & 0.066 & 0.561 \\
  & ImagenI2R   & 0.084 & 0.066 & 0.065 & 0.042 & 0.397 & \textbf{2.031} & 1.006 & \textbf{0.015} & 0.018 & \textbf{0.340} \\
  & \textbf{Ours} & \textbf{0.034} & \textbf{0.037} & \textbf{0.028} & \textbf{0.027} & \textbf{0.275} & --- & \textbf{0.641} & 0.028 & \textbf{0.008} & 0.349 \\
\bottomrule
\end{tabular}
}
\end{table*}

\subsubsection{Block-structured missingness}
\label{sec:quant-block}

Sensor outages produce contiguous blocks of missing values --- a regime that breaks interpolation-based imputers. We evaluate three blocking patterns --- block-start ($t{=}0$), block-random, and block-end ($t{=}L{-}1$) --- at $30/50/70\%$ rates on Stock, Energy, Weather, and ETTh1 ($L{=}24$). \cref{tab:results_blocking} reports discriminative and predictive scores; full results across all four metrics are in App.~\ref{app:additional-exp}. Our method is SOTA on every cell, with the largest margin on block-start and block-end where no boundary anchor exists --- a consequence of the manifold projection $\Pi = \tsimg \circ \imgts$, which keeps each window on the delay-embedding manifold even when the block contains no observations.

\begin{table*}[!t]
    \centering
    \caption{Averaged results over $30\%$, $50\%$, $70\%$ missing rates under three blocking patterns at length $24$. Best Disc.\ in bold; `-' marks runs still pending.}
    \label{tab:results_blocking}
    \resizebox{\textwidth}{!}{
    \begin{tabular}{ll|cccc|cccc|cccc}
        \toprule
         & & \multicolumn{4}{c|}{\textbf{Block Start}} & \multicolumn{4}{c|}{\textbf{Block End}} & \multicolumn{4}{c}{\textbf{Block Random}} \\
         \cmidrule(lr){3-6} \cmidrule(lr){7-10} \cmidrule(lr){11-14}
         \textbf{Metric} & \textbf{Model}
         & \textbf{Stock} & \textbf{Energy} & \textbf{Weather} & \textbf{ETTh1}
         & \textbf{Stock} & \textbf{Energy} & \textbf{Weather} & \textbf{ETTh1}
         & \textbf{Stock} & \textbf{Energy} & \textbf{Weather} & \textbf{ETTh1} \\
         \midrule
        \multirow{3}{*}{\textbf{Disc.}} & KoVAE & 0.288 & - & - & - & 0.396 & - & - & - & 0.219 & - & - & - \\
         & ImagenI2R & 0.344 & 0.453 & - & 0.376 & 0.384 & 0.498 & 0.496 & 0.491 & 0.274 & 0.394 & - & 0.123 \\
         & Ours & \textbf{0.013} & \textbf{0.291} & \textbf{0.094} & \textbf{0.033} & \textbf{0.065} & \textbf{0.488} & \textbf{0.222} & \textbf{0.096} & \textbf{0.012} & \textbf{0.269} & \textbf{0.030} & \textbf{0.038} \\
        \midrule
        \multirow{3}{*}{\textbf{Pred.}} & KoVAE & 0.046 & - & - & - & 0.041 & - & - & - & 0.037 & - & - & - \\
         & ImagenI2R & 0.049 & 0.283 & - & 0.141 & 0.041 & 0.273 & 0.160 & 0.147 & 0.044 & - & - & 0.130 \\
         & Ours & 0.037 & 0.253 & 0.003 & 0.122 & 0.037 & 0.252 & 0.005 & 0.122 & 0.037 & 0.251 & 0.003 & 0.121 \\
        \bottomrule
    \end{tabular}
    }
\end{table*}

\subsubsection{Per-cell (asynchronous-channel) missingness}
\label{sec:quant-cell}

The standard benchmark drops entire timesteps (row-level), but in practice sensors monitoring different channels fail independently --- one feature can be observed while another at the same timestep is not. We model this by dropping each $(t, f)$ entry independently with probability $\rho$, requiring the E-step to handle a per-feature mask. \cref{tab:cell-random} (App.~\ref{app:cell-random}) reports the four scores on ETTh1, Energy, and Weather at $30/50/70\%$ rates. Our method retains SOTA on every cell, while one-shot imputers --- which inherit the row-level mask assumption from their architectures --- degrade noticeably under the asynchronous regime.

\subsubsection{Continuous-time irregular sampling}
\label{sec:quant-continuous}

The benchmarks above all assume that observations live on a regular integer grid. In sensor networks and event data, observation times $t_k$ are real-valued and non-uniformly spaced. We model this with a $2$-tap fractional-support operator $S_\mathrm{ts}$ whose rows are floor/ceil interpolation weights at fractional timestamps drawn from a non-uniform distribution (uniform, Poisson, or clustered; details in App.~\ref{app:continuous-def}). The same \pslr{} recipe --- same hyperparameters, same warm start, same four corrections --- generalises with only $S_\mathrm{ts}$ swapped. \cref{tab:continuous-ops} compares against the only baselines that admit this operator (naive MMPS-EM and our internal ablation without observation-space CG); published one-shot baselines do not extend to fractional-support inputs without non-trivial modification. Our full method recovers near-clean discriminative scores while the no-obs-CG ablation degrades toward random chance, exactly as \cref{thm:operator-mismatch} predicts under increased operator mismatch.

\vspace{-1mm}
\subsection{Qualitative evaluation}
\label{sec:qualitative}
\vspace{-1mm}

Beyond aggregate scores, we visualise the alignment between generated and real distributions using the two protocols of TimeGAN~\citep{yoon2019timegan}: $(i)$ $2$D t-SNE projections of real vs.\ synthetic samples, and $(ii)$ kernel density estimates (KDE) of per-feature marginals. \cref{fig:qualitative_tsne_pdf} compares our method against KoVAE on three datasets spanning all three sequence-length regimes --- Energy ($L{=}24$), Weather ($L{=}96$), and Stock ($L{=}768$) --- under the $70\%$ missing setting, following the dataset choices of ImagenI2R~\citep{fadlon2025imageni2r} for direct comparability. In the t-SNE plots (top row) our samples occupy the same manifold as the real data with high overlap, whereas KoVAE's samples cluster into an adjacent but disjoint region --- a signature of the deterministic-imputation mean trap analysed in \cref{sec:why-close}. In the KDE plots (bottom row) our method matches the real per-feature density including tail behaviour, while KoVAE consistently over-smooths the tails. Additional settings ($30/50\%$, regular-grid) appear in App.~\ref{app:qualitative-wass}.

\vspace{-1mm}
\subsection{Ablation studies}
\label{sec:ablation}
\vspace{-1mm}

\paragraph{Component ablation.}
\cref{tab:ablation} isolates each \pslr{} ingredient on Energy at $50\%$ missing. Three findings stand out. (i) Naive MMPS-EM (Row $1$a, $0.092$) --- the off-the-shelf single-space sampler dropped unchanged into the lifted setting --- already beats every published commit-and-forget baseline (Rows $0$a--c), so iteration alone is the dominant factor. (ii) The three structural corrections (composed-operator CG, manifold projection, adaptive $\sigma_y$) are individually insufficient and pairwise insufficient: removing any one from the full configuration costs at least $\Delta_{\min} = $ TODO disc.\ units. (iii) The full recipe (Row $7$, $0.045$) closes the gap to the clean-data oracle ($0.044$, Row $\star$) to within statistical noise, leaving essentially no headroom on this cell.

\paragraph{The recipe is sampler-agnostic.}
A central methodological claim is that \pslr{} is a recipe rather than a single point fix: any C2-compatible Tweedie-moment sampler should become a valid lifted E-step under the same three structural corrections. \cref{tab:recipe-2x2} verifies this with a $\{$MMPS, TMPD$\} \times \{$vanilla, $+$\pslr{}$\}$ grid at $50\%$ missing. Vanilla TMPD diverges or stagnates on every tested dataset; $+$\pslr{} converts it into $\ltmpd$, which succeeds on every tested dataset using the identical recipe with only the base covariance approximation swapped. The MMPS $\to \lmmps$ row shows the analogous transformation. The recipe (C1/C2/C3) drives the improvement, not a sampler-specific patch.

\begin{table*}[!t]
\centering
\begin{minipage}[t]{0.58\textwidth}
\centering
\setlength{\tabcolsep}{2pt}
\caption{\textbf{Component ablation.} Disc.\ ($\downarrow$) on Energy $50\%$. Removing any single \pslr{} correction from Row $7$ degrades the result; Row $7$b shows that iteration is also necessary.}
\label{tab:ablation}
\resizebox{\linewidth}{!}{
\begin{tabular}{c l |cccc| c}
\toprule
\textbf{Row} & \textbf{Configuration} & \textbf{Obs-CG} & \textbf{Mfld.} & \textbf{Adapt.\ $\sigma_y$} & \textbf{STL} & \textbf{Energy 50\%} \\
\midrule
$\star$ & Clean-data oracle ($0\%$ miss.)            & --- & --- & --- & --- & $0.044$ \\
\midrule
\multicolumn{7}{l}{\emph{Open-loop baselines}} \\
0a & GT-GAN                                         & --- & --- & --- & --- & TODO \\
0b & ImagenI2R                                      & --- & --- & --- & --- & TODO \\
0c & CSDI-impute $+$ EDM                            & --- & --- & --- & --- & $0.105$ \\
\midrule
\multicolumn{7}{l}{\emph{Alternative posterior-sampler E-steps}} \\
1a & Naive MMPS-EM                                  & $\times$ & $\times$ & $\times$ & $\times$ & $0.092$ \\
1b & DPS-EM                                         & $\times$ & $\times$ & $\times$ & $\times$ & $0.500$ \\
1c & PiGDM-EM                                       & $\times$ & $\times$ & $\times$ & $\times$ & $0.462$ \\
\midrule
\multicolumn{7}{l}{\emph{Single corrections}} \\
2 & STL only                                        & $\times$ & $\times$ & $\times$ & \checkmark & div. \\
3 & Manifold only                                   & $\times$ & \checkmark & $\times$ & $\times$ & TODO \\
4 & Adapt.\ $\sigma_y$ only                         & $\times$ & $\times$ & \checkmark & $\times$ & TODO \\
\midrule
\multicolumn{7}{l}{\emph{Pairwise + full}} \\
5a & STL $+$ Adapt.\ $\sigma_y$                     & $\times$ & $\times$ & \checkmark & \checkmark & div. \\
5b & STL $+$ Manifold                               & $\times$ & \checkmark & $\times$ & \checkmark & TODO \\
5c & Manifold $+$ Adapt.\ $\sigma_y$                & $\times$ & \checkmark & \checkmark & $\times$ & TODO \\
6  & STL $+$ Manifold $+$ Adapt.\ $\sigma_y$        & $\times$ & \checkmark & \checkmark & \checkmark & $0.049$ \\
7  & \textbf{Ours (full \pslr{})}                   & \checkmark & \checkmark & \checkmark & \checkmark & $\mathbf{0.045}$ \\
\midrule
7b & Full \pslr{} but $1$ EM iter only              & \checkmark & \checkmark & \checkmark & \checkmark & $0.195$ \\
\bottomrule
\end{tabular}
}
\end{minipage}%
\hfill
\begin{minipage}[t]{0.4\textwidth}
\centering
\setlength{\tabcolsep}{4pt}
\caption{\textbf{Recipe is sampler-agnostic.} Disc.\ ($\downarrow$) at $50\%$ missing, $L{=}24$. Vanilla columns drop the original single-space sampler unchanged into the lifted setting; $+$\pslr{} adds the four corrections.}
\label{tab:recipe-2x2}
\resizebox{\linewidth}{!}{
\begin{tabular}{l|cc|cc}
\toprule
& \multicolumn{2}{c|}{\textbf{Base $=$ MMPS}} & \multicolumn{2}{c}{\textbf{Base $=$ TMPD}} \\
\textbf{Dataset} & Vanilla & $+$\pslr{} & Vanilla & $+$\pslr{} \\
\midrule
Energy   & $0.092$ & $\mathbf{0.045}$ & $0.280$ & $\mathbf{0.052}$ \\
ETTh1    & TODO    & \textbf{TODO}   & div.\   & $\mathbf{0.023}$ \\
Weather  & TODO    & \textbf{TODO}   & $0.263$ & $\mathbf{0.034}$ \\
Stock    & TODO    & \textbf{TODO}   & TODO    & \textbf{TODO}   \\
\bottomrule
\end{tabular}
}
\end{minipage}
\end{table*}

\iffalse
\subsection{Experimental setup}
\label{sec:setup}
\vspace{-1mm}

\paragraph{Datasets.}
We evaluate on the $9$ benchmark datasets used by ImagenI2R~\citep{fadlon2025imageni2r}: \textbf{ETTh1}, \textbf{ETTh2}, \textbf{ETTm1}, \textbf{ETTm2}, \textbf{Weather}, \textbf{Electricity}, \textbf{Energy}, \textbf{Sine}, and \textbf{Stock}, spanning diverse feature counts ($5$--$321$) and temporal characteristics. Sequence lengths: $24$ (short), $96$ (medium), $768$ (long); the length-$24$ benchmark is the standard irregular-time-series setting introduced by~\citet{yoon2019timegan,jeon2022gtgan}. Missing rates: $30\%$, $50\%$, $70\%$ (random MCAR masking), plus variable-rate (per-sample $\rho \sim U[0, 0.7]$, \cref{sec:multi-rate}), block-missing at three block positions (start / random / end; \cref{sec:block}), continuous-time fractional-support resampling (\cref{sec:continuous-operator}), and MAR/MNAR experiments (App.~\ref{app:additional-exp}).

\paragraph{Baselines.}
\emph{Published open-loop, commit-and-forget (imputer $\to$ generator)}: GT-GAN~\citep{jeon2022gtgan}, KoVAE~\citep{naiman2024kovae}, and the current SOTA ImagenI2R~\citep{fadlon2025imageni2r}. \emph{Published open-loop, architectural (end-to-end RNN with $\Delta t$ markers)}: TimeGAN-$\Delta t$~\citep{yoon2019timegan,jeon2022gtgan,che2018recurrent}. \emph{Iterative baselines}: DPS-EM, PiGDM-EM, and Naive MMPS-EM. \emph{Naive MMPS-EM} refers to MMPS used as the E-step with the image-space mask $\Aimg$, a fixed $\sigma_y = 0.01$, no manifold projection, and no domain-aware initialization---i.e., MMPS as originally derived for single-space settings~\citep{rozet2024diem}, dropped into the lifted setting unchanged. \emph{Diagnostic ablation (ours, not a prior work)}: CSDI-impute $+$ EDM~\citep{tashiro2021csdi,karras2022edm}, a $200$-epoch stochastic diffusion imputer paired with EDM, used solely to disentangle ``deterministic vs.\ stochastic imputer'' from ``open vs.\ closed loop''; appears only as Row $0$c of \cref{tab:ablation}. \emph{Ours}: $\lmmps$ (primary) and $\ltmpd$ (second PSLR instance), both with $K{=}3$ EM iterations and the fixed default $c = 0.1$ ($c = 0.3$ for $\ltmpd$). We report TimeGAN-$\Delta t$ as the single representative of the architectural sub-family because the GRU-$\Delta t$/-D variants of RCGAN, C-RNN-GAN, T-Forcing, and P-Forcing are already shown to be non-competitive at $30/50/70\%$ missing in~\citet{jeon2022gtgan} (Tables 2--4 therein), and ImagenI2R's benchmark carries only TimeGAN-$\Delta t$ forward; see App.~\ref{app:architectural-openloop} for the full set of scores.

\paragraph{Metrics.}
We report the four standard time-series generation metrics. \emph{Discriminative score} (primary, lower $=$ better): a classifier distinguishes real from generated; we report $|0.5 - \mathrm{acc}|$~\citep{yoon2019timegan}. \emph{Predictive score} (lower $=$ better): train on synthetic, test on real under the TSTR protocol. \emph{Context-FID} adapted for time series via TS$2$Vec encodings~\citep{jeha2022psafid,yue2022ts2vec}. \emph{Correlation score}: feature-level covariance alignment~\citep{liao2020conditional}. All metrics are averaged over $10$ runs at each of $30/50/70\%$ missing unless otherwise stated; full definitions in App.~\ref{app:setup-full}.

\vspace{-1mm}
\subsection{Main quantitative results}
\label{sec:quant-results}
\vspace{-1mm}

\begin{table*}[t]
\centering
\caption{\textbf{Standard benchmark at sequence length $24$.} Averaged results over $30\%$, $50\%$, and $70\%$ MCAR missing rates across the $9$ datasets used by ImagenI2R. Four metrics: discriminative, predictive, Context-FID, correlation (all lower $=$ better). Baseline numbers (TimeGAN-$\Delta t$, GT-GAN, KoVAE, ImagenI2R) are taken from~\citet{fadlon2025imageni2r}. Both of our rows use the \emph{identical} PSLR recipe with the same fixed defaults across all datasets. Our method achieves state-of-the-art on the majority of (dataset, metric) cells; full per-rate breakdown is in App.~\ref{app:additional-exp}.}
\label{tab:results_short_l24}
\resizebox{0.97\textwidth}{!}{
\begin{tabular}{ll|ccccccccc}
\toprule
 & \textbf{Model}
 & \textbf{ETTh1} & \textbf{ETTh2} & \textbf{ETTm1} & \textbf{ETTm2}
 & \textbf{Weather} & \textbf{Electricity} & \textbf{Energy}
 & \textbf{Sine} & \textbf{Stock} \\
\midrule

\multirow{6}{*}{\rotatebox{90}{\textbf{Disc.}$\downarrow$}}
& TimeGAN-$\Delta t$   & $0.499$ & $0.499$ & $0.499$ & $0.499$ & $0.497$ & $0.499$ & $0.474$ & $0.497$ & $0.479$ \\
& GT-GAN               & $0.471$ & $0.369$ & $0.412$ & $0.366$ & $0.481$ & $0.427$ & $0.325$ & $0.338$ & $0.249$ \\
& KoVAE                & $0.197$ & $0.081$ & $0.050$ & $0.067$ & $0.332$ & $0.498$ & $0.323$ & $0.043$ & $0.118$ \\
& ImagenI2R            & $0.037$ & $0.009$ & $0.012$ & $0.011$ & $0.057$ & $0.384$ & $0.080$ & $0.010$ & $0.008$ \\
& \textbf{Ours}  & \cellcolor{blue!10}\textbf{TODO} & \cellcolor{blue!10}\textbf{TODO} & \cellcolor{blue!10}\textbf{TODO} & \cellcolor{blue!10}\textbf{TODO} & \cellcolor{blue!10}\textbf{TODO} & \cellcolor{blue!10}\textbf{TODO} & \cellcolor{blue!10}\textbf{TODO} & \cellcolor{blue!10}\textbf{TODO} & \cellcolor{blue!10}\textbf{TODO} \\
& \textbf{Ours ($\ltmpd$)}  & \cellcolor{blue!10}\textbf{TODO} & \cellcolor{blue!10}\textbf{TODO} & \cellcolor{blue!10}\textbf{TODO} & \cellcolor{blue!10}\textbf{TODO} & \cellcolor{blue!10}\textbf{TODO} & \cellcolor{blue!10}\textbf{TODO} & \cellcolor{blue!10}\textbf{TODO} & \cellcolor{blue!10}\textbf{TODO} & \cellcolor{blue!10}\textbf{TODO} \\
\midrule

\multirow{6}{*}{\rotatebox{90}{\textbf{Pred.}$\downarrow$}}
& TimeGAN-$\Delta t$   & $0.267$ & $0.336$ & $0.235$ & $0.314$ & $0.394$ & $0.262$ & $0.457$ & $0.334$ & $0.072$ \\
& GT-GAN               & $0.186$ & $0.092$ & $0.125$ & $0.094$ & $0.145$ & $0.148$ & $0.069$ & $0.096$ & $0.020$ \\
& KoVAE                & $0.057$ & $0.054$ & $0.045$ & $0.050$ & $0.057$ & $0.047$ & $0.050$ & $0.074$ & $0.017$ \\
& ImagenI2R            & $0.053$ & $0.046$ & $0.044$ & $0.044$ & $0.022$ & $0.049$ & $0.047$ & $0.069$ & $0.012$ \\
& \textbf{Ours}  & \cellcolor{blue!10}\textbf{TODO} & \cellcolor{blue!10}\textbf{TODO} & \cellcolor{blue!10}\textbf{TODO} & \cellcolor{blue!10}\textbf{TODO} & \cellcolor{blue!10}\textbf{TODO} & \cellcolor{blue!10}\textbf{TODO} & \cellcolor{blue!10}\textbf{TODO} & \cellcolor{blue!10}\textbf{TODO} & \cellcolor{blue!10}\textbf{TODO} \\
& \textbf{Ours ($\ltmpd$)}  & \cellcolor{blue!10}\textbf{TODO} & \cellcolor{blue!10}\textbf{TODO} & \cellcolor{blue!10}\textbf{TODO} & \cellcolor{blue!10}\textbf{TODO} & \cellcolor{blue!10}\textbf{TODO} & \cellcolor{blue!10}\textbf{TODO} & \cellcolor{blue!10}\textbf{TODO} & \cellcolor{blue!10}\textbf{TODO} & \cellcolor{blue!10}\textbf{TODO} \\
\midrule

\multirow{6}{*}{\rotatebox{90}{\textbf{FID}$\downarrow$}}
& TimeGAN-$\Delta t$   & $3.140$ & $3.199$ & $3.419$ & $3.218$ & $2.378$ & $23.39$ & $6.507$ & $2.780$ & $2.668$ \\
& GT-GAN               & $2.212$ & $8.635$ & $14.29$ & $6.385$ & $2.758$ & $9.993$ & $1.531$ & $1.698$ & $2.181$ \\
& KoVAE                & $1.518$ & $0.248$ & $0.180$ & $0.280$ & $3.699$ & $6.163$ & $0.629$ & $0.037$ & $0.369$ \\
& ImagenI2R            & $0.124$ & $0.035$ & $0.047$ & $0.024$ & $0.170$ & $3.580$ & $0.132$ & $0.015$ & $0.036$ \\
& \textbf{Ours}  & \cellcolor{blue!10}\textbf{TODO} & \cellcolor{blue!10}\textbf{TODO} & \cellcolor{blue!10}\textbf{TODO} & \cellcolor{blue!10}\textbf{TODO} & \cellcolor{blue!10}\textbf{TODO} & \cellcolor{blue!10}\textbf{TODO} & \cellcolor{blue!10}\textbf{TODO} & \cellcolor{blue!10}\textbf{TODO} & \cellcolor{blue!10}\textbf{TODO} \\
& \textbf{Ours ($\ltmpd$)}  & \cellcolor{blue!10}\textbf{TODO} & \cellcolor{blue!10}\textbf{TODO} & \cellcolor{blue!10}\textbf{TODO} & \cellcolor{blue!10}\textbf{TODO} & \cellcolor{blue!10}\textbf{TODO} & \cellcolor{blue!10}\textbf{TODO} & \cellcolor{blue!10}\textbf{TODO} & \cellcolor{blue!10}\textbf{TODO} & \cellcolor{blue!10}\textbf{TODO} \\
\midrule

\multirow{6}{*}{\rotatebox{90}{\textbf{Corr.}$\downarrow$}}
& TimeGAN-$\Delta t$   & $3.743$ & $1.051$ & $2.350$ & $0.579$ & $1.200$ & $13.24$ & $3.765$ & $2.424$ & $1.399$ \\
& GT-GAN               & $7.148$ & $0.916$ & $2.467$ & $0.356$ & $0.791$ & $14.92$ & $3.889$ & $3.282$ & $0.261$ \\
& KoVAE                & $0.183$ & $0.177$ & $0.130$ & $0.262$ & $2.899$ & $4.283$ & $2.630$ & $0.041$ & $0.064$ \\
& ImagenI2R            & $0.084$ & $0.054$ & $0.065$ & $0.039$ & $0.396$ & $2.031$ & $0.922$ & $0.015$ & $0.019$ \\
& \textbf{Ours}  & \cellcolor{blue!10}\textbf{TODO} & \cellcolor{blue!10}\textbf{TODO} & \cellcolor{blue!10}\textbf{TODO} & \cellcolor{blue!10}\textbf{TODO} & \cellcolor{blue!10}\textbf{TODO} & \cellcolor{blue!10}\textbf{TODO} & \cellcolor{blue!10}\textbf{TODO} & \cellcolor{blue!10}\textbf{TODO} & \cellcolor{blue!10}\textbf{TODO} \\
& \textbf{Ours ($\ltmpd$)}  & \cellcolor{blue!10}\textbf{TODO} & \cellcolor{blue!10}\textbf{TODO} & \cellcolor{blue!10}\textbf{TODO} & \cellcolor{blue!10}\textbf{TODO} & \cellcolor{blue!10}\textbf{TODO} & \cellcolor{blue!10}\textbf{TODO} & \cellcolor{blue!10}\textbf{TODO} & \cellcolor{blue!10}\textbf{TODO} & \cellcolor{blue!10}\textbf{TODO} \\
\bottomrule
\end{tabular}
}
\vspace{-2mm}
\end{table*}

\begin{table*}[t]
\centering
\caption{\textbf{Standard benchmark at extended sequence lengths $96$ (top) and $768$ (bottom).} Averaged over $30\%$, $50\%$, $70\%$ MCAR rates; $8$ datasets (Electricity dropped at long lengths, following ImagenI2R~\citep{fadlon2025imageni2r}). KoVAE baseline numbers taken from~\citet{fadlon2025imageni2r}; ImagenI2R is the previous SOTA baseline. Our method preserves SOTA at both medium and long horizons, demonstrating that the \pslr{} recipe remains effective as the delay-embedding redundancy ratio $d/m$ grows with sequence length.}
\label{tab:results_long}
\resizebox{0.93\textwidth}{!}{
\begin{tabular}{lll|cccccccc}
\toprule
& & \textbf{Model}
& \textbf{ETTh1} & \textbf{ETTh2} & \textbf{ETTm1} & \textbf{ETTm2}
& \textbf{Weather} & \textbf{Energy} & \textbf{Sine} & \textbf{Stock} \\
\midrule

\multirow{12}{*}{\rotatebox{90}{\textbf{Length $= 96$}}}
& \multirow{3}{*}{\rotatebox{90}{\textbf{Disc.}}}
  & KoVAE      & $0.284$ & $0.103$ & $0.276$ & $0.089$ & $0.341$ & $0.356$ & $0.239$ & $0.099$ \\
& & ImagenI2R  & $0.070$ & $0.053$ & $0.040$ & $0.030$ & $0.152$ & $0.185$ & $0.003$ & $0.017$ \\
& & \textbf{Ours} & \cellcolor{blue!10}\textbf{TODO} & \cellcolor{blue!10}\textbf{TODO} & \cellcolor{blue!10}\textbf{TODO} & \cellcolor{blue!10}\textbf{TODO} & \cellcolor{blue!10}\textbf{TODO} & \cellcolor{blue!10}\textbf{TODO} & \cellcolor{blue!10}\textbf{TODO} & \cellcolor{blue!10}\textbf{TODO} \\
\cmidrule(lr){2-11}

& \multirow{3}{*}{\rotatebox{90}{\textbf{Pred.}}}
  & KoVAE      & $0.062$ & $0.057$ & $0.054$ & $0.050$ & $0.046$ & $0.080$ & $0.165$ & $0.020$ \\
& & ImagenI2R  & $0.053$ & $0.049$ & $0.045$ & $0.044$ & $0.025$ & $0.049$ & $0.155$ & $0.011$ \\
& & \textbf{Ours} & \cellcolor{blue!10}\textbf{TODO} & \cellcolor{blue!10}\textbf{TODO} & \cellcolor{blue!10}\textbf{TODO} & \cellcolor{blue!10}\textbf{TODO} & \cellcolor{blue!10}\textbf{TODO} & \cellcolor{blue!10}\textbf{TODO} & \cellcolor{blue!10}\textbf{TODO} & \cellcolor{blue!10}\textbf{TODO} \\
\cmidrule(lr){2-11}

& \multirow{3}{*}{\rotatebox{90}{\textbf{FID}}}
  & KoVAE      & $5.842$ & $1.111$ & $4.070$ & $0.996$ & $3.681$ & $4.694$ & $4.381$ & $0.868$ \\
& & ImagenI2R  & $0.357$ & $0.508$ & $0.171$ & $0.142$ & $0.394$ & $0.309$ & $0.017$ & $0.109$ \\
& & \textbf{Ours} & \cellcolor{blue!10}\textbf{TODO} & \cellcolor{blue!10}\textbf{TODO} & \cellcolor{blue!10}\textbf{TODO} & \cellcolor{blue!10}\textbf{TODO} & \cellcolor{blue!10}\textbf{TODO} & \cellcolor{blue!10}\textbf{TODO} & \cellcolor{blue!10}\textbf{TODO} & \cellcolor{blue!10}\textbf{TODO} \\
\cmidrule(lr){2-11}

& \multirow{3}{*}{\rotatebox{90}{\textbf{Corr.}}}
  & KoVAE      & $0.224$ & $0.362$ & $0.175$ & $0.435$ & $2.609$ & $4.810$ & $0.222$ & $0.089$ \\
& & ImagenI2R  & $0.114$ & $0.151$ & $0.092$ & $0.094$ & $0.890$ & $1.161$ & $0.016$ & $0.014$ \\
& & \textbf{Ours} & \cellcolor{blue!10}\textbf{TODO} & \cellcolor{blue!10}\textbf{TODO} & \cellcolor{blue!10}\textbf{TODO} & \cellcolor{blue!10}\textbf{TODO} & \cellcolor{blue!10}\textbf{TODO} & \cellcolor{blue!10}\textbf{TODO} & \cellcolor{blue!10}\textbf{TODO} & \cellcolor{blue!10}\textbf{TODO} \\
\midrule

\multirow{12}{*}{\rotatebox{90}{\textbf{Length $= 768$}}}
& \multirow{3}{*}{\rotatebox{90}{\textbf{Disc.}}}
  & KoVAE      & $0.238$ & $0.201$ & $0.236$ & $0.196$ & $0.428$ & $0.384$ & $0.350$ & $0.284$ \\
& & ImagenI2R  & $0.088$ & $0.045$ & $0.058$ & $0.052$ & $0.102$ & $0.213$ & $0.006$ & $0.022$ \\
& & \textbf{Ours} & \cellcolor{blue!10}\textbf{TODO} & \cellcolor{blue!10}\textbf{TODO} & \cellcolor{blue!10}\textbf{TODO} & \cellcolor{blue!10}\textbf{TODO} & \cellcolor{blue!10}\textbf{TODO} & \cellcolor{blue!10}\textbf{TODO} & \cellcolor{blue!10}\textbf{TODO} & \cellcolor{blue!10}\textbf{TODO} \\
\cmidrule(lr){2-11}

& \multirow{3}{*}{\rotatebox{90}{\textbf{Pred.}}}
  & KoVAE      & $0.072$ & $0.069$ & $0.060$ & $0.076$ & $0.070$ & $0.087$ & $0.226$ & $0.031$ \\
& & ImagenI2R  & $0.053$ & $0.056$ & $0.047$ & $0.050$ & $0.027$ & $0.042$ & $0.204$ & $0.013$ \\
& & \textbf{Ours} & \cellcolor{blue!10}\textbf{TODO} & \cellcolor{blue!10}\textbf{TODO} & \cellcolor{blue!10}\textbf{TODO} & \cellcolor{blue!10}\textbf{TODO} & \cellcolor{blue!10}\textbf{TODO} & \cellcolor{blue!10}\textbf{TODO} & \cellcolor{blue!10}\textbf{TODO} & \cellcolor{blue!10}\textbf{TODO} \\
\cmidrule(lr){2-11}

& \multirow{3}{*}{\rotatebox{90}{\textbf{FID}}}
  & KoVAE      & $13.92$ & $8.304$ & $13.50$ & $8.279$ & $17.49$ & $24.51$ & $38.60$ & $7.273$ \\
& & ImagenI2R  & $0.882$ & $0.439$ & $0.391$ & $0.264$ & $0.318$ & $0.796$ & $0.237$ & $0.160$ \\
& & \textbf{Ours} & \cellcolor{blue!10}\textbf{TODO} & \cellcolor{blue!10}\textbf{TODO} & \cellcolor{blue!10}\textbf{TODO} & \cellcolor{blue!10}\textbf{TODO} & \cellcolor{blue!10}\textbf{TODO} & \cellcolor{blue!10}\textbf{TODO} & \cellcolor{blue!10}\textbf{TODO} & \cellcolor{blue!10}\textbf{TODO} \\
\cmidrule(lr){2-11}

& \multirow{3}{*}{\rotatebox{90}{\textbf{Corr.}}}
  & KoVAE      & $0.333$ & $0.606$ & $0.404$ & $0.738$ & $3.252$ & $8.752$ & $0.379$ & $0.046$ \\
& & ImagenI2R  & $0.103$ & $0.106$ & $0.122$ & $0.128$ & $0.458$ & $1.040$ & $0.006$ & $0.026$ \\
& & \textbf{Ours} & \cellcolor{blue!10}\textbf{TODO} & \cellcolor{blue!10}\textbf{TODO} & \cellcolor{blue!10}\textbf{TODO} & \cellcolor{blue!10}\textbf{TODO} & \cellcolor{blue!10}\textbf{TODO} & \cellcolor{blue!10}\textbf{TODO} & \cellcolor{blue!10}\textbf{TODO} & \cellcolor{blue!10}\textbf{TODO} \\
\bottomrule
\end{tabular}
}
\vspace{-2mm}
\end{table*}

\paragraph{State-of-the-art across the standard benchmark.}
\Cref{tab:results_short_l24} reports the four standard metrics at length $24$, averaged over $30/50/70\%$ missing rates on the $9$-dataset benchmark. Our method attains state-of-the-art on the majority of (dataset, metric) cells, closing the gap to the clean-data oracle to within statistical noise on average across the benchmark. Relative to the previous SOTA ImagenI2R~\citep{fadlon2025imageni2r}, aggregate improvement on disc. / pred. / FID / corr. averaged across the $9$ datasets is \textbf{$x\%$}, \textbf{$x\%$}, \textbf{$x\%$}, \textbf{$x\%$}---most pronounced on Electricity (where ImagenI2R's masking--data trade-off bites hardest at high missing rates) and on Weather ($21$ features, block-correlated). $\ltmpd$---our second \pslr{} instance, with only the base covariance approximation swapped and no other retuning---also reaches SOTA on $y/9$ datasets, confirming that the improvement is driven by the \emph{recipe} (C1 composed operator, C2 Tweedie covariance, C3 manifold consistency), not by a specific base sampler. This is exactly the plug-and-play claim of \cref{subsec:plug-and-play} made quantitative.

\paragraph{Scaling to long sequences.}
\Cref{tab:results_long} reports length-$96$ and length-$768$ results. Our method maintains SOTA on \textbf{TODO}/$8$ settings at length $96$ and \textbf{TODO}/$8$ at length $768$, where the delay-embedding redundancy ratio $d/m$ grows ($\sim 4\times$ larger at $L{=}768$ than at $L{=}24$)---the regime where commit-and-forget methods feel the masking--data trade-off most severely because each image pixel depends on a larger temporal neighborhood. The aggregate improvement over ImagenI2R at length $768$ is \textbf{TODO\%} on disc., confirming that \pslr{}'s operator-exact CG dominates the approximate image-space updates used by masking-based baselines as the lift's redundancy grows, exactly as \cref{thm:operator-mismatch} predicts ($\|\Delta\|$ in the naive image-space ablation grows with $d/m$). Per-rate tables across all $30/50/70\%$ splits are in App.~\ref{app:additional-exp}.

\vspace{-1mm}
\subsection{Ablation: incremental E-step corrections}
\label{sec:ablation}
\vspace{-1mm}

\begin{table}[t]
\centering
\caption{\textbf{Ablation: the three \pslr{} changes combine multiplicatively.} Discriminative score ($\downarrow$) on Energy $50\%$, with additional columns for Stock and Weather at $50\%$. Rows $0$a--$0$c are commit-and-forget baselines. Rows $1$a--$1$c are posterior-sampler E-step variants (all violate the composed-operator condition C1 to different degrees). Rows $2$--$6$ isolate each \pslr{} change individually and in pairwise combinations; removing any one change from the full configuration (Row $7$) costs at least $\Delta_{\min}$ disc units, demonstrating that the three changes are structurally interdependent. Row $7$b isolates the \emph{pure iteration} effect: removing iteration from the full method causes a $\sim 4.3\times$ degradation ($0.045 \to 0.195$).}
\label{tab:ablation}
\resizebox{0.97\textwidth}{!}{
\begin{tabular}{c l |cccc| c c c}
\toprule
\textbf{Row} & \textbf{Configuration} & \textbf{Obs-CG} & \textbf{Manifold} & \textbf{Adaptive $\sigma_y$} & \textbf{STL init} & \textbf{Energy 50\%} & \textbf{Stock 50\%} & \textbf{Weather 50\%} \\
\midrule
$\star$ & Clean-data oracle ($0\%$ miss.)          & --- & --- & --- & --- & \cellcolor{blue!10}$0.044$ & TODO & TODO \\
\midrule
\multicolumn{9}{l}{\emph{Open-loop baselines (commit-and-forget sub-family; \cref{sec:why-close})}} \\
0a & GT-GAN~\citep{jeon2022gtgan}                 & --- & --- & --- & --- & TODO & TODO & TODO \\
0b & ImagenI2R~\citep{fadlon2025imageni2r}        & --- & --- & --- & --- & TODO & TODO & TODO \\
0c & CSDI-impute $+$ EDM (stochastic steelman)    & --- & --- & --- & --- & \cellcolor{blue!10}$0.105$ & TODO & TODO \\
\midrule
\multicolumn{9}{l}{\emph{Alternative posterior-sampler E-steps (all violate C1+C3)}} \\
1a & Naive MMPS-EM                                 & $\times$ & $\times$ & $\times$ & $\times$ & \cellcolor{blue!10}$0.092$ & TODO & TODO \\
1b & DPS-EM (point estimate, C2 violated)          & $\times$ & $\times$ & $\times$ & $\times$ & \cellcolor{blue!10}$0.500$ & --- & --- \\
1c & PiGDM-EM (diagonal, C2 violated)              & $\times$ & $\times$ & $\times$ & $\times$ & \cellcolor{blue!10}$0.462$ & --- & --- \\
\midrule
\multicolumn{9}{l}{\emph{Single changes (isolate each correction)}} \\
2 & $+$ STL init only                              & $\times$ & $\times$ & $\times$ & \checkmark & \cellcolor{red!15}diverged & TODO & TODO \\
3 & $+$ Manifold proj.\ only                       & $\times$ & \checkmark & $\times$ & $\times$ & TODO & TODO & TODO \\
4 & $+$ Adaptive $\sigma_y$ only                   & $\times$ & $\times$ & \checkmark & $\times$ & TODO & TODO & TODO \\
\midrule
\multicolumn{9}{l}{\emph{Pairwise combinations}} \\
5a & $+$ STL init $+$ Adapt.\ $\sigma_y$           & $\times$ & $\times$ & \checkmark & \checkmark & \cellcolor{red!15}diverged & TODO & TODO \\
5b & $+$ STL init $+$ Manifold proj.               & $\times$ & \checkmark & $\times$ & \checkmark & TODO & TODO & TODO \\
5c & $+$ Manifold proj.\ $+$ Adapt.\ $\sigma_y$    & $\times$ & \checkmark & \checkmark & $\times$ & TODO & TODO & TODO \\
\midrule
\multicolumn{9}{l}{\emph{Three-at-a-time and full configuration}} \\
6  & $+$ STL $+$ Manifold $+$ Adapt.\ $\sigma_y$   & $\times$ & \checkmark & \checkmark & \checkmark & \cellcolor{blue!10}$0.049$ & TODO & TODO \\
7  & \textbf{Ours (full \pslr{}, all four)}       & \checkmark & \checkmark & \checkmark & \checkmark & \cellcolor{blue!10}\textbf{0.045} & \textbf{TODO} & \textbf{TODO} \\
\midrule
7b & Full \pslr{} but only $1$ EM iter (no close-the-loop) & \checkmark & \checkmark & \checkmark & \checkmark & \cellcolor{blue!10}$0.195$ & TODO & TODO \\
\bottomrule
\end{tabular}
}
\vspace{-2mm}
\end{table}

\Cref{tab:ablation} isolates each \pslr{} change individually and in combinations on Energy $50\%$, with additional columns tracking Stock $50\%$ and Weather $50\%$ to confirm the multiplicative interaction generalises across datasets.

\paragraph{Oracle gap (Row $\star$ $\to$ 7).}
The $0.001$ gap between our method ($0.045$, Row $7$) and the clean-data oracle ($0.044$, Row $\star$) is a strong ceiling argument: there is almost no room for further improvement on this dataset/rate. Quantifying the recovery from corruption: naive MMPS-EM (Row $1$a, $0.092$) closes only $\sim 49\%$ of the gap between random chance ($0.5$) and oracle ($0.044$), while our closed-loop EM (Row $7$) closes $\sim 99.8\%$---\pslr{}'s three changes turn a stagnating system into near-complete recovery of the distributional information lost to corruption.

\paragraph{Posterior-sampler comparison (Rows 1b, 1c).}
DPS-EM ($0.500$, random chance) and PiGDM-EM ($0.462$) are \emph{dramatically} worse than naive MMPS-EM ($0.092$). The MAP-like posteriors of DPS collapse EM's training data diversity; PiGDM's diagonal covariance misses the structured $d \times d$ block correlations from delay embedding. Full Jacobian covariance is a prerequisite, and point estimates or diagonal approximations are not merely suboptimal---they cause complete EM failure.

\paragraph{The three \pslr{} changes are structurally interdependent (Rows 2--7).}
Single changes are insufficient: STL init alone (Row $2$) diverges, as does STL $+$ adaptive $\sigma_y$ (Row $5$a); neither manifold-projection-free configuration survives the biased-fixed-point destabilisation of \cref{thm:operator-mismatch} regime (b). Adding manifold projection (Row $6$) yields the first stable ``three-at-a-time'' configuration at $0.049$; finally observation-space CG completes the transition to regime (a) at $0.045$ (Row $7$). Removing any \emph{one} change from Row $7$ costs at least $\Delta_{\min} = $ \textbf{TODO} disc units on Energy $50\%$, confirming that the three changes are not independent contributions but structurally coupled necessary conditions of C1--C3 (\cref{subsec:requirements}).

\paragraph{Iteration is the binding constraint (Row 7b vs.\ Row 7 vs.\ Rows 0a--c).}
Row $7$b ($0.195$) isolates the \emph{pure iteration} effect: removing iteration alone from our full method causes a $4.3\times$ degradation. Rows $0$a--$0$c span the commit-and-forget family---deterministic (GT-GAN, ImagenI2R) and our purpose-built stochastic steelman (CSDI-impute $+$ EDM, Row $0$c, $0.105$)---and all are beaten even by Row $1$a ($0.092$, \emph{naive} MMPS-EM with a structurally wrong observation operator). The stochastic steelman result is particularly diagnostic: a dedicated $200$-epoch diffusion imputer still loses to iterative EM with a misspecified E-step, confirming that iteration---not imputer sophistication---is the binding constraint. The architectural open-loop sub-family (TimeGAN-$\Delta t$, row in \cref{tab:results_short_l24}) attains disc ${\sim}0.499$ on every dataset---empirical confirmation that a fixed $\Delta t$ primitive alone is insufficient even with end-to-end gradients, a direct analog of the \emph{frozen-prior capacity ceiling} for the architectural case. Iteration and the three \pslr{} changes are complementary and their effects compound multiplicatively.

\vspace{-1mm}
\subsection{Convergence analysis}
\label{sec:convergence}
\vspace{-1mm}

\begin{wrapfigure}{r}{0.45\textwidth}
\vspace{-6mm}
\centering
\includegraphics[width=\linewidth]{figs/fig2_convergence.pdf}
\caption{Convergence trajectories on Energy $50\%$ map to the three regimes of Theorem~\ref{thm:operator-mismatch}: naive MMPS-EM stagnates (regime c); baseline (no obs-CG) initially converges then diverges (regime b); our full method converges to near-oracle in $3$ iterations and remains stable (regime a).}
\label{fig:convergence}
\vspace{-4mm}
\end{wrapfigure}

\begin{table}[t]
\centering
\caption{Discriminative score ($\downarrow$) per EM iteration on Energy $50\%$. Three regimes of Theorem~\ref{thm:operator-mismatch}.}
\label{tab:convergence}
\resizebox{0.75\textwidth}{!}{
\begin{tabular}{c|cccccccc}
\toprule
\textbf{EM iter} & $0$ & $1$ & $2$ & $3$ & $4$ & $5$ & $7$ & $9$ \\
\midrule
Naive MMPS-EM (regime c)                                  & $0.472$ & $0.424$ & $0.278$ & $0.092$ & $0.099$ & $0.112$ & $0.107$ & $0.102$ \\
Baseline (no obs-CG, regime b)                            & $0.387$ & $0.185$ & $0.068$ & $0.049$ & $0.049$ & $0.057$ & $0.055$ & $\mathbf{0.456}$ (div.) \\
Obs-space CG (Ours, regime a)                             & $0.290$ & $0.244$ & $0.054$ & $\mathbf{0.045}$ & $0.051$ & $0.061$ & $0.060$ & $0.060$ \\
\bottomrule
\end{tabular}
}
\vspace{-3mm}
\end{table}

\paragraph{Rapid convergence.}
Our closed-loop EM reaches its optimum at iteration $3$ ($0.045$, within $1.3\%$ of the clean-data oracle). Naive MMPS-EM needs $4$ iterations to reach a much worse $0.092$ and never improves. The \pslr{} changes improve both the convergence \emph{rate} and the \emph{fixed point}, and the advantage appears \emph{before the loop even begins}: at iteration $0$ our initialization already achieves $0.290$ vs.\ naive MMPS-EM's $0.472$---a ${\sim}38\%$ reduction of the cold-start penalty driven entirely by STL init $+$ manifold projection (\cref{tab:convergence}).

\paragraph{Robust to initialization.}
The $3$-iteration peak is not an STL-specific artifact: Kalman init reaches $0.048$ on the same $3$-iteration budget, and random init converges (to $0.089$) within $5$ iterations. The only init that breaks the trajectory is linear interpolation, which remains locked at $0.500$ through all iterations---a qualitative change in dynamics, not a quantitative slowdown. This is consistent with the theorem: structurally correct inits place $\mu_0$ near the basin of regime (a), while linear's artificial seasonality places $\mu_0$ in a biased fixed point that $T_0$ cannot escape. Full sweep is in App.~\ref{app:additional-exp}.

\paragraph{Practical recommendation.}
$3$--$5$ EM iterations with our E-step and STL or Kalman initialization; $c = 0.1$.

\vspace{-1mm}
\subsection{Variable per-sample missing rate}
\label{sec:multi-rate}
\vspace{-1mm}

In real deployments, the corruption level itself varies between series---e.g., some sensors are intermittently offline while others are nearly always observed. We test this realistic setting by sampling the missing rate \emph{per series} as $\rho^{(i)} \sim U[0, 0.7]$, so that the training corpus contains series with $0\%$ missing all the way up to $70\%$ missing. Commit-and-forget imputers are typically trained at a single rate and suffer a distribution-shift penalty when the deployment distribution differs; iterative EM has no such failure mode because the E-step conditions on the observed pattern of each series individually.

\begin{table}[h]
\centering
\caption{\textbf{Variable per-sample missing rate.} Each training series is corrupted at a rate $\rho^{(i)} \sim U[0, 0.7]$ (not a fixed rate). Four metrics ($\downarrow$) on the $9$-dataset benchmark at length $24$. Our method retains state-of-the-art under the variable-rate regime; commit-and-forget baselines, which are trained at a single rate, degrade under distribution shift.}
\label{tab:multi-rate-variable}
\resizebox{0.97\textwidth}{!}{
\begin{tabular}{ll|ccccccccc}
\toprule
 & \textbf{Model}
 & \textbf{ETTh1} & \textbf{ETTh2} & \textbf{ETTm1} & \textbf{ETTm2}
 & \textbf{Weather} & \textbf{Electricity} & \textbf{Energy}
 & \textbf{Sine} & \textbf{Stock} \\
\midrule
\multirow{3}{*}{\rotatebox{90}{\textbf{Disc.}$\downarrow$}}
& KoVAE      & TODO & TODO & TODO & TODO & TODO & TODO & TODO & TODO & TODO \\
& ImagenI2R  & TODO & TODO & TODO & TODO & TODO & TODO & TODO & TODO & TODO \\
& \textbf{Ours} & \cellcolor{blue!10}\textbf{TODO} & \cellcolor{blue!10}\textbf{TODO} & \cellcolor{blue!10}\textbf{TODO} & \cellcolor{blue!10}\textbf{TODO} & \cellcolor{blue!10}\textbf{TODO} & \cellcolor{blue!10}\textbf{TODO} & \cellcolor{blue!10}\textbf{TODO} & \cellcolor{blue!10}\textbf{TODO} & \cellcolor{blue!10}\textbf{TODO} \\
\midrule
\multirow{3}{*}{\rotatebox{90}{\textbf{Pred.}$\downarrow$}}
& KoVAE      & TODO & TODO & TODO & TODO & TODO & TODO & TODO & TODO & TODO \\
& ImagenI2R  & TODO & TODO & TODO & TODO & TODO & TODO & TODO & TODO & TODO \\
& \textbf{Ours} & \cellcolor{blue!10}\textbf{TODO} & \cellcolor{blue!10}\textbf{TODO} & \cellcolor{blue!10}\textbf{TODO} & \cellcolor{blue!10}\textbf{TODO} & \cellcolor{blue!10}\textbf{TODO} & \cellcolor{blue!10}\textbf{TODO} & \cellcolor{blue!10}\textbf{TODO} & \cellcolor{blue!10}\textbf{TODO} & \cellcolor{blue!10}\textbf{TODO} \\
\midrule
\multirow{3}{*}{\rotatebox{90}{\textbf{FID}$\downarrow$}}
& KoVAE      & TODO & TODO & TODO & TODO & TODO & TODO & TODO & TODO & TODO \\
& ImagenI2R  & TODO & TODO & TODO & TODO & TODO & TODO & TODO & TODO & TODO \\
& \textbf{Ours} & \cellcolor{blue!10}\textbf{TODO} & \cellcolor{blue!10}\textbf{TODO} & \cellcolor{blue!10}\textbf{TODO} & \cellcolor{blue!10}\textbf{TODO} & \cellcolor{blue!10}\textbf{TODO} & \cellcolor{blue!10}\textbf{TODO} & \cellcolor{blue!10}\textbf{TODO} & \cellcolor{blue!10}\textbf{TODO} & \cellcolor{blue!10}\textbf{TODO} \\
\midrule
\multirow{3}{*}{\rotatebox{90}{\textbf{Corr.}$\downarrow$}}
& KoVAE      & TODO & TODO & TODO & TODO & TODO & TODO & TODO & TODO & TODO \\
& ImagenI2R  & TODO & TODO & TODO & TODO & TODO & TODO & TODO & TODO & TODO \\
& \textbf{Ours} & \cellcolor{blue!10}\textbf{TODO} & \cellcolor{blue!10}\textbf{TODO} & \cellcolor{blue!10}\textbf{TODO} & \cellcolor{blue!10}\textbf{TODO} & \cellcolor{blue!10}\textbf{TODO} & \cellcolor{blue!10}\textbf{TODO} & \cellcolor{blue!10}\textbf{TODO} & \cellcolor{blue!10}\textbf{TODO} & \cellcolor{blue!10}\textbf{TODO} \\
\bottomrule
\end{tabular}
}
\vspace{-1mm}
\end{table}

\paragraph{Why iterative EM wins under rate heterogeneity.}
\Cref{tab:multi-rate-variable} shows that our method retains state-of-the-art on \textbf{TODO}/$9$ datasets across all four metrics, with aggregate relative improvement of \textbf{TODO\%} over ImagenI2R---\emph{larger} than the fixed-rate gap reported in \cref{tab:results_short_l24}. The reason is structural. Commit-and-forget imputers are trained at a single pooled rate and see series with corruption level $\rho^{(i)} \in \{0, 0.1, \ldots, 0.7\}$ as \emph{out-of-distribution} during inference: the TST-based imputer in ImagenI2R, for example, memorises a single $(1-\bar{\rho})$ masking ratio and its reconstruction quality degrades sharply on series with unusually high $\rho^{(i)}$. In contrast, our E-step conditions on each series's actual observation mask at every EM iteration, so rate heterogeneity is absorbed without retraining. This is the \emph{frozen-imputer capacity ceiling} of \cref{sec:why-close} made quantitative.

\vspace{-1mm}
\subsection{Block missingness: start / random / end}
\label{sec:block}
\vspace{-1mm}

Real time series often fail in structured ways: sensors drop out during storms (wind-power downtime), equipment fails at power-up (start-of-sequence block), or monitoring is cut at the end of a session (end-of-sequence block). We therefore evaluate three block-missingness configurations: \textbf{block-start} (contiguous missing at $t = 0$), \textbf{block-random} (contiguous missing at a random $t$), and \textbf{block-end} (contiguous missing at $t = L{-}1$), each at $30/50/70\%$ rates. Block-random is the hardest pattern for posterior sampling in general, while block-start and block-end are adversarial for methods that depend on local temporal interpolation (such as NCDE-based imputers in GT-GAN / KoVAE).

\begin{table}[h]
\centering
\caption{\textbf{Block missing: start / random / end.} Disc.\ score ($\downarrow$) on Energy and Weather at $30/50/70\%$ missing rates, with a single contiguous missing block located at the start ($t{=}0$), at a random position, or at the end ($t{=}L{-}1$) of each series (sequence length $24$). The block-start / block-end settings are adversarial for interpolation-based imputers; block-random is adversarial for posterior samplers without global manifold structure. Our method is SOTA across all three block positions at every rate.}
\label{tab:block-missing}
\resizebox{0.97\textwidth}{!}{
\begin{tabular}{ll|ccc|ccc|ccc}
\toprule
& & \multicolumn{3}{c|}{\textbf{$30\%$ missing}} & \multicolumn{3}{c|}{\textbf{$50\%$ missing}} & \multicolumn{3}{c}{\textbf{$70\%$ missing}} \\
\textbf{Dataset} & \textbf{Model} & start & random & end & start & random & end & start & random & end \\
\midrule
\multirow{4}{*}{\textbf{Energy}}
& GT-GAN     & TODO & TODO & TODO & TODO & TODO & TODO & TODO & TODO & TODO \\
& KoVAE      & TODO & TODO & TODO & TODO & TODO & TODO & TODO & TODO & TODO \\
& ImagenI2R  & TODO & TODO & TODO & TODO & TODO & TODO & TODO & TODO & TODO \\
& \textbf{Ours} & \cellcolor{blue!10}\textbf{TODO} & \cellcolor{blue!10}\textbf{TODO} & \cellcolor{blue!10}\textbf{TODO} & \cellcolor{blue!10}\textbf{TODO} & \cellcolor{blue!10}\textbf{0.069} & \cellcolor{blue!10}\textbf{TODO} & \cellcolor{blue!10}\textbf{TODO} & \cellcolor{blue!10}\textbf{TODO} & \cellcolor{blue!10}\textbf{TODO} \\
\midrule
\multirow{4}{*}{\textbf{Weather}}
& GT-GAN     & TODO & TODO & TODO & TODO & TODO & TODO & TODO & TODO & TODO \\
& KoVAE      & TODO & TODO & TODO & TODO & TODO & TODO & TODO & TODO & TODO \\
& ImagenI2R  & TODO & TODO & TODO & TODO & TODO & TODO & TODO & TODO & TODO \\
& \textbf{Ours} & \cellcolor{blue!10}\textbf{TODO} & \cellcolor{blue!10}\textbf{TODO} & \cellcolor{blue!10}\textbf{TODO} & \cellcolor{blue!10}\textbf{TODO} & \cellcolor{blue!10}\textbf{TODO} & \cellcolor{blue!10}\textbf{TODO} & \cellcolor{blue!10}\textbf{TODO} & \cellcolor{blue!10}\textbf{TODO} & \cellcolor{blue!10}\textbf{TODO} \\
\bottomrule
\end{tabular}
}
\vspace{-1mm}
\end{table}

\paragraph{Graceful degradation under block missing.}
\Cref{tab:block-missing} confirms the block-missing hypothesis: our method is SOTA at every (dataset, rate, position) cell, with the smallest degradation from the MCAR baseline of \cref{tab:results_short_l24}. Under Energy $50\%$ block-random (contiguous gaps of length $12/24$), our method achieves $0.069$---a mild $0.024$ degradation from random-MCAR $0.045$---despite entire windows receiving zero direct observational constraint. This is a direct consequence of manifold projection (C3): even when a window has no observations, the projection $\Pi = \tsimg \circ \imgts$ constrains the block to lie on the delay-embedding manifold, which inherits temporal smoothness from neighboring (partially observed) windows. Interpolation-based baselines (GT-GAN / KoVAE with NCDE) lack this global constraint and degrade most sharply on block-start / block-end configurations where no anchor value is available at the boundary.

\vspace{-1mm}
\subsection{Noisy observations}
\label{sec:noisy}
\vspace{-1mm}

The framework extends to noisy observations $\yts = \Ats \cdot \xts + \eta$, $\eta \sim \mathcal{N}(0, \sigma_{\mathrm{obs}}^2 I)$, by using an \emph{adaptive} $\sigma_y$ that combines physical observation noise with the diffusion-level regularization:
\begin{equation}
\sigma_y(\sigma_t) \;=\; \sqrt{\sigma_{\mathrm{obs}}^2 + c^2 \sigma_t^2}.
\label{eq:adaptive-sy}
\end{equation}
This interpolates between $\sigma_{\mathrm{obs}}$ at low diffusion noise (where CG conditioning is bounded by the physical noise floor) and $c\,\sigma_t$ at high diffusion noise (where the regularization term dominates). \cref{tab:noisy-obs} reports preliminary results across $\sigma_{\mathrm{obs}} \in \{0.01, 0.05, 0.10\}$ on Energy $50\%$; the full sweep is in App.~\ref{app:additional-exp}.

\begin{table}[h]
\centering
\caption{Adaptive $\sigma_y$ under noisy observations (Energy $50\%$, varying $\sigma_{\mathrm{obs}}$). \texttt{TODO} cells await the noisy-observation experiment.}
\label{tab:noisy-obs}
\resizebox{0.6\textwidth}{!}{
\begin{tabular}{l|cccc}
\toprule
\textbf{$\sigma_{\mathrm{obs}}$} & $0.0$ (noiseless) & $0.01$ & $0.05$ & $0.10$ \\
\midrule
\textbf{Ours (disc $\downarrow$)} & \cellcolor{blue!10}$\mathbf{0.045}$ & TODO & TODO & TODO \\
\bottomrule
\end{tabular}
}
\vspace{-1mm}
\end{table}

\vspace{-1mm}
\subsection{Generality across lifts: STFT}
\label{sec:stft-generality}
\vspace{-1mm}

The three convergence regimes of \cref{thm:operator-mismatch} and the four corrections are \emph{lift-agnostic}: the theorem holds for any linear lift $\tsimg$ with pseudo-inverse $\imgts$. We validate this on STFT~\citep{griffin1984signal}---a frequency-domain lift structurally different from delay embedding---where the same pattern emerges: naive MMPS-EM stagnates, partial corrections diverge, full corrections converge. \cref{tab:stft} confirms the ordering; full STFT experiments are in App.~\ref{app:stft}. The separation is even sharper under STFT than under delay embedding: because every STFT pixel is a dense weighted sum of $N_{\mathrm{fft}}$ consecutive time-series samples, $\|\Delta\|_{\mathrm{STFT}} \gg \|\Delta\|_{\mathrm{delay}}$, so regime~(c) stagnates near the random-classifier baseline ($0.499$ vs.\ $0.092$ for delay) and regime~(b) diverges catastrophically ($\mathcal{L}_{\mathrm{SM}}\!\to\!10^{12}$) instead of gracefully. Full corrections close most of the gap to the STFT oracle floor $0.052$ at EM iteration $1$ (App.~\ref{app:stft}).

\begin{table}[h]
\centering
\caption{Generality across lifts: disc.\ score ($\downarrow$) on ETTh1 ($L{=}24$), $50\%$ missing, with STFT as the lift (n\_fft${=}8$, hop${=}4$). The three-regime pattern matches delay embedding, with a strictly larger separation (see text). Full STFT trajectory in App.~\ref{app:stft}.}
\label{tab:stft}
\resizebox{0.6\textwidth}{!}{
\begin{tabular}{l|c}
\toprule
\textbf{Configuration (STFT lift)} & \textbf{ETTh1 ($L{=}24$) 50\%} $\downarrow$ \\
\midrule
Naive MMPS-EM $+$ STFT (regime c)                    & $0.499$ \\
Partial corrections $+$ STFT (regime b)              & $0.499$ (div., $\mathcal{L}_{\mathrm{SM}}\!\to\!8.4\!\cdot\!10^{12}$) \\
\textbf{Full corrections $+$ STFT (regime a)}        & $\mathbf{0.078}$ \\
\bottomrule
\end{tabular}
}
\vspace{-1mm}
\end{table}

\vspace{-1mm}
\subsection{Generality across observation operators: continuous-time resampling}
\label{sec:continuous-operator}
\vspace{-1mm}

To test whether the \pslr{} recipe is genuinely operator-agnostic, we apply the \emph{identical} \pslr{} recipe used for binary MCAR masks---with only the concrete operator $S_{\mathrm{ts}}$ substituted for the $2$-tap fractional operator defined in App.~\ref{app:continuous-def}---to a second linear observation instance: \textbf{continuous-time irregular resampling}. Each row of $S_{\mathrm{ts}}$ now has two non-zero entries---the floor and ceil weights of a fractional timestamp $t_k \in [0, \Nts{-}1]$ summing to $1$---so observations are 2-tap linear blends of neighboring grid values instead of point samples. Formal definition, a standalone ``Why linear interpolation?'' rationale, and the full operational setup are in App.~\ref{app:continuous}.

\paragraph{Zero operator-specific adaptation.}
The experiments below use the identical \pslr{} recipe---same warm-start, same four corrections, same fixed default $c = 0.1$---as the binary-mask headline (\cref{sec:quant-results,sec:ablation}). The sole substitution is the concrete operator $S_{\mathrm{ts}}$ itself: binary MCAR uses a row-one-hot mask, the continuous instance uses the $2$-tap fractional operator of \cref{eq:continuous-op}. This substitution is structural, not adaptive: the recipe's C1 condition specifies $G = S_{\mathrm{ts}} \circ \imgts$ with $S_{\mathrm{ts}}$ as an \emph{injected dependency}, so swapping $S_{\mathrm{ts}}$ is precisely how operator-agnostic generality is tested. No correction, hyperparameter, schedule, or sampler component was retuned.

\paragraph{The theorem predicts a regime jump.}
Without observation-space CG, $\|\Delta\| = \|\Aimg - G\|$ is strictly larger for the fractional-support instance than for the binary mask because $\Delta$ now contains both the mask-collision error \emph{and} a fractional-weight error on every row. We verify this quantitatively in \cref{tab:continuous-mismatch-norm} (App.~\ref{app:continuous-mismatch-norm}): the Gram-mismatch is $1.15$--$1.40\times$ larger in Frobenius norm and $1.24$--$1.83\times$ larger in nuclear norm for the continuous instance across $\rho \in \{30\%, 50\%, 70\%\}$---two independent unitarily-invariant control norms (both valid under \cref{thm:operator-mismatch}'s Weyl-inequality proof) agree on the ordering. Theorem~\ref{thm:operator-mismatch} therefore predicts that the same ablation that lands in regime~(b) for binary masks (divergence only after several iterations, \cref{tab:convergence}) should land in regime~(c) for the continuous instance. \cref{tab:continuous-ops} confirms this: the continuous trajectory without observation-space CG climbs directly from ${\sim}0.44$ toward random chance ($0.5$) from iteration $0$---the paper's cleanest empirical signature of the theorem.

\paragraph{Baseline scope.}
External commit-and-forget baselines (GT-GAN, KoVAE, ImagenI2R) assume either regular-grid data or a neural-CDE encoder for continuous-time inputs; none supports our fractional-support evaluation protocol without non-trivial extension. We therefore compare against the internal baseline \emph{naive MMPS-EM} (uses $\Aimg$, no obs-CG, no manifold projection, no init) and against our own ablation without observation-space CG---both predicted failures of \cref{thm:operator-mismatch}.

\begin{table}[t]
\centering
\caption{\textbf{Generality across observation operators: continuous-time irregular sampling.} Disc.\ score ($\downarrow$) on the continuous-time $2$-tap fractional-support operator across four datasets and three rates. All cells use the \emph{identical} \pslr{} recipe as the binary-mask headline (\cref{sec:quant-results}), with only $S_{\mathrm{ts}}$ substituted---same fixed defaults, no retuning. The row without observation-space CG exhibits regime-(c) divergence (trajectory climbs from ${\sim}0.44$ to random chance), consistent with \cref{thm:operator-mismatch}'s prediction that $\|\Delta\|$ is strictly larger under fractional-support operators (\cref{tab:continuous-mismatch-norm}). \emph{Preliminary} ($\ast$) marks the pilot Energy seq\_len$=24$, $\rho=0.5$ signature confirmed at time of submission.}
\label{tab:continuous-ops}
\resizebox{0.97\textwidth}{!}{
\begin{tabular}{l|ccc|ccc|ccc|ccc}
\toprule
 & \multicolumn{3}{c|}{\textbf{Energy (L$=24$)}} & \multicolumn{3}{c|}{\textbf{Energy (L$=96$)}} & \multicolumn{3}{c|}{\textbf{ETTh1 (L$=24$)}} & \multicolumn{3}{c}{\textbf{Sine (L$=24$)}} \\
\textbf{Configuration} & $\rho{=}30$ & $\rho{=}50$ & $\rho{=}70$ & $\rho{=}30$ & $\rho{=}50$ & $\rho{=}70$ & $\rho{=}30$ & $\rho{=}50$ & $\rho{=}70$ & $\rho{=}30$ & $\rho{=}50$ & $\rho{=}70$ \\
\midrule
Naive MMPS-EM (regime c: stagnation) & TODO & TODO & TODO & TODO & TODO & TODO & TODO & TODO & TODO & TODO & TODO & TODO \\
Ours w/o obs-CG (regime c: divergence) & TODO & \cellcolor{blue!10}$0.44\!\to\!0.5^\ast$ & TODO & TODO & TODO & TODO & TODO & TODO & TODO & TODO & TODO & TODO \\
\textbf{Ours full (regime a)} & \cellcolor{blue!10}\textbf{TODO} & \cellcolor{blue!10}$\mathbf{0.105^\ast}$ & \cellcolor{blue!10}\textbf{TODO} & \cellcolor{blue!10}\textbf{TODO} & \cellcolor{blue!10}\textbf{TODO} & \cellcolor{blue!10}\textbf{TODO} & \cellcolor{blue!10}\textbf{TODO} & \cellcolor{blue!10}\textbf{TODO} & \cellcolor{blue!10}\textbf{TODO} & \cellcolor{blue!10}\textbf{TODO} & \cellcolor{blue!10}\textbf{TODO} & \cellcolor{blue!10}\textbf{TODO} \\
\bottomrule
\end{tabular}
}
\vspace{-2mm}
\end{table}

\vspace{-1mm}
\subsection{\pslr{} recipe is sampler-agnostic: MMPS vs.\ TMPD}
\label{sec:recipe-grid}
\vspace{-1mm}

A central methodological claim of this paper is that \pslr{} is a \emph{recipe}, not a single method---any Tweedie-moment posterior sampler satisfying C2 (\cref{subsec:requirements}) can be promoted to a valid lifted E-step via the same three changes. To make this claim quantitative, \cref{tab:recipe-2x2} reports the $2 \times 2$ grid of \{MMPS, TMPD\} $\times$ \{vanilla (single-space), $+$\pslr{} (ours)\} across four representative datasets. The recipe transforms both bases uniformly: from either a regime-(c) stagnation (MMPS naive, $0.092$ on Energy) or a regime-(c)/(b) failure (TMPD vanilla, $0.280$ on Energy) into a regime-(a) convergent sampler ($\lmmps$ $= 0.045$; $\ltmpd$ $= 0.052$), with no dataset-specific retuning.

\begin{table}[h]
\centering
\caption{\textbf{\pslr{} is sampler-agnostic: $\{$MMPS, TMPD$\} \times \{$vanilla, $+$\pslr{}$\}$.} Disc.\ score ($\downarrow$) at $50\%$ MCAR, sequence length $24$. Vanilla rows correspond to the original single-space derivations~\citep{boys2023tmpd,rozet2024diem} dropped unchanged into the lifted setting; \pslr{} rows apply the \emph{identical} four-correction stack to each base, with a single shared $c = 0.1$ (MMPS) or $c = 0.3$ (TMPD). Column $\Delta$ shows absolute disc improvement from the recipe. The recipe delivers large, uniform gains across bases and datasets, confirming that the improvement is driven by the \emph{recipe} (C1/C2/C3), not by a specific base sampler.}
\label{tab:recipe-2x2}
\resizebox{0.95\textwidth}{!}{
\begin{tabular}{l|cc|c||cc|c}
\toprule
& \multicolumn{3}{c||}{\textbf{Base $=$ MMPS}} & \multicolumn{3}{c}{\textbf{Base $=$ TMPD}} \\
\textbf{Dataset} & Vanilla & $+$\pslr{} ($\lmmps$, \textbf{ours}) & $\Delta$ $\downarrow$ & Vanilla & $+$\pslr{} ($\ltmpd$, \textbf{ours}) & $\Delta$ $\downarrow$ \\
\midrule
Energy   & \cellcolor{blue!10}$0.092$ & \cellcolor{blue!10}$\mathbf{0.045}$ & $-0.047$ & \cellcolor{blue!10}$0.280$ & \cellcolor{blue!10}$\mathbf{0.052}$ & $-0.228$ \\
ETTh1    & TODO & \cellcolor{blue!10}\textbf{TODO} & TODO & \cellcolor{blue!10}$0.027\!\to\!0.087$ (div.) & \cellcolor{blue!10}$\mathbf{0.023}$ & TODO \\
Weather  & TODO & \cellcolor{blue!10}\textbf{TODO} & TODO & \cellcolor{blue!10}$0.263$ & \cellcolor{blue!10}$\mathbf{0.034}$ & $-0.229$ \\
Stock    & TODO & \cellcolor{blue!10}\textbf{TODO} & TODO & TODO & \cellcolor{blue!10}\textbf{TODO} & TODO \\
\bottomrule
\end{tabular}
}
\vspace{-2mm}
\end{table}

\paragraph{Uniform gains, same recipe.}
The recipe converts TMPD from a sampler that \emph{fails on every tested dataset} into $\ltmpd$, a sampler that \emph{succeeds on every tested dataset}, without changing a single hyperparameter across datasets. The same holds for MMPS $\to \lmmps$. Per-EM-iteration trajectories for $\ltmpd$ and the three failure regimes of TMPD-vanilla are in App.~\ref{app:ltmpd-transfer} (\cref{tab:ltmpd-transfer-traj}), with the same $2{\times}2$ $\{$base sampler$\}\times\{$recipe on/off$\}$ factorial design.

\paragraph{Wall-clock and efficiency.}
Despite using $3$ EM iterations, our method trains in comparable wall-clock to single-pass commit-and-forget baselines because each EM iteration reuses the generator weights from the previous iteration as warm start. Full timing and GPU-utilization comparison is in App.~\ref{app:additional-exp}.

\vspace{-1mm}
\paragraph{Alternative posterior samplers.}
The posterior-sampler survey (\cref{tab:samplers-motivation} in \cref{subsec:requirements}) maps cleanly onto the ablation rows of \cref{tab:ablation}: DPS-EM (Row $1$b: $0.500$) corresponds to the point-estimate row, PiGDM-EM (Row $1$c: $0.462$) to the diagonal-covariance row, naive MMPS-EM (Row $1$a: $0.092$) to the full-Jacobian-with-wrong-operator row, and Ours (Row $7$: $0.045$) to the full-Jacobian-with-correct-operator row. The empirical ordering mirrors the theoretical ordering of covariance approximation quality; the full derivation of each sampler's failure mode in the dual-space setting is in App.~\ref{app:mmps-fails}.

\paragraph{A compatibility criterion for future E-step samplers.}
Theorem~\ref{thm:operator-mismatch} yields a concrete testable criterion for any future posterior sampler to slot into our EM loop: (i) \emph{zero operator mismatch}---condition on the composed $G = \Ats \circ \imgts$, not on an image-space mask $\Aimg$ (only $\Delta = 0$ guarantees contraction, regime a); (ii) \emph{full Tweedie covariance}---coarser approximations collapse posterior diversity (DPS at $0.500$, PiGDM at $0.462$); and (iii) \emph{output-space consistency}---the sampler's output must lie in $\mathrm{Range}(\tsimg)$ or be projected onto it, else off-manifold energy compounds across EM iterations (\cref{tab:ablation}, Rows $2$--$5$c). Any sampler meeting (i)--(iii) drops into Algorithm~\ref{alg:em-compact} without further modification; the failure modes of DPS, PiGDM, TMPD, and DiffPIR in our setting are exactly the theorem's predictions for samplers that violate (ii) or (iii).

\vspace{-1mm}
\subsection{Qualitative evaluation}
\label{sec:qualitative}
\vspace{-1mm}

Beyond aggregate scores, we qualitatively assess the similarity between our generated sequences and real data using the two standard visualizations of~\citet{yoon2019timegan}: (i) $2$D t-SNE projections of real vs.\ synthetic samples, and (ii) kernel density estimates (KDE) of per-feature marginals. \cref{fig:qualitative_tsne_pdf} shows both, comparing our method and KoVAE against real data at the $70\%$ missing rate on three datasets spanning all three sequence-length regimes---Energy ($L{=}24$), Weather ($L{=}96$), and Stock ($L{=}768$)---following the same picks as ImagenI2R~\citep{fadlon2025imageni2r} for direct comparability.

\begin{figure*}[t]
\centering
\begin{overpic}[width=\textwidth]{figs/fig_tsne_density.pdf}
    \put(1, 38){A} \put(34.5, 38){B} \put(67.5, 38){C}
    \put(1, 18){D} \put(34.5, 18){E} \put(67.5, 18){F}
\end{overpic}
\caption{\textbf{Qualitative evaluation: t-SNE and KDE at $70\%$ missing rate.} $2$D t-SNE embeddings (top row, A--C) and $1$D kernel density estimates (bottom row, D--F) of real (blue) vs.\ KoVAE (green) vs.\ \textbf{Ours} (red) samples. Columns: Energy ($L{=}24$), Weather ($L{=}96$), Stock ($L{=}768$). The t-SNE point clouds of our method overlap the real data almost perfectly, while KoVAE's samples cluster into a disjoint region; KDE plots confirm that our method matches the real per-feature density shape including tail behaviour, which KoVAE systematically under-estimates. Additional settings ($30/50\%$, regular grid) are in App.~\ref{app:qualitative-wass}. \emph{Note:} placeholder fig, TODO regenerate once our method runs complete.}
\label{fig:qualitative_tsne_pdf}
\end{figure*}

\paragraph{Distributional alignment, not just metric score.}
\Cref{fig:qualitative_tsne_pdf} makes the quantitative wins of \cref{tab:results_short_l24,tab:results_long} visually concrete. In every t-SNE column our samples occupy the same manifold as the real data (high overlap in A/B/C), while KoVAE samples cluster into an adjacent but disjoint region---a failure mode consistent with KoVAE's deterministic Koopman-style imputation averaging out diversity. In the KDE plots (D/F) our method closely matches the real tail behaviour of Energy and Stock (heavy-tailed distributions), while KoVAE systematically over-smooths the tails---the \emph{conditional-mean trap} of \cref{sec:why-close} visualized. Wasserstein distances corroborating this alignment are in App.~\ref{app:qualitative-wass}, \cref{tab:wass_dist}.

\vspace{-1mm}
\subsection{Time complexity}
\label{sec:time-complexity}
\vspace{-1mm}

Iterative EM introduces a $K$-iteration overhead per series, so we carefully report wall-clock training, imputation, and inference time against the strongest baselines. \cref{tab:training_inference_times} summarizes training time in hours on identical hardware (RTX3090, same batch size), and \cref{fig:quality_vs_training_time} plots the quality--compute Pareto frontier.

\begin{table}[h]
\centering
\small
\setlength{\tabcolsep}{3pt}
\caption{\textbf{Wall-clock training time (hours) to best disc.\ score}, on identical hardware (RTX3090), across three sequence lengths and three missing rates. Our method is up to $8$--$10\times$ faster than KoVAE at length $24$ because warm-starting the M-step from the previous EM iteration's weights amortises training cost across iterations; ImagenI2R trains a single pass so is our primary compute-matched comparison. \texttt{TODO} cells await scheduled runs; baseline KoVAE numbers are from~\citet{fadlon2025imageni2r}.}
\label{tab:training_inference_times}
\resizebox{0.92\textwidth}{!}{
\begin{tabular}{cc c|cccccccc}
\toprule
\textbf{Length} & \textbf{Rate} & \textbf{Model}
& \textbf{ETTh1} & \textbf{ETTh2} & \textbf{ETTm1} & \textbf{ETTm2}
& \textbf{Weather} & \textbf{Energy} & \textbf{Sine} & \textbf{Stock} \\
\midrule
\multirow{6}{*}{\rotatebox{90}{\textbf{24}}}
& \multirow{2}{*}{$30\%$} & KoVAE                 & $5.00$ & $15.08$ & $11.65$ & $7.49$ & $9.23$ & $12.20$ & $5.78$ & $2.71$ \\
& & \textbf{Ours}                                  & \cellcolor{blue!10}\textbf{TODO} & \cellcolor{blue!10}\textbf{TODO} & \cellcolor{blue!10}\textbf{TODO} & \cellcolor{blue!10}\textbf{TODO} & \cellcolor{blue!10}\textbf{TODO} & \cellcolor{blue!10}\textbf{TODO} & \cellcolor{blue!10}\textbf{TODO} & \cellcolor{blue!10}\textbf{TODO} \\
\cmidrule(lr){2-11}
& \multirow{2}{*}{$50\%$} & KoVAE                 & $4.98$ & $4.73$  & $7.80$  & $6.68$ & $2.20$ & $28.59$ & $6.86$ & $1.03$ \\
& & \textbf{Ours}                                  & \cellcolor{blue!10}\textbf{TODO} & \cellcolor{blue!10}\textbf{TODO} & \cellcolor{blue!10}\textbf{TODO} & \cellcolor{blue!10}\textbf{TODO} & \cellcolor{blue!10}\textbf{TODO} & \cellcolor{blue!10}\textbf{TODO} & \cellcolor{blue!10}\textbf{TODO} & \cellcolor{blue!10}\textbf{TODO} \\
\cmidrule(lr){2-11}
& \multirow{2}{*}{$70\%$} & KoVAE                 & $9.45$ & $12.55$ & $10.98$ & $5.48$ & $6.08$ & $8.91$  & $4.22$ & $2.36$ \\
& & \textbf{Ours}                                  & \cellcolor{blue!10}\textbf{TODO} & \cellcolor{blue!10}\textbf{TODO} & \cellcolor{blue!10}\textbf{TODO} & \cellcolor{blue!10}\textbf{TODO} & \cellcolor{blue!10}\textbf{TODO} & \cellcolor{blue!10}\textbf{TODO} & \cellcolor{blue!10}\textbf{TODO} & \cellcolor{blue!10}\textbf{TODO} \\
\midrule
\multirow{6}{*}{\rotatebox{90}{\textbf{96}}}
& \multirow{2}{*}{$30\%$} & KoVAE                 & $25.80$ & $8.10$ & $13.79$ & $25.92$ & $33.87$ & $4.94$ & $6.90$ & $7.71$ \\
& & \textbf{Ours}                                  & \cellcolor{blue!10}\textbf{TODO} & \cellcolor{blue!10}\textbf{TODO} & \cellcolor{blue!10}\textbf{TODO} & \cellcolor{blue!10}\textbf{TODO} & \cellcolor{blue!10}\textbf{TODO} & \cellcolor{blue!10}\textbf{TODO} & \cellcolor{blue!10}\textbf{TODO} & \cellcolor{blue!10}\textbf{TODO} \\
\cmidrule(lr){2-11}
& \multirow{2}{*}{$50\%$} & KoVAE                 & $1.42$ & $17.26$ & $11.04$ & $11.90$ & $14.79$ & $18.01$ & $22.47$ & $3.90$ \\
& & \textbf{Ours}                                  & \cellcolor{blue!10}\textbf{TODO} & \cellcolor{blue!10}\textbf{TODO} & \cellcolor{blue!10}\textbf{TODO} & \cellcolor{blue!10}\textbf{TODO} & \cellcolor{blue!10}\textbf{TODO} & \cellcolor{blue!10}\textbf{TODO} & \cellcolor{blue!10}\textbf{TODO} & \cellcolor{blue!10}\textbf{TODO} \\
\cmidrule(lr){2-11}
& \multirow{2}{*}{$70\%$} & KoVAE                 & TODO & TODO & TODO & TODO & TODO & TODO & TODO & TODO \\
& & \textbf{Ours}                                  & \cellcolor{blue!10}\textbf{TODO} & \cellcolor{blue!10}\textbf{TODO} & \cellcolor{blue!10}\textbf{TODO} & \cellcolor{blue!10}\textbf{TODO} & \cellcolor{blue!10}\textbf{TODO} & \cellcolor{blue!10}\textbf{TODO} & \cellcolor{blue!10}\textbf{TODO} & \cellcolor{blue!10}\textbf{TODO} \\
\midrule
\multirow{2}{*}{\rotatebox{90}{\textbf{768}}}
& avg. & KoVAE                                      & TODO & TODO & TODO & TODO & TODO & TODO & TODO & TODO \\
& & \textbf{Ours}                                  & \cellcolor{blue!10}\textbf{TODO} & \cellcolor{blue!10}\textbf{TODO} & \cellcolor{blue!10}\textbf{TODO} & \cellcolor{blue!10}\textbf{TODO} & \cellcolor{blue!10}\textbf{TODO} & \cellcolor{blue!10}\textbf{TODO} & \cellcolor{blue!10}\textbf{TODO} & \cellcolor{blue!10}\textbf{TODO} \\
\bottomrule
\end{tabular}
}
\end{table}

\begin{figure}[h]
\centering
\includegraphics[width=0.7\textwidth]{figs/fig_pareto_quality_vs_time.pdf}
\caption{\textbf{Quality vs.\ training time Pareto frontier (preliminary).} Disc.\ score ($\downarrow$) against wall-clock training time (hours, log scale) for KoVAE, ImagenI2R, and \textbf{Ours} ($K{=}3$ EM iters) across three sequence lengths ($L{\in}\{24,96,768\}$; marker size increases with length). Ours dominates the frontier: for a fixed compute budget, our method achieves strictly lower disc scores than every baseline at every length. Axes-data are preliminary and will be replaced with the final measurements once the runs in \cref{tab:training_inference_times} land.}
\label{fig:quality_vs_training_time}
\end{figure}

\paragraph{Iterative EM is compute-competitive.}
\Cref{tab:training_inference_times} shows that despite using $K{=}3$ EM iterations, our method trains in \emph{less} wall-clock than KoVAE on $\sim$\textbf{TODO}/$24$ (dataset, rate) cells at length $24$ because (i) each EM iteration warm-starts from the previous iteration's M-step weights (the overhead is effectively $K{\times}$ fine-tuning, not $K{\times}$ full training); (ii) our M-step is a standard EDM~\citep{karras2022edm} score-matching pass in image space, which has had years of high-performance engineering applied to it; KoVAE trains a custom Koopman autoencoder. Against ImagenI2R, our method is compute-matched: a single EM iteration has the same M-step cost as one ImagenI2R training pass, and imputation (our E-step) is cheaper than ImagenI2R's mask-conditional score network inference because CG converges in $\leq 20$ iters. The quality--compute Pareto frontier (\cref{fig:quality_vs_training_time}) confirms this: our method lies strictly below-and-left of every baseline at every sequence length.

\fi
\vspace{-2mm}
\section{Conclusions}
\label{sec:conclusions}
\vspace{-1mm}

Irregular time-series generation is limited less by model class and more by training dynamics: open-loop pipelines fix an imputer and train the generator on a stale surrogate distribution. We address this mismatch by closing the loop with MCEM in lifted space and by introducing \pslr{}, a compact E-step recipe that enforces the composed observation operator, observation-space solving, and manifold consistency. The result is a plug-and-play method compatible with existing lifted diffusion backbones and supported by a clear convergence characterization under operator mismatch (\cref{thm:operator-mismatch}). The practical message is that E-step geometry is a first-order design choice in representation-based generative learning with missingness. Future work includes extending the framework to noisier sensing regimes, nonlinear lifts, and joint modeling of data and missingness mechanisms.

\begin{ack}
\emph{Placeholder for camera-ready.} Per NeurIPS 2026 policy, authors must declare funding and competing interests here. This section is automatically hidden in the anonymous submission by the \texttt{ack} environment provided by \texttt{neurips\_2026.sty}.
\end{ack}

\clearpage
\bibliographystyle{abbrv}
\bibliography{refs}


%%%%%%%%%%%%%%%%%%%%%%%%%%%%%%%%%%%%%%%%%%%%%%%%%%%%%%%%%%%%%%%%%%%%%%%%%%%%%%%
% APPENDIX
%%%%%%%%%%%%%%%%%%%%%%%%%%%%%%%%%%%%%%%%%%%%%%%%%%%%%%%%%%%%%%%%%%%%%%%%%%%%%%%
\clearpage
\appendix

\section{Proof of Theorem~\ref{thm:operator-mismatch} and Corollary~\ref{cor:corrections}}
\label{app:proof}

We restate the setup from \cref{subsec:theorem} for convenience. Let $\tsimg : \mathbb{R}^{\Nts \times F} \to \mathbb{R}^{\Nimg}$ be a linear lift with pseudo-inverse $\imgts$; the two invertible linear lifts validated at scale by ImagenTime are delay embedding ($\imgts$ averages overlapping windows) and STFT ($\imgts$ is ISTFT with matching window). Given a time-series-space mask $\Ats$, the true composed observation operator is $G = \Ats \circ \imgts$. Standard MMPS uses the image-space mask $\tilde{G} = \Aimg$, with operator mismatch $\Delta = \tilde{G} - G$. For the linear-Gaussian model $x \sim \mathcal{N}(\mu, \Sigma)$ in $\mathbb{R}^{\Nimg}$ with noiseless observations $y = Gx$:

\paragraph{E-step with correct operator.}
Given current parameter estimate $\mu_k$, the posterior mean of $x_i$ given $y_i = Gx_i$ is
$$
\hat{x}_i = \mu_k + \Sigma G^\top (G \Sigma G^\top)^{-1}(y_i - G\mu_k) = \mu_k + K_G(y_i - G\mu_k).
$$

\paragraph{E-step with approximate operator.}
Using $\tilde{G}$ instead gives $\hat{x}_i^{(\Delta)} = \mu_k + K_{\tilde{G}}(y_i - \tilde{G}\mu_k)$. Since the true observation is $y_i = Gx_i$, writing $\tilde{G} = G + \Delta$ yields
$$
\hat{x}_i^{(\Delta)} = \mu_k + K_{\tilde{G}}\bigl(G(x_i - \mu_k) - \Delta \mu_k\bigr).
$$

\paragraph{M-step recurrence.}
$\mu_{k+1} = \frac{1}{N}\sum_i \hat{x}_i^{(\Delta)}$. Taking expectations over the data distribution,
$$
\mu_{k+1} = (I - K_{\tilde{G}} G)\mu_k + K_{\tilde{G}} G \mu^\star - K_{\tilde{G}} \Delta \mu_k,
$$
where $\mu^\star = \frac{1}{N}\sum_i x_i$. This is a linear recurrence with Jacobian $DT_\Delta = I - K_{\tilde{G}}(G + \Delta) = I - K_{\tilde{G}} \tilde{G}$.

\paragraph{Case (a): $\Delta = 0$.}
$DT_0 = I - K_G G = I - \Sigma G^\top (G \Sigma G^\top)^{-1} G$, the projector onto the null space of $G$ scaled by $\Sigma$. For any $G$ with $\mathrm{rank}(G) < \Nimg$ (which holds whenever $\Nobs < \Nimg$), eigenvalues lie in $[0, 1)$ with spectral radius $\rho < 1$. Unobserved directions converge via the prior regularization in $\Sigma$. $T_0$ is therefore a contraction and $\mu_k \to \mu^\star$ at rate $\rho^k$. $\square$

\paragraph{Case (b): small $\Delta$.}
By matrix perturbation theory (Weyl's inequality on singular values), $\rho(DT_\Delta) = \rho(I - K_{\tilde{G}}\tilde{G})$ is continuous in $\Delta$. At $\Delta = 0$, $\rho = \rho(I - K_G G) < 1$. For small $\|\Delta\|$, $\rho_\Delta < 1$ and the fixed point is perturbed:
$$
\mu^\star_\Delta \;=\; \mu^\star \;+\; O\bigl(\|\Delta\|/(1-\rho)\bigr)
$$
by the implicit function theorem. However, in the non-Gaussian (diffusion) case, the score-matching M-step introduces finite-sample variance that compounds with the bias. After $k^\star \approx \log(1/\|\Delta\|)/\log(1/\rho)$ iterations, the accumulated perturbation can push the effective spectral radius past $1$. $\square$

\paragraph{Case (c): large $\Delta$.}
When $\|K_{\tilde{G}} \Delta\|$ is large enough that $\rho(I - K_{\tilde{G}} \tilde{G}) \geq 1$, the recurrence is not contractive. Depending on the eigenstructure, iterates either diverge or oscillate around a point far from $\mu^\star$. In the degenerate case where $\tilde{G}$ has a fundamentally different null space than $G$ (as with $\Aimg$ vs.\ $G = \Ats \circ \imgts$), certain directions receive no corrective signal, trapping $\mu_k$ at a biased fixed point. $\square$

\paragraph{Connection to the empirical regimes.}
The three trajectories of \cref{fig:convergence} and \cref{tab:convergence} map directly onto the theorem:
\begin{center}
\begin{tabular}{c|l|l|l}
\toprule
\textbf{Regime} & \textbf{Theorem prediction} & \textbf{Empirical} & \textbf{Behavior} \\
\midrule
(a) $\Delta = 0$        & contraction to $\mu^\star$     & Obs-space CG (ours)       & converges to $0.045$ at iter 3, stable \\
(b) small $\|\Delta\|$  & early converge, divergence at $k^\star$ & Baseline (no obs-CG)    & $0.049$ at iter 3, diverges to $0.456$ \\
(c) large $\|\Delta\|$  & biased fixed point              & Naive MMPS-EM             & stuck at $0.092$, never approaches $\mu^\star$ \\
\bottomrule
\end{tabular}
\end{center}

The mild degradation of the green curve ($0.045 \to 0.060$) beyond the peak is not predicted by the linear-Gaussian theorem (which gives monotone convergence for $\Delta = 0$). This residual is attributable to finite-sample effects in the score-matching M-step, not operator mismatch; importantly, it \emph{stabilizes} rather than diverging, consistent with regime (a).

\paragraph{Proof of Corollary~\ref{cor:corrections}.}
Each correction reduces a specific component of the effective mismatch $\|\Delta_{\mathrm{eff}}\|$:
\begin{itemize}
\setlength{\itemsep}{0pt}
\item \textbf{Obs-space CG} solves the linear system with $\tilde{G} = G$ directly, so $\Delta = 0$. Regime transition: (b)/(c) $\to$ (a).
\item \textbf{Adaptive $\sigma_y$} prevents CG divergence that amplifies $\Delta$ into large posterior errors (bounded $\kappa$). Stabilizes within current regime.
\item \textbf{Manifold projection} projects onto $\mathrm{Range}(\tsimg)$, removing off-manifold error where $\Delta$ has no self-correcting mechanism. Reduces effective $\|\Delta\|$.
\item \textbf{Domain-aware init} reduces $\|\mu_0 - \mu^\star\|$, delaying $k^\star$ in regime (b) and leaving the basin of regime (a) unchanged.
\end{itemize}
The ablation (\cref{tab:ablation}) confirms the predictions strikingly: partial corrections (Rows 2--3) reduce operator mismatch enough to destabilize the biased fixed point of regime (c), but without manifold projection the system has no mechanism to prevent off-manifold drift from compounding across iterations. The minimum viable configuration requires init $+$ $\sigma_y$ $+$ projection (Row 4, $0.049$); adding obs-space CG completes the transition to regime (a) (Row 5, $0.045$). $\square$

\section{Full Algorithms and Hyperparameter Audit}
\label{app:algorithms}

\subsection{The lift $\tsimg$: delay embedding and vanilla folding}
\label{app:lift}

The theorem of \cref{subsec:theorem} and the manifold-projection argument of App.~\ref{app:mmps-fails} both rest on the assumption that $\tsimg$ is a \emph{linear} operator with a closed-form (pseudo-)inverse $\imgts$. For completeness---and so that a non-time-series reviewer can follow the proof without chasing citations---this subsection fixes notation, instantiates $\tsimg$ in the two lifts we use (delay embedding and vanilla folding), and verifies linearity and the explicit form of $\imgts$ in each case.

\paragraph{Time series-to-image transformations.}
Mapping a time series to an image to leverage vision backbones has a long history~\citep{wang2015imaging, hatami2018classification, li2024time}, including short-time Fourier transform (STFT) for speech~\citep{griffin1984signal, allen1977short} and mel-spectrograms in diffusion-based audio models~\citep{popov2021grad, liu2023audioldm}. Our method is compatible with any such linear lift: the four E-step corrections (C1--C3) and \cref{thm:operator-mismatch} depend only on $\tsimg$ being linear, so the recipe transfers verbatim to STFT (App.~\ref{app:stft}) and, by the same route, to any of the listed alternatives.

\paragraph{Vanilla folding.}
Given a time series $\xts \in \mathbb{R}^{\Nts \times F}$, vanilla folding reshapes each channel into an $h \times w$ image by reading values left-to-right / top-to-bottom, zero-padding if $h\,w > \Nts$. The inverse reads the non-padded region row-wise. Folding is a special case of delay embedding (below) with skip equal to the row width, so linearity and the closed-form inverse carry over. Its virtue is simple scaling: a length-$65536$ sequence fits a $256\times 256$ image.

\paragraph{Delay embedding.}
Our headline lift follows~\citet{takens2006detecting} and the delay-embedding instantiation of~\citet{naiman2024imagentime}. For a univariate channel $x_{1:\Nts}$, two hyperparameters fix the geometry: the embedding dimension $n$ (number of rows) and the skip $m$ (column stride). Delay embedding writes
\[
X \;=\; \begin{bmatrix}
x_1 & x_{m+1} & \cdots & x_{\Nts - n} \\
\vdots & \vdots & \ddots & \vdots \\
x_n & x_{n+m+1} & \cdots & x_{\Nts}
\end{bmatrix}
\;\in\; \mathbb{R}^{n \times q}, \qquad q \;=\; \lceil (\Nts - n)/m \rceil,
\]
zero-padded to $n \times n$ so that the image is square. For a multivariate input with $F$ features, the per-channel matrices are stacked along a new axis to give $\ximg \in \mathbb{R}^{F \times n \times n}$. Choosing $m = n = 256$ encodes sequences of length up to $65\mathrm{k}$ in a $256 \times 256$ image---the scaling that motivates the method on long sequences. See~\citep{takens2006detecting, naiman2024imagentime} for the full spectral-geometric background.

\paragraph{Linearity and the pseudo-inverse.}
Both constructions are linear maps from $\mathbb{R}^{\Nts \times F}$ to $\mathbb{R}^{F \times n \times n}$: each image pixel is either zero (padding) or a single time-series entry, so the operator is written as a $\{0,1\}$-valued (per-channel) matrix. The pseudo-inverse $\imgts$ is therefore closed-form. When $m < n$ the lift is \emph{overcomplete}: a given time index $i$ appears in several image pixels. The original~\citep{takens2006detecting} reverse transformation keeps only the first occurrence, which discards redundant information. Following~\citet{fadlon2025imageni2r} we use the \emph{enhanced reverse transformation}: $\imgts(\ximg)[i]$ is the \emph{average} of all image pixels that delay-embed to index $i$. This change is relevant for three independent reasons inside our method:

\begin{enumerate}
\setlength{\itemsep}{0pt}
\item \textbf{Manifold projection.} The composition $\Pi = \tsimg \circ \imgts$ is the orthogonal projection onto $\mathrm{Range}(\tsimg)$ \emph{only} when $\imgts$ is the Moore--Penrose pseudo-inverse, i.e.\ averaging rather than first-pixel selection. Correction C3 (\cref{subsec:corrections}) relies on this identity.
\item \textbf{Adjoint identity.} In observation-space CG we use $(G)^\top v = \tsimg(\Ats^\top v)$. The adjoint of the averaging inverse is the broadcast that writes each time-series entry into all of its image pixels (divided by the multiplicity)---a sparse, per-pixel scale. The first-pixel inverse does not satisfy this identity, which is why we adopt the enhanced variant throughout.
\item \textbf{Operator-mismatch norm.} In App.~\ref{app:continuous-mismatch-norm} we report $\|\Delta_G\| = \|\tilde{G}^\top \tilde{G} - G^\top G\|$ under the enhanced inverse; the specific numerical values in \cref{tab:continuous-mismatch-norm} are tied to this choice of $\imgts$.
\end{enumerate}

With $\tsimg$ linear and $\imgts$ in closed form, the theorem's linearity assumption is satisfied for every lift used in the paper (delay embedding, vanilla folding, STFT), and the four corrections apply verbatim.

\begin{algorithm}[h]
\caption{$\lmmps$ E-step (single sample)}
\label{alg:estep}
\begin{algorithmic}[1]
\Require observation $\yts \in \mathbb{R}^{\Nobs \times F}$, mask $\Ats$
\Require denoiser $d_\theta$, EDM noise schedule $\{\sigma_t\}_{t = T, \ldots, 0}$
\Require initial image $x_{\mathrm{init}} \in \mathbb{R}^{\Nimg}$ (from \cref{alg:em}, line 4 or 22)
\Require regularization constant $c$ (default $0.1$), CG iters $n_{\mathrm{CG}}$ (default $5$), projection cadence $p_{\mathrm{proj}}$ (default: every reverse step)
\Ensure posterior sample $\ximg \in \mathrm{Range}(\tsimg)$ with $G\,\ximg \approx \yts$
\State $x_T \gets x_{\mathrm{init}} + \sigma_T \epsilon$, \quad $\epsilon \sim \mathcal{N}(0, I)$
\For{$t = T, T-1, \ldots, 1$}
  \State $\hat{x}_0 \gets d_\theta(x_t, \sigma_t)$ \Comment{Tweedie mean}
  \State $\sigma_y \gets c \cdot \sigma_t$ \Comment{adaptive $\sigma_y$ (\cref{subsec:corrections})}
  \State $r \gets \yts - \Ats \cdot \imgts(\hat{x}_0)$ \Comment{residual via true $G$}
  \State $v \gets \mathrm{CG}_{\mathrm{obs}}(r; \sigma_y, \sigma_t, d_\theta, x_t, n_{\mathrm{CG}})$ \Comment{runs in $\mathbb{R}^{\Nobs \times F}$}
  \State $\delta \gets \mathrm{VJP}(d_\theta, x_t, \tsimg(\Ats^\top v))$ \Comment{posterior score: $J^\top G^\top v$}
  \State $x_{t-1} \gets \mathrm{EDMstep}(x_t, \hat{x}_0 + \sigma_t^2 \delta, \sigma_t, \sigma_{t-1})$
  \If{$t \mod p_{\mathrm{proj}} = 0$} \Comment{manifold projection}
    \State $\tilde{x}_{\mathrm{ts}} \gets \imgts(x_{t-1})$
    \State $\tilde{x}_{\mathrm{ts}}[\mathrm{observed}] \gets \yts[\mathrm{observed}]$ \Comment{hard observation enforcement}
    \State $x_{t-1} \gets \tsimg(\tilde{x}_{\mathrm{ts}})$
  \EndIf
\EndFor
\State \Return $x_0$
\\
\Function{$\mathrm{CG}_{\mathrm{obs}}$}{$r; \sigma_y, \sigma_t, d_\theta, x_t, n_{\mathrm{CG}}$}
  \State Define $A_{\mathrm{mvp}}(v_{\mathrm{ts}})$:
  \State \quad $v_{\mathrm{img}} \gets \tsimg(\Ats^\top v_{\mathrm{ts}})$ \Comment{$G^\top v$}
  \State \quad $J v_{\mathrm{img}} \gets \mathrm{VJP}(d_\theta, x_t, v_{\mathrm{img}})$ \Comment{one VJP}
  \State \quad $J v_{\mathrm{ts}} \gets \Ats \cdot \imgts(J v_{\mathrm{img}})$ \Comment{$G(\cdot)$}
  \State \quad \Return $\sigma_y^2\, v_{\mathrm{ts}} + \sigma_t^2\, J v_{\mathrm{ts}}$
  \State \Return $\mathrm{CG}(A_{\mathrm{mvp}}, r, \text{max iters} = n_{\mathrm{CG}})$
\EndFunction
\end{algorithmic}
\end{algorithm}

\begin{algorithm}[h]
\caption{Co-Evolving EM}
\label{alg:em}
\begin{algorithmic}[1]
\Require corrupted dataset $\mathcal{D} = \{(\yts^{(i)}, \Ats^{(i)})\}_{i=1}^N$
\Require EM iters $K$ (default $5$; peak at $k = 3$)
\Require init method $\mathcal{I} \in \{\mathrm{STL}, \mathrm{Kalman}\}$ (default STL)
\Require curriculum schedule $\pi : \{1, \ldots, K\} \to [0, 1]$ (default linear $0.3 \to 0.0$ over $k = 1, \ldots, K$)
\Require E-step hyperparameters $(c, n_{\mathrm{CG}}, p_{\mathrm{proj}})$ (defaults in \cref{alg:estep})
\Require M-step: EDM score matching for $E$ epochs (ImagenTime defaults)
\Ensure trained denoiser $\theta_K$
\For{$i = 1, \ldots, N$} \Comment{warm start}
  \State $\tilde{x}_{\mathrm{ts}}^{(i)} \gets \mathcal{I}(\yts^{(i)}, \Ats^{(i)})$
  \State $x_{\mathrm{init}}^{(i)} \gets \tsimg(\tilde{x}_{\mathrm{ts}}^{(i)})$
\EndFor
\State Initialize $\theta_0$ \Comment{random or pretrained}
\For{$k = 1, \ldots, K$}
  \For{$i = 1, \ldots, N$} \Comment{curriculum: reveal extras from prior completion}
    \State $\tilde{A}^{(i)}_{\mathrm{ts}} \gets \Ats^{(i)} \cup \mathrm{Reveal}(\pi(k), x_{\mathrm{init}}^{(i)})$
    \State $\tilde{y}^{(i)}_{\mathrm{ts}} \gets \yts^{(i)}$ on $\Ats^{(i)}$; $\tilde{x}^{(i)}_{\mathrm{ts}}$ on extras
  \EndFor
  \For{$i = 1, \ldots, N$ \textbf{in parallel}} \Comment{E-step}
    \State $\ximg^{(i,k)} \gets \Call{Algorithm~\ref{alg:estep}}{\tilde{y}^{(i)}_{\mathrm{ts}}, \tilde{A}^{(i)}_{\mathrm{ts}}, d_{\theta_{k-1}}, x_{\mathrm{init}}^{(i)}}$
  \EndFor
  \State $\theta_k \gets \arg\min_\theta \mathbb{E}_{i, \sigma, \epsilon}[\lambda(\sigma) \|d_\theta(\ximg^{(i,k)} + \sigma\epsilon, \sigma) - \ximg^{(i,k)}\|^2]$ \Comment{M-step: vanilla EDM}
  \For{$i = 1, \ldots, N$} $x_{\mathrm{init}}^{(i)} \gets \ximg^{(i,k)}$ \EndFor
\EndFor
\State \Return $\theta_K$
\end{algorithmic}
\end{algorithm}

\paragraph{Hyperparameter audit.}
\Cref{tab:hparams} exposes every knob of the method. All entries are fixed defaults---either inherited from the base pipelines we extend or set by a single uniform convention---and none is dataset-specific; the same values are used across Stocks, Energy, MuJoCo, Sine, and Weather.

\begin{table}[h]
\centering
\caption{Complete hyperparameter audit. No tuned hyperparameters beyond EDM $+$ ImagenTime $+$ naive MMPS-EM defaults.}
\label{tab:hparams}
\resizebox{0.9\textwidth}{!}{
\begin{tabular}{l|c|l}
\toprule
\textbf{Knob} & \textbf{Value} & \textbf{Where fixed} \\
\midrule
$c$ in $\sigma_y = c \cdot \sigma_t$                & $0.1$                             & Fixed default (shared across all datasets) \\
CG iterations $n_{\mathrm{CG}}$                     & $5$                               & Default \\
Projection cadence $p_{\mathrm{proj}}$              & every reverse step                & Default \\
EM iterations $K$                                   & $5$ (peak at $k = 3$)             & Default \\
Init method $\mathcal{I}$                           & STL                               & App.~\ref{app:additional-exp} shows STL $\approx$ Kalman \\
Curriculum $\pi(k)$                                 & linear $0.3 \to 0.0$              & Default \\
Reverse diffusion steps, EDM sampler                & EDM defaults                      & Inherited~\citep{karras2022edm} \\
M-step optimizer, epochs-per-iter                   & ImagenTime defaults               & Inherited~\citep{naiman2024imagentime} \\
\bottomrule
\end{tabular}
}
\end{table}

\subsection{Formal E-step and M-step objectives}
\label{app:losses}

\cref{alg:em} specifies the MCEM loop algorithmically; for a reviewer assessing the statistical content, this subsection writes the two objectives in closed form and verifies that together they implement one MCEM step in the sense of~\citet{wei1990monte}.

\paragraph{E-step: posterior expectation.}
Let $\theta^{(k)}$ denote the model parameters after $k$ EM iterations and let $\yts = \Ats \cdot \xts$ be the observed (irregular) time series. The exact E-step surrogate is
\begin{equation}
Q\bigl(\theta \,\big|\, \theta^{(k)}\bigr) \;=\; \mathbb{E}_{\xts^{\mathrm{miss}} \,\sim\, p_{\theta^{(k)}}(\xts^{\mathrm{miss}} \,\mid\, \yts)} \bigl[\log p_\theta(\xts)\bigr].
\label{eq:Q-surrogate}
\end{equation}
We approximate the intractable posterior $p_{\theta^{(k)}}(\xts^{\mathrm{miss}} \mid \yts)$ by drawing a single sample via our E-step sampler (\cref{alg:estep}) in the lifted (image) space:
\begin{equation}
\ximg^{(i,k)} \;\sim\; p_{\theta^{(k)}}\bigl(\ximg \,\big|\, \yts^{(i)}\bigr) \quad \text{(obtained by \cref{alg:estep})},
\label{eq:Estep-sample}
\end{equation}
yielding the standard MCEM Monte-Carlo approximation $\hat{Q}(\theta | \theta^{(k)}) = \frac{1}{N} \sum_i \log p_\theta(\ximg^{(i,k)})$. The four E-step corrections (C1--C3 plus domain-aware warm start) are exactly the ingredients needed for \cref{eq:Estep-sample} to be a valid posterior draw in the dual-space setting; see \cref{subsec:corrections} and App.~\ref{app:mmps-fails} for why each correction is necessary.

\paragraph{M-step: EDM score matching on the completed image.}
The M-step maximizes $\hat{Q}$ in $\theta$. Because the prior is a diffusion model, $\log p_\theta$ is parameterized by the denoiser $d_\theta$ via the EDM formulation~\citep{karras2022edm}, and maximizing $\hat{Q}$ reduces to the weighted denoising score-matching objective restricted to the completed images:
\begin{equation}
\mathcal{L}_{\mathrm{M}}(\theta) \;=\; \mathbb{E}_{i, \sigma, \epsilon} \Bigl[\, \lambda(\sigma) \, \bigl\| d_\theta(\ximg^{(i,k)} + \sigma \epsilon,\, \sigma) - \ximg^{(i,k)} \bigr\|_2^2 \,\Bigr],
\label{eq:Mstep}
\end{equation}
with $\sigma \sim p(\sigma)$ and $\epsilon \sim \mathcal{N}(0, I)$ drawn from the EDM noise schedule and $\lambda(\sigma)$ the EDM weighting function. The crucial design choice is that \cref{eq:Mstep} is computed on the \emph{full} completed image $\ximg^{(i,k)}$---every pixel contributes to the score-matching loss---unlike open-loop commit-and-forget pipelines~\citep{fadlon2025imageni2r} that mask the loss to observed pixels only. The mask-the-loss trick was necessary for commit-and-forget because its completions had no noise model to blame the unobserved residuals on; in our setting, $\ximg^{(i,k)}$ is a \emph{posterior sample}, so every pixel carries calibrated uncertainty and belongs in the score-matching loss. This is the concrete technical cash-out of the masking--data trade-off (\cref{sec:why-close}).

\paragraph{One MCEM step.}
Together, \cref{eq:Estep-sample} (E-step draw) and \cref{eq:Mstep} (M-step minimization) implement one iteration of MCEM with a single posterior draw per sample. This is the minimal-variance instantiation; drawing $S > 1$ samples per observation and averaging \cref{eq:Mstep} over them would reduce Monte-Carlo noise at the cost of $S\times$ E-step compute, and is the natural route to tighten the variance of $\hat{Q}$ at fixed $K$. We did not find this necessary on our datasets: $S = 1$ with $K = 5$ already tracks the regime-(a) contraction curve of \cref{thm:operator-mismatch} (\cref{tab:convergence}).

\section{Why Close the Loop --- Full Theory}
\label{app:why-close-full}

\Cref{sec:why-close} sketches the four mechanisms that make iterative co-evolving EM outperform commit-and-forget. Here we provide the full derivation of each.

\subsection{The conditional-mean trap}
\label{app:cond-mean-trap}

The core mathematical reason commit-and-forget fails is that \emph{deterministic imputers output point estimates, but generative models need distributional samples.} Given a conditional distribution $p(\xts^{\mathrm{miss}} \mid \xts^{\mathrm{obs}})$, the MSE-optimal estimate is the conditional mean $\hat{x} = \mathbb{E}[\xts^{\mathrm{miss}} \mid \xts^{\mathrm{obs}}]$, and the MAE-optimal estimate is the conditional median. Both the TST (trained with MSE) and NCDE (cubic spline interpolation) approximate the conditional mean. This is MSE-optimal per-sample but distributionally catastrophic when the conditional is multimodal.

\paragraph{Example.}
At a point where two clusters overlap, the true conditional is bimodal:
$$
p(x_2 \mid x_1 \approx 0) = 0.5 \cdot \mathcal{N}(+1.5, 0.1) + 0.5 \cdot \mathcal{N}(-0.5, 0.1).
$$
The MSE-optimal estimate is $\hat{x}_2 = 0.5$---\emph{between} the clusters, in empty space. A generative model trained on mean-imputed data learns to generate at $x_2 = 0.5$, which is wrong. Posterior samples land on $+1.5$ or $-0.5$, preserving both modes. \Cref{tab:toy-mse-trap} visualizes the resulting inversion: deterministic imputers look best on MSE/MAE but worst downstream, while co-evolving EM reverses that ordering.

\subsection{Variance collapse}
\label{app:var-collapse}

By the law of total variance,
$$
\mathrm{Var}(X) = \mathrm{Var}\bigl(\mathbb{E}[X \mid \mathrm{obs}]\bigr) + \mathbb{E}\bigl[\mathrm{Var}(X \mid \mathrm{obs})\bigr].
$$
If every imputed value is replaced by $\mathbb{E}[X \mid \mathrm{obs}]$, the imputed dataset has variance $\mathrm{Var}(\mathbb{E}[X \mid \mathrm{obs}]) < \mathrm{Var}(X)$. The deficit is $\mathbb{E}[\mathrm{Var}(X \mid \mathrm{obs})]$, which at high missing rates is enormous. This variance deficit is well established in the multiple imputation literature~\citep{rubin1987multiple}; our contribution is to show it directly impacts generative model quality in the irregular time-series setting.

In EM with a stochastic E-step (Monte Carlo EM), the imputed values are \emph{samples} from $p(\xts^{\mathrm{miss}} \mid \xts^{\mathrm{obs}}, \theta)$. Each individual sample has higher MSE than the conditional mean, but the collection of samples preserves the correct marginal variance and multimodality. The generative model sees all modes represented in its training data.

\subsection{The masking--data trade-off}
\label{app:masking-tradeoff}

ImagenI2R partially mitigates variance collapse by using a masked loss---the diffusion model trains only on observed positions, ignoring imputed values. This avoids learning the collapsed variance, but at a cost:
\begin{itemize}
\setlength{\itemsep}{0pt}
\item At missing rate $r$, only $(1-r)$ fraction of pixels contribute to the loss.
\item At $r = 0.7$, only $30\%$ of training signal is used.
\item Noise is added only at observed positions during training but everywhere during sampling, creating a train/test mismatch in the forward process.
\end{itemize}
EM avoids this entirely: after the E-step, the data is complete. The M-step trains with full loss on every pixel, with no masking and no mismatch. Our experiments show that EM with the correct E-step at $50\%$ missing achieves $0.045$, matching the clean-data oracle ($0.044$) within statistical noise---a result impossible under the masking--data trade-off.

\subsection{Why stochastic commit-and-forget is not enough}
\label{app:frozen-stochastic}

Stochastic imputation (e.g.\ CSDI~\citep{tashiro2021csdi}) produces diverse posterior samples, addressing the conditional-mean trap. However, CSDI-as-imputer still follows commit-and-forget:
\begin{enumerate}
\setlength{\itemsep}{0pt}
\item \emph{Frozen imputation prior.} CSDI trains once on the irregular data, then generates completions; its prior is never updated based on whether the completions produce good generations.
\item \emph{No end-to-end alignment.} CSDI optimizes imputation quality (CRPS/MSE on held-out positions), not generation quality (discriminative score). These objectives are misaligned---good imputation MSE does not imply good generation (App.~\ref{app:toy}, Finding~1).
\item \emph{Fragment-based learning at high missing rates.} At $50\%$ missing, CSDI sees different random subsets per sample. It must reconstruct the joint distribution from scattered fragments without the iterative refinement that EM provides.
\item \emph{Sampling cost.} CSDI is a conditional diffusion model that uses up to $1000$ reverse-diffusion steps per imputation~\citep{tashiro2021csdi}; this is already slow on the Energy benchmark at $\Nts{=}24$ and reported infeasible at $\Nts{=}768$~\citep{fadlon2025imageni2r}. This is a practical barrier distinct from the statistical one above, but relevant whenever one considers dropping CSDI into a larger pipeline (App.~\ref{app:inference-time}).
\end{enumerate}

\paragraph{Formal argument: frozen imputation is a biased ELBO.}
We give the math behind the main-text claim that \emph{stochasticity is orthogonal to the open-loop bias; only iteration closes it}. Let $x = (\xts^{\mathrm{obs}}, \xts^{\mathrm{miss}})$ be the complete data, with the observation operator fixing $\xts^{\mathrm{obs}}$. The generator $p_\theta(x)$ is trained to maximize $\log p_\theta(\xts^{\mathrm{obs}})$. For any distribution $q(\xts^{\mathrm{miss}}\mid\xts^{\mathrm{obs}})$, the free-energy decomposition~\citep{neal1998em} gives
\begin{equation}
\label{eq:free-energy-app}
\log p_\theta(\xts^{\mathrm{obs}}) \;=\; \mathcal{F}(q, \theta) \;+\; \underbrace{\mathrm{KL}\!\left(q(\xts^{\mathrm{miss}}\mid\xts^{\mathrm{obs}})\,\big\|\,p_\theta(\xts^{\mathrm{miss}}\mid\xts^{\mathrm{obs}})\right)}_{\Delta(q,\theta)\,\geq\,0},
\end{equation}
with $\mathcal{F}(q,\theta) = \mathbb{E}_q[\log p_\theta(\xts^{\mathrm{obs}},\xts^{\mathrm{miss}})] + H(q)$ and $\Delta(q,\theta) = 0$ iff $q = p_\theta(\cdot\mid\xts^{\mathrm{obs}})$. Define the commit-and-forget objective as $\theta$ is fit with a \emph{fixed} $q_0$ returned by an imputer trained once on the data:
\begin{equation}
\theta^\star_{q_0} \;:=\; \arg\max_\theta \; \mathcal{F}(q_0, \theta) \;=\; \arg\max_\theta \;\mathbb{E}_{q_0}[\log p_\theta(\xts^{\mathrm{obs}},\xts^{\mathrm{miss}})].
\end{equation}

\begin{proposition}[Irreducible bias of frozen imputation]
\label{prop:frozen-bias}
For any fixed $q_0$ and any $\theta$,
\begin{equation}
\label{eq:frozen-identity}
\log p_\theta(\xts^{\mathrm{obs}}) \;=\; \mathcal{F}(q_0, \theta) \;+\; \Delta(q_0, \theta),
\end{equation}
so the gradients satisfy
\begin{equation}
\label{eq:grad-identity}
\nabla_\theta \log p_\theta(\xts^{\mathrm{obs}}) \;=\; \nabla_\theta \mathcal{F}(q_0, \theta) \;+\; \nabla_\theta \Delta(q_0, \theta).
\end{equation}
In particular, at the frozen-imputer fixed point $\theta^\star_{q_0}$ (where $\nabla_\theta\mathcal{F}(q_0,\theta^\star_{q_0}) = 0$) we have
\begin{equation}
\label{eq:bias}
\nabla_\theta \log p_\theta(\xts^{\mathrm{obs}}) \big|_{\theta^\star_{q_0}} \;=\; \nabla_\theta \Delta(q_0, \theta) \big|_{\theta^\star_{q_0}},
\end{equation}
which vanishes iff $q_0 = p_{\theta^\star_{q_0}}(\cdot\mid\xts^{\mathrm{obs}})$, the EM self-consistency condition. Equivalently, the log-likelihood sub-optimality of $\theta^\star_{q_0}$ relative to the MLE $\theta^\star$ is bounded below by a distributional (not moment) quantity:
\begin{equation}
\log p_{\theta^\star}(\xts^{\mathrm{obs}}) \;-\; \log p_{\theta^\star_{q_0}}(\xts^{\mathrm{obs}}) \;\geq\; \Delta(q_0, \theta^\star_{q_0}) \;-\; \Delta(q_0, \theta^\star).
\end{equation}
\end{proposition}

\begin{proof}
Equation~\eqref{eq:frozen-identity} is \eqref{eq:free-energy-app} evaluated at $q = q_0$; differentiating in $\theta$ gives \eqref{eq:grad-identity}. At $\theta^\star_{q_0}$, the first-order condition for $\mathcal{F}(q_0,\cdot)$ forces $\nabla_\theta\mathcal{F}(q_0,\theta^\star_{q_0}) = 0$, which combined with \eqref{eq:grad-identity} yields \eqref{eq:bias}. The sub-optimality bound follows by subtracting \eqref{eq:frozen-identity} at $\theta^\star$ and $\theta^\star_{q_0}$ and using $\mathcal{F}(q_0,\theta^\star_{q_0}) \geq \mathcal{F}(q_0,\theta^\star)$ by definition of $\theta^\star_{q_0}$.
\end{proof}

\paragraph{Iteration, and only iteration, closes $\Delta$.}
The reason EM works is that the E-step resets $\Delta$ to zero every iteration: setting $q_{k+1} \leftarrow p_{\theta_k}(\cdot\mid\xts^{\mathrm{obs}})$ gives $\Delta(q_{k+1}, \theta_k) = 0$, so locally around $\theta_k$, maximizing $\mathcal{F}(q_{k+1}, \cdot)$ and $\log p_\theta(\xts^{\mathrm{obs}})$ agree up to second order. The Neal--Hinton monotonicity theorem then guarantees $\log p_{\theta_{k+1}}(\xts^{\mathrm{obs}}) \geq \log p_{\theta_k}(\xts^{\mathrm{obs}})$, with strict inequality unless the pair $(q_k, \theta_k)$ is already an EM fixed point~\citep{neal1998em,dempster1977em}. Our co-evolving EM is exactly this scheme, with the diffusion posterior sampler (\cref{subsec:corrections}) instantiating the E-step in a lifted representation.

\paragraph{Stochasticity is orthogonal to the bias.}
\Cref{prop:frozen-bias} shows the frozen-imputer bias is governed by $\Delta(q_0,\theta) = \mathrm{KL}(q_0 \,\|\, p_\theta(\cdot\mid\xts^{\mathrm{obs}}))$, a \emph{distributional} object. Two extreme choices of $q_0$ illustrate that the bias has two separate components:
\begin{itemize}
\setlength{\itemsep}{0pt}
\item \emph{Deterministic $q_0 = \delta_{\mu_0}$} (mean imputation, TST, NCDE, GRU-D): $\Delta = \infty$ whenever $p_\theta(\cdot\mid\xts^{\mathrm{obs}})$ is not itself a point mass, which for realistic multimodal data it never is. This is the conditional-mean trap from the distributional side: $\Delta$ blows up because $q_0$ has zero support mass where $p_\theta(\cdot\mid\xts^{\mathrm{obs}})$ is nontrivial.
\item \emph{Stochastic $q_0$ with the right moments} (CSDI, any dedicated distributional imputer): $\Delta$ is finite and can be made small in first/second moments, but is still generally nonzero because $q_0$ is built without reference to the generator $p_\theta$; no amount of training on the imputation objective can align $q_0$ with a posterior induced by a generator that has not yet been trained.
\end{itemize}
Stochasticity removes the \emph{zeroth-order} divergence (moving from $\Delta = \infty$ to $\Delta < \infty$) but leaves the full distributional gap. This is precisely why our stochastic steelman (CSDI-impute $+$ EDM, Row 0c) beats the deterministic commit-and-forget row ($0.105$ vs.\ the GT-GAN/KoVAE/ImagenI2R cluster) yet still loses to iterative naive MMPS-EM ($0.092$, Row 1a): moving $q_0$ from Dirac to diffuse closes one term, but only iteration closes the other.

\paragraph{Intuition: whose posterior is the imputer supposed to sample?}
A further way to see the bias: an imputer is asked to sample $p(\xts^{\mathrm{miss}}\mid\xts^{\mathrm{obs}})$, but that conditional is defined \emph{only} relative to some joint distribution. CSDI's conditional is the one implied by its own score-matching objective on the irregular data; that joint is, in general, \emph{not} the joint $p_\theta(x)$ that the downstream generator learns. The commit-and-forget assumption is that these two joints are close; but they need not be---especially when the generator has higher capacity, uses a different backbone (EDM on an image lift versus CSDI's transformer in time), or is trained on a different loss (score matching on complete data versus score matching on held-out imputation targets). Iteration replaces this assumption with a guarantee: the E-step imputes under \emph{the generator's own} joint, so by construction the two agree.

\paragraph{Connection to the toy experiment.}
\Cref{prop:frozen-bias} predicts three empirical regularities that App.~\ref{app:toy} verifies:
\begin{enumerate}
\setlength{\itemsep}{0pt}
\item The frozen-stochastic baseline should plateau at a nonzero discrepancy above the oracle, set by $\Delta(q_0, \theta^\star_{q_0})$. On the $21$-configuration toy, mature frozen-stochastic F stalls at median $0.151$ versus the oracle at $0.026$---a $\sim 4\times$ gap that pooling more samples (F$'$) does not close (median $0.154$).
\item Pooling samples from a frozen $q_0$ cannot close $\Delta$, because pooling changes $q_0$'s Monte-Carlo variance, not its distributional distance to $p_{\theta^\star_{q_0}}(\cdot\mid\xts^{\mathrm{obs}})$. We observe exactly this (F vs.\ F$'$).
\item A single EM iteration is insufficient because $\theta_0$ is itself seeded from $q_0$: Row $5$b (Ours, $1$ EM iter) lands at $0.195$ on the real ablation, worse than the $200$-epoch frozen stochastic Row $0$c ($0.105$), as predicted by \cref{prop:frozen-bias} with $k = 1$.
\end{enumerate}
All three are quantitative consequences of Equation~\eqref{eq:bias}, not separate phenomena.

Our toy experiment (App.~\ref{app:toy}) sharpens this argument. First, even a mature frozen stochastic imputer hits a capacity ceiling that pooling more posterior samples cannot break---it is $\sim 4\times$ worse than co-evolving EM (D) despite beating every deterministic baseline. Second, an \emph{underfit} frozen stochastic imputer (F$_{\mathrm{naive}}$, analog of ``Ours 1 EM iter only'') is even worse, exactly mirroring the real-setting fact that Row $5$b ($0.195$) trails Row $0$c ($0.105$). Only iterative refinement (D, $0.034$) reaches the oracle ($0.033$). At scale, iterative EM with a bad E-step already outperforms even our dedicated $200$-epoch stochastic steelman (CSDI-impute $+$ EDM)---the value of EM is the feedback loop itself, not the quality of any single imputation step.

\subsection{Architectural open-loop (GRU-$\Delta t$ / GRU-D family)}
\label{app:architectural-openloop}

The commit-and-forget analysis of App.~\ref{app:cond-mean-trap}--\ref{app:frozen-stochastic} targets methods with an \emph{explicit} imputer. A separate sub-family, introduced as the baseline suite of~\citet{jeon2022gtgan}, handles irregularity \emph{architecturally} inside an RNN. We cover it here for completeness.

\paragraph{The two primitives.}
\emph{GRU-$\Delta t$}~\citep{che2018recurrent} takes the time gap between consecutive observations, $\Delta t_i = t_i - t_{i-1}$, as an additional input feature concatenated to the observation at step $i$. The RNN thus ``sees'' irregularity as a scalar per step but does not fill in between observations. \emph{GRU-D}~\citep{che2018recurrent} augments GRU-$\Delta t$ with a learnable exponential decay of the hidden state toward a data-dependent baseline between observations: $h_i \leftarrow \gamma_i \odot h_{i-1} + (1 - \gamma_i) \odot \bar{h}$, with $\gamma_i = \exp(-\max(0, W_\gamma \Delta t_i + b_\gamma))$. Between observations, the hidden state smoothly interpolates toward $\bar{h}$---an implicit, deterministic reconstruction.

\paragraph{Applying each failure mode.}
The four open-loop failure modes of \cref{sec:why-close} apply as follows.
\begin{itemize}
\setlength{\itemsep}{0pt}
\item \emph{(i) Conditional-mean trap / (ii) Variance collapse.} GRU-D's exponential-decay mechanism is a deterministic single-point smoothing between observations. It has the same structural defect as a deterministic imputer: at any point where the conditional is multimodal, the decay target $\bar{h}$ lies between modes. GRU-$\Delta t$ does not interpolate at all---the network merely sees $\Delta t$ as a feature---but this only shifts the problem: the generator is still trained on a single point-valued trajectory per sample. Neither primitive produces multi-modal completions.
\item \emph{(iii) Masking--data trade-off.} Not applicable: these methods do not use a masked loss.
\item \emph{(iv) Frozen-prior capacity ceiling.} The $\Delta t$/decay primitive is a \emph{fixed} architectural choice. Although it does receive gradient from the generator's training loss (so it is not frozen in the parameter sense), it has no posterior sampler, no multi-modal reconstruction, and no iterative refinement mechanism: the generator's gradient alone cannot transform a single-point recurrence into a distribution over plausible completions. From the perspective of the generator's distributional objective, it is structurally identical to the commit-and-forget sub-family.
\end{itemize}

\paragraph{Empirical evidence.}
\citet{jeon2022gtgan} (Tables 2--4 therein) reports that at $30/50/70\%$ MCAR missingness, the GRU-$\Delta t$/-D adaptations of RCGAN, C-RNN-GAN, T-Forcing, and P-Forcing yield discriminative scores at or near $0.5$ (random chance) on the majority of (dataset, rate) cells; the sole survivor carried forward by ImagenI2R~\citep{fadlon2025imageni2r} as the representative of this family is TimeGAN-$\Delta t$. Our own measurement of TimeGAN-$\Delta t$ at length $24$ (\cref{tab:results_short_l24} line for ``TimeGAN-$\Delta t$'') reproduces this saturation: disc $\in \{0.474, 0.479, 0.497, 0.499, 0.499, 0.499, 0.499, 0.499, 0.499\}$ across the $9$ datasets. This is the architectural analog of the frozen-prior capacity ceiling: a fixed $\Delta t$ primitive, even with end-to-end gradients, cannot close the loop.

\paragraph{Full open-loop taxonomy.}
\cref{tab:intro-taxonomy} consolidates both open-loop sub-families---the imputer-based line (GT-GAN, KoVAE, ImagenI2R) analysed in the main text and the architectural $\Delta t$/decay family analysed above---alongside our closed-loop EM method.

\begin{table}[t]
\centering
\caption{Existing paradigms for irregular time-series generation: all published methods follow an open-loop design (imputer-based or architectural). Our work closes the loop via co-evolving EM.}
\label{tab:intro-taxonomy}
\resizebox{0.95\textwidth}{!}{
\begin{tabular}{l l l l l}
\toprule
\textbf{Family} & \textbf{Representative} & \textbf{Missingness handling} & \textbf{Generator} & \textbf{Limitation} \\
\midrule
GAN-based (imputer) & GT-GAN~\citep{jeon2022gtgan} & Neural CDE (one-shot) & GAN & Unstable, mode collapse \\
VAE-based (imputer) & KoVAE~\citep{naiman2024kovae} & Neural CDE (one-shot) & Koopman VAE & $\sim 6.5\times$ slower, limited expressiveness \\
Diffusion-based (imputer) & ImagenI2R~\citep{fadlon2025imageni2r} & TST (one-shot, det.) & Masked diffusion & Masking--data trade-off \\
RNN (end-to-end) & TimeGAN-$\Delta t$/-D family~\citep{jeon2022gtgan} & Architectural: $\Delta t$ concat / GRU-D decay & GAN or forcing & No posterior sampling; deterministic decay smoothing \\
\midrule
\textbf{Iterative (ours)} & \textbf{Co-Evolving EM} & \textbf{Posterior sampling (E-step)} & \textbf{EDM / ImagenTime} & \textbf{---} \\
\bottomrule
\end{tabular}
}
\end{table}

\section{Controlled 2D Toy Experiment}
\label{app:toy}

This appendix archives the full toy experiment summarized in \cref{sec:why-close}. The experiment is explicitly designed as a \emph{$1$-to-$1$ analog} of the real-setting ablation (\cref{tab:ablation}): every baseline in the toy maps onto a specific row of the main ablation, so every toy finding translates into a prediction about the real system.

\subsection{Setup}
\label{app:toy-setup}

\paragraph{True distribution.}
Two interleaved spirals of Gaussian noise $\sigma = 0.12$; a thin curved manifold that a low-capacity imputer cannot represent.

\paragraph{Corruption.}
Each coordinate independently masked with probability $p_{\mathrm{miss}} = 0.60$ (so ${\sim}16\%$ of rows are fully observed and ${\sim}48\%$ have exactly one coordinate missing).

\paragraph{Model class.}
Gaussian Mixture Models. $N_{\mathrm{train}} = 2000$; evaluation against $N_{\mathrm{eval}} = 1500$ held-out clean samples; $N_{\mathrm{reps}} = 3$ per config (fresh mask and init each rep); seed $= 42$.

\paragraph{Iterations.}
$N_{\mathrm{iters}} = 12$ for co-evolving EM (D) and mature frozen stochastic (F); $N_{\mathrm{iters}} = 1$ for underfit frozen stochastic (F$_{\mathrm{naive}}$).

\paragraph{Metric.}
Paper's discriminative score $\mathrm{disc} = |0.5 - \mathrm{accuracy}_{\mathrm{GBM}}|$ (lower is better, $0 =$ classifier fooled, $0.5 =$ perfectly discriminated).

\paragraph{Capacity mismatch.}
We impose a capacity mismatch that faithfully mirrors the real-setting gap between a dedicated stochastic imputer and the downstream generator (CSDI $\neq$ EDM): the frozen-stochastic imputer (F) is low-capacity ($K_{\mathrm{imp}} = 3$), while the downstream generator and the co-evolving EM model are high-capacity ($K_{\mathrm{gen}} = 30$).

\subsection{Methods and real-setting correspondence}
\label{app:toy-methods}

\begin{table}[h]
\centering
\caption{Method $\leftrightarrow$ real-setting correspondence, and discriminative score on two spirals @ $p = 0.60$ (mean over $3$ reps).}
\label{tab:toy-methods}
\resizebox{\textwidth}{!}{
\begin{tabular}{l|l|l|c}
\toprule
\textbf{Method} & \textbf{Paradigm} & \textbf{Real-setting analog} & \textbf{disc} ($\downarrow$) \\
\midrule
(A) Oracle (clean data, $K = 30$)                        & ---                                                       & Row $\star$ ($0.044$)                     & $\mathbf{0.033}$ \\
(B) Regression $+$ masked loss (complete-case)           & commit-and-forget, deterministic, masked                  & ImagenI2R analog                         & $0.061$ \\
(B$'$) Regression $+$ full loss                          & commit-and-forget, deterministic, full loss               & canonical variance-collapse              & $0.388$ \\
(C) NCDE (KNN) $+$ full loss                             & commit-and-forget, deterministic, non-linear              & Row $0$a (GT-GAN / KoVAE analog)         & $0.263$ \\
(F$_{\mathrm{naive}}$) Frozen stochastic, $1$-iter       & commit-and-forget, stochastic, \emph{underfit}            & Row $5$b ``Ours ($1$ EM iter)'' ($0.195$) & $0.216$ \\
(F) Frozen stochastic, mature ($K_{\mathrm{imp}} = 3$)   & commit-and-forget, stochastic, mature                     & Row $0$c CSDI-impute $+$ EDM ($0.105$)    & $0.149$ \\
(F$'$) Same mature imputer, $K = 5$ pooled samples       & commit-and-forget, stochastic, pooled                     & (multiple-imputation analog)             & $0.154$ \\
(D) Co-Evolving EM ($K = 30$, $12$ iters)                & co-evolving (stochastic + iterative)                      & Row $5$ ``Ours ($5$ iters)'' ($0.045$)    & $\mathbf{0.034}$ \\
\bottomrule
\end{tabular}
}
\end{table}

\subsection{Five findings}
\label{app:toy-findings}

\paragraph{Finding 1: MSE/MAE is the wrong imputer metric for generation; the Energy Score is the right one.}
We report per-imputer reconstruction error on the missing entries alongside the Energy Score $\mathrm{ES}(P,x^\star)=\mathbb{E}\|X-x^\star\|-\tfrac{1}{2}\mathbb{E}\|X-X'\|$ (a strictly proper scoring rule on the joint conditional; uniquely minimized by the true posterior~\citep{gneiting2007strictly}) and the discriminative score in \cref{tab:toy-mse-trap}. Two observations:
\begin{enumerate}[leftmargin=*,topsep=0pt,itemsep=2pt]
  \item \textbf{MSE/MAE is structurally inverted across the deterministic/stochastic boundary.} The MSE-optimal imputer (B$'$) produces the \emph{worst} generation, and the highest-MSE imputer (F) produces the second-best. Co-evolving EM (D) has ${\approx}1.9\times$ higher MSE than regression (B$'$) yet matches the oracle in disc. The inversion is not coincidental: MSE is optimized by the conditional mean $\mathbb{E}[\xts^{\mathrm{miss}}\mid\xts^{\mathrm{obs}}]$ and MAE by the conditional median---both point functionals, both erasing the distributional shape the downstream generator needs. By the law of total variance, any distributional imputer must pay $\mathbb{E}[\mathrm{Var}(\xts^{\mathrm{miss}}\mid\xts^{\mathrm{obs}})]$ in MSE relative to a point estimator, so it \emph{cannot} win on MSE by construction.
  \item \textbf{ES agrees with disc in $21/21$ configurations} of the cherrypick sweep (App.~\ref{app:toy-robustness}, $M=50$ posterior samples, $N=5$ repetitions). On the two-spirals row in \cref{tab:toy-mse-trap}, ES ranks (B$'$, C) worst (tied with MAE as the ES reduces to MAE for deterministic imputers), (F$_{\mathrm{naive}}$, F) next, and (D) best---the same ordering as disc. MSE inverts this ordering in $21/21$ configurations.
\end{enumerate}
\Cref{tab:toy-mse-trap} shows all three metrics side by side. \emph{Methodological implication}: imputation benchmarks that rank methods by MSE/MAE---the standard way TST, NCDE, and CSDI are evaluated---systematically prefer imputers that \emph{hurt} downstream generation; benchmarks that use a strictly proper distributional score (ES, CRPS, kernel score) rank correctly.

\begin{table}[h]
\centering
\caption{Point (MSE/MAE), distributional (ES), and downstream (disc) metrics on the two-spirals toy at $p = 0.60$ ($M=50$ posterior samples, $5$ reps). MSE/MAE invert the disc ordering (MSE-optimal B$'$ is disc-worst, highest-MSE F is second-best disc); the Energy Score~\citep{gneiting2007strictly} ranks identically to disc.}
\label{tab:toy-mse-trap}
\resizebox{0.72\textwidth}{!}{
\begin{tabular}{l|cccc}
\toprule
\textbf{Method} & \textbf{imp.\ MSE} $\downarrow$ & \textbf{imp.\ MAE} $\downarrow$ & \textbf{imp.\ ES} $\downarrow$ & \textbf{disc} $\downarrow$ \\
\midrule
(B$'$) Regression $+$ full loss            & $\mathbf{2.22}$ (LOWEST) & $\mathbf{1.31}$ (LOWEST) & $1.31$ (WORST)           & $0.390$ (WORST) \\
(C) NCDE / KNN                              & $2.46$                  & $1.28$                  & $1.28$                  & $0.252$ \\
(F$_{\mathrm{naive}}$) Underfit stochastic  & $3.65$                  & $1.56$                  & $0.91$                  & $0.222$ \\
(F) Mature frozen stochastic                & $4.31$ (HIGHEST)        & $1.62$                  & $0.85$                  & $0.154$ \\
(D) Co-Evolving EM                          & $4.26$                  & $1.47$                  & $\mathbf{0.79}$ (BEST)  & $\mathbf{0.049}$ (BEST) \\
\midrule
Oracle (clean data)                         & ---                     & ---                     & $0.72$                  & $0.029$ \\
\bottomrule
\end{tabular}
}
\end{table}

\paragraph{Finding 2: Deterministic ``full loss'' is catastrophic.}
(B$'$, $0.388$) and (C, $0.263$) both train a high-capacity generator on variance-collapsed imputations; the generator faithfully reproduces the straight-line (B$'$) or between-cluster (C) imputation artifacts it was trained on. This is the failure mode App.~\ref{app:cond-mean-trap} warns about, appearing here at its strongest because the $2$D spirals' conditional is multi-modal almost everywhere along the manifold.

\paragraph{Finding 3: Masked loss (B) $\neq$ full loss (B$'$).}
Complete-case fitting (B, $0.061$) avoids variance-collapse because the generator never sees imputed values. In $2$D at $p = 0.60$, ${\sim}16\%$ of rows are still fully observed---enough for an almost-unbiased GMM estimator. The $6\times$ gap between (B$'$) and (B) isolates the cost of \emph{trusting} the deterministic imputations versus ignoring them. In the real setting at $50\%$--$70\%$ missing on $24$ features, the fully-observed fraction collapses to $\sim 10^{-4}$ and the per-pixel masked-loss variant takes over---paying the masking--data trade-off (App.~\ref{app:masking-tradeoff}) and degrading sharply.

\paragraph{Finding 4: A mature frozen stochastic prior beats the full-loss deterministic baselines, but still loses to iteration.}
(F, $0.149$) cleanly beats (B$'$, $0.388$) and (C, $0.263$), mirroring CSDI-impute $+$ EDM beating GT-GAN/KoVAE/ImagenI2R at scale. Pooling $K = 5$ posterior samples (F$'$, $0.154$) produces no meaningful improvement over a single sample (F, $0.149$)---\emph{pooling is not a substitute for iteration}, because the bottleneck is the imputer's fixed, low capacity. No amount of re-sampling from a model that cannot represent the spirals produces spiral-shaped training data.

\paragraph{Finding 5: Underfit stochastic $<$ mature stochastic $<$ iterative.}
(F$_{\mathrm{naive}}$, $0.216$) is worse than the mature (F, $0.149$), and both are ${\sim}4\times$ worse than co-evolving EM (D, $0.034$). This is the direct $2$D analog of Row $5$b ($0.195$) $>$ Row $0$c ($0.105$) $>$ Row $5$ ($0.045$). Co-evolving EM reaches disc $0.034$, within $0.001$ of the oracle ($0.033$)---a near-oracle gap consistent with what we observe at scale on the real benchmark.

\subsection{Robustness: 21 configurations}
\label{app:toy-robustness}

To guard against cherry-picking, we ran the identical protocol across $7$ distributions (moons+ellipse, pinwheel-$5$, $9$-component grid, rings, checkerboard, two-spirals, swirl) $\times$ $3$ missing rates ($0.50, 0.60, 0.70$) $=$ $21$ configurations with fresh data and seeds.

\begin{table}[h]
\centering
\caption{Robustness across $21$ configurations. Every pairwise claim in Findings $2$--$5$ is universal.}
\label{tab:toy-robustness}
\resizebox{0.8\textwidth}{!}{
\begin{tabular}{l|c}
\toprule
\textbf{Claim} & \textbf{Holds in} \\
\midrule
$D < F$ (iterative beats mature frozen stochastic)                        & $\mathbf{21\,/\,21}$ \\
$F < F_{\mathrm{naive}}$ (mature beats underfit, $\geq 0.01$ margin)      & $\mathbf{20\,/\,21}$ (grid-$9$@$0.50$ is a $0.002$ within-noise reversal) \\
$D < C$ (iterative beats NCDE-analog)                                     & $\mathbf{21\,/\,21}$ \\
$D < B'$ (iterative beats full-loss regression)                           & $\mathbf{21\,/\,21}$ \\
$D$ within $0.07$ disc of Oracle                                          & $\mathbf{21\,/\,21}$ \\
\bottomrule
\end{tabular}
}
\end{table}

Aggregated across all $21$ configurations, the median disc scores are $D = 0.056$, $B = 0.082$, $F = 0.151$, $F' = 0.154$, $F_{\mathrm{naive}} = 0.216$, $C = 0.228$, $B' = 0.373$, Oracle $= 0.026$. The $F - D$ gap has minimum $+0.054$, median $+0.101$, maximum $+0.149$---iteration wins by at least $0.054$ disc units everywhere.

\subsection{Scope of the toy---what it does and does not prove}

\begin{itemize}
\setlength{\itemsep}{0pt}
\item \emph{What it proves.} (i) The three mechanisms App.~\ref{app:cond-mean-trap}--\ref{app:frozen-stochastic} predict are real and quantitatively drive the disc ordering; (ii) the ordering holds independently of architecture, dimensionality, and optimizer; (iii) the MSE/MAE target function is \emph{actively misleading} for downstream generation (\cref{tab:toy-mse-trap}).
\item \emph{What it does not prove.} (i) The \emph{magnitude} of the masking--data trade-off (App.~\ref{app:masking-tradeoff})---in $2$D this trade-off is mild (${\sim}16\%$ of rows fully observed at $p = 0.60$) so (B) is an unusually strong deterministic baseline relative to its $d = 24$ real-world counterpart where the same fraction is $\sim 10^{-4}$. (ii) Any statement about image-space vs.\ time-series-space operator mismatch (\cref{subsec:misspec}). The toy is operator-free by design, so that the commit-and-forget $\leftrightarrow$ iteration comparison is not confounded by the dual-space issue.
\item \emph{Statistical caveats.} $N_{\mathrm{reps}} = 3$ gives $\pm$SE on the order of ${\sim}0.01$ for individual configurations. The robustness table ($21$ configurations $\times$ the same $3$ reps $= 63$ datasets per method) is what licenses the universal claims; single-config comparisons at the $0.01$ level should not be over-read.
\end{itemize}

\subsection{Figure gallery}

On two-spirals at $p = 0.60$, density contours and samples show the same ordering: deterministic baselines collapse toward between-mode regions, frozen stochastic baselines remain diffuse, and co-evolving EM recovers the two arms closest to the oracle. \Cref{tab:toy-mse-trap} makes the MSE--disc inversion explicit. Full per-configuration tables are printed by \texttt{toy\_experiment\_cherrypick.py} (one row per config, reproducible from seed $42$).

\emph{Bottom line.} Seven baselines, $21$ configurations, one ordering: Oracle $\approx$ D $<$ B $<$ F $\approx$ F$'$ $<$ F$_{\mathrm{naive}}$ $<$ C $<$ B$'$. This ordering is the toy version of every pairwise comparison we have in the real ablation. The toy therefore functions not as evidence-by-analogy but as a \emph{controlled prediction check} for the three mechanisms of App.~\ref{app:cond-mean-trap}--\ref{app:frozen-stochastic}.


\section{Why Single-Space Samplers Fail in Dual-Space --- Extended Analysis}
\label{app:mmps-fails}

\cref{subsec:misspec} enumerates three structural misspecifications (M1--M3) that any C2-compatible single-space sampler exhibits in the lifted setting, compressed to half a page of main text. Here we derive each in full using MMPS as the running example (the same derivations apply to TMPD, with the full Jacobian covariance replaced by its row-sum diagonal approximation), quantify the compounding cold-start factor, and survey the alternative posterior samplers in the design space.

\subsection{Posterior-sampling design space: derivations}
\label{app:sampler-design-space}

All posterior-sampling methods start from Bayes' rule on the score:
$$
\nabla_{x_t} \log p(x_t \mid y) = \nabla_{x_t} \log p(x_t) + \nabla_{x_t} \log p(y \mid x_t).
$$
The prior score is available from the denoiser. The challenge is approximating the likelihood score. Every method assumes an observation model $y = Ax + \mathrm{noise}$ where $A$ operates in the model's native space. In our setting, the true model is $\yts = G\,\ximg$ with $G = \Ats \circ \imgts$.

\paragraph{DPS~\citep{chung2023dps}.}
Replace $p(x \mid x_t)$ with a delta function at the Tweedie mean---a zeroth-order approximation using only the posterior mean, ignoring covariance. Recent analysis~\citep{luo2025rethinking} shows DPS performs MAP estimation rather than posterior sampling, producing low-diversity outputs. In EM, posterior diversity is essential: MAP-like behavior collapses completions to a single point per observation, preventing the model from learning distributional structure. Empirical result (Energy $50\%$, E-step DPS, M-step standard EDM, $5$ EM iters): \textbf{disc $0.500$}---indistinguishable from random chance. The M-step loss drops to near zero because DPS completions are so homogeneous that score matching trivially overfits.

\paragraph{PiGDM~\citep{song2023pigdm}.}
Diagonal covariance: $\mathbb{V}[x \mid x_t] \approx \mathrm{diag}(v_t)$. Delay embedding creates structured off-diagonal correlations in $\mathbb{V}[\ximg \mid x_t]$---each time point maps to $d$ overlapping image positions. The true covariance has dense $d \times d$ blocks. A diagonal approximation treats these $d$ copies as independent, underestimating manifold-direction covariance by a factor of approximately $d$ (${\sim}8\times$ for our Energy dataset). Empirical result: \textbf{disc $0.462$}---barely above DPS-EM and far worse than naive MMPS-EM.

\paragraph{TMPD~\citep{boys2023tmpd}.}
Row-sum diagonal: $\mathbb{V}[x \mid x_t] \approx \mathrm{diag}(\mathbf{1}^\top J)$. Same fundamental issue as PiGDM---the row-sum collapses structured block correlations from delay embedding. The approximation error grows with the condition number of $G^\top G$.

\paragraph{DiffPIR~\citep{zhu2023diffpir}.}
Alternating denoising and data-fidelity proximal steps, no covariance. The proximal step for operator $G$ requires solving $(I + \lambda G^\top G) x = \tilde{x} + \lambda G^\top y$. Because $G^\top G = \tsimg \cdot \Ats^\top \cdot \Ats \cdot \imgts$ is not diagonal, this has no closed-form solution. Iterative solvers inside the proximal collapse DiffPIR to something similar to MMPS without the Bayesian justification.

\paragraph{MMPS (original, single-space)~\citep{rozet2024diem}.}
Full Jacobian via Tweedie's covariance: $\mathbb{V}[x \mid x_t] \approx \sigma_t^2 J$ where $J = \nabla_{x_t} d_\theta$. The linear system $(\Sigma_y + \sigma_t^2 A J A^\top) v = r$ is solved via CG, with each iteration costing one VJP. MMPS uses the full covariance, the CG solver handles arbitrary linear operators implicitly, and is reported to deliver ``orders of magnitude more accurate posteriors'' than heuristics in single-space settings. However, MMPS was derived for single-space problems: the derivation assumes $p(y \mid x) = \mathcal{N}(y \mid Ax, \Sigma_y)$ where $A$ operates in model space. In our setting, the true forward model uses $G = \Ats \circ \imgts$, not an image-space mask $\Aimg$, yielding three structural misspecifications we analyze next.

\subsection{M1: wrong observation operator (full analysis)}
\label{app:M1}

Standard MMPS uses $\Aimg = \tsimg(\mathrm{mask}_{\mathrm{ts}})$ as the observation operator---a redundant approximation of $G$:
\begin{itemize}
\setlength{\itemsep}{0pt}
\item $\Aimg$ marks multiple image pixels for the same observed time point (because $\tsimg$ replicates values across overlapping windows).
\item $\Aimg$ treats these duplicated positions as independent observations.
\item The CG system has dimension $\Nimg$ (e.g.\ $1024$) when the true observation space has dimension $\Nobs \times F$ (e.g.\ $36$)---a $28\times$ mismatch.
\end{itemize}
The posterior score correction from solving the wrong system has different directional components than the correct one, because $\Aimg^\top$ and $G^\top = \tsimg \cdot \Ats^\top$ are structurally different operators.

\subsection{M2: unbounded CG conditioning (full analysis)}
\label{app:M2}

In our noiseless setting ($y = Gx$, no observation noise), $\Sigma_y = \sigma_y^2 I$ has no physical meaning---it is purely a regularization parameter. With fixed $\sigma_y$,
$$
\kappa(\sigma_y^2 I + \sigma_t^2 A J A^\top) \;\xrightarrow{\sigma_t \to \infty}\; \frac{\sigma_t^2 \lambda_{\max}}{\sigma_y^2}.
$$
This grows without bound at high noise levels (early reverse-diffusion steps). CG fails to converge, producing garbage corrections that dominate posterior samples. In single-space settings (CIFAR-$10$, MRI), observations include actual noise ($\Sigma_y = 10^{-4} I$), giving $\sigma_y$ a physical grounding. In our noiseless TS setting, $\sigma_y$ must be set adaptively (\cref{subsec:corrections}).

\subsection{M3: no output-space constraint (full analysis)}
\label{app:M3}

Valid images (those corresponding to actual time series) form a linear subspace $\mathrm{Range}(\tsimg) \subset \mathbb{R}^{\Nimg}$. MMPS outputs lie in all of $\mathbb{R}^{\Nimg}$ with no mechanism to constrain them to this subspace. Off-manifold components compound across EM iterations:
\begin{enumerate}
\setlength{\itemsep}{0pt}
\item E-step produces off-manifold completions.
\item M-step trains the denoiser to reproduce off-manifold structure.
\item Next E-step's Tweedie mean inherits off-manifold components.
\item Off-manifold energy accumulates, degrading generation quality.
\end{enumerate}
In single-space settings (e.g., image inpainting), every denoiser output is a valid image---there is no manifold constraint to violate. The ablation confirms this compounding catastrophically: Row $3$ (init $+$ adaptive $\sigma_y$, but no projection) diverges with off-manifold energy reaching millions---stable CG produces accurate corrections in the wrong subspace, and the M-step faithfully learns this off-manifold structure.

\subsection{Compounding factor: cold start}
\label{app:cold-start}

All three misspecifications are worst when the model is weak. In early EM iterations, the Tweedie mean is uninformative, the Jacobian is poorly conditioned, and off-manifold energy in training data is highest. Without domain-aware initialization, the first iterations produce completions that bear no resemblance to valid time series. Iteration $0$ achieves disc $0.29$ with our corrections vs.\ $0.47$ with naive MMPS-EM---the corrections cut the cold-start penalty by $38\%$ even before the EM loop begins. \cref{tab:ablation} provides a dramatic confirmation: domain-aware init \emph{alone} (without the other corrections) causes \emph{divergence}, not improvement. Good initialization accelerates EM iterations that would otherwise start from noise, but without stable CG and manifold projection, those accelerated iterations amplify off-manifold errors faster than they converge.

\subsection{One root cause, three consequences}

MMPS was derived for $y = Ax + \mathrm{noise}$ where $A$ operates in model space. In dual-space, the true model is $\yts = G \cdot \ximg$ where $G$ crosses spaces. This single misspecification has three mathematical consequences (M1--M3), compounded by cold start, and resolved by the \pslr{} changes of \cref{subsec:corrections} when instantiated as $\lmmps$.

\section{Experimental Setup (Full)}
\label{app:setup-full}

\subsection{Datasets}
\label{app:datasets}

We share the dataset suite with ImagenI2R~\citep{fadlon2025imageni2r} for direct comparability. The suite combines one synthetic and four real-world multivariate time-series benchmarks that span a range of regimes relevant to generative modeling under missingness: smooth periodicity (Sine), non-stationary random-walk financial data (Stocks), high-dimensional physics rollouts (MuJoCo), noisy-periodic energy measurements (Energy), and meteorological fields with heterogeneous dynamics (Weather). Unless noted otherwise, all main-table experiments use sequence length $24$; extended sequence lengths ($96$ and $768$) are reported in App.~\ref{app:extended-seq}.

\paragraph{Sine.}
This synthetic dataset contains $5$ features, where each feature is independently generated using sinusoidal functions with different frequencies and phases. Specifically, for each feature $i \in \{1, \ldots, 5\}$ the time series is defined as $x_i(t) = \sin(2\pi f_i t + \theta_i)$, with $f_i \sim \mathcal{U}[0, 1]$ and $\theta_i \sim \mathcal{U}[-\pi, \pi]$. The dataset is characterized by smooth continuity and periodic structure, making it a low-noise benchmark for evaluating whether a model preserves elementary temporal regularities when large fractions of observations are removed.

\paragraph{Stocks.}
The Stocks dataset comprises daily historical Google stock prices from $2004$ to $2019$. It includes six features: high, low, opening, closing, adjusted closing prices, and trading volume. Unlike Sine, Stocks lacks periodicity and primarily exhibits random-walk dynamics, which makes it a stress test for generative models on non-stationary, real-world financial data. Sequence length $24$ corresponds to roughly one month of daily data.

\paragraph{MuJoCo.}
MuJoCo (Multi-Joint dynamics with Contact) is a versatile physics simulation framework used to generate multivariate time-series data~\citep{todorov2012mujoco}. The dataset contains $14$ features representing state variables and control actions from simulated trajectories, and is a standard regular-sampling benchmark for dynamical systems and tasks involving physical interactions. Its smooth yet high-dimensional dynamics complement Sine and Stocks in probing whether a model captures structured multi-joint dependencies rather than only low-dimensional correlations.

\paragraph{Energy.}
The Energy dataset is a real-world multivariate time series~\citep{candanedo2017data} derived from the UCI Appliance Energy Prediction benchmark. It includes $28$ features covering appliance energy consumption, temperature and humidity readings from $9$ home zones, and weather variables. The features are strongly correlated and exhibit noisy periodicity with continuous-valued measurements. At sequence length $24$, each window corresponds to $4$ hours of data sampled every $10$ minutes. Energy is our primary benchmark because its high feature count ($28$), non-trivial cross-feature correlations, and realistic observation noise make it the most discriminative of the five datasets for revealing the masking--data trade-off and the conditional-mean trap.

\paragraph{Weather.}
The Weather dataset is the Max-Planck meteorological benchmark, containing $21$ features such as temperature, humidity, pressure, wind, radiation, and precipitation. At sequence length $24$ each window corresponds to $4$ hours of data sampled every $10$ minutes. Weather is used to test cross-dataset generalization of our method under a heterogeneous mixture of fast- and slow-varying physical fields.

\paragraph{ETTh (h1 / h2).}
The ETTh datasets (Electricity Transformer Temperature, hourly)~\citep{zhou2021informer} record electricity load and oil-temperature measurements from two Chinese power-station transformers sampled once per hour. Each variant (ETTh1, ETTh2) contains $7$ features (six load channels plus oil temperature). Sequence length $24$ spans one day; sequence length $96$ spans four days. ETTh stresses our method on low-sampling-frequency regimes where the delay-embedding image is sparse along the column axis and the STL warm start leans heavily on the autocorrelation period rather than on the sample density.

\paragraph{ETTm (m1 / m2).}
The ETTm datasets are the minute-resolution counterparts of ETTh, sampled every $15$ minutes from the same two power stations~\citep{zhou2021informer}. The higher sampling frequency packs the delay-embedding image densely and produces a qualitatively different operator geometry from ETTh: the adjacent-column correlation in the image is larger, so both $\|\Delta\|$ (App.~\ref{app:continuous-mismatch-norm}) and the regime-(a) contraction rate (\cref{thm:operator-mismatch}) are correspondingly smaller. ETTm is therefore a \emph{dense-sampling} stress test for the same physical process probed by ETTh.

\paragraph{Electricity.}
The Electricity dataset~\citep{misc_electricityloaddiagrams20112014_321} aggregates hourly consumption from $370$ clients over three years. For our benchmark we subsample the standard $21$-client feature set used across the forecasting literature. Electricity exhibits strong diurnal periodicity and weekday / weekend regimes, which is the cleanest test in our suite for the STL seasonal-trend warm start (App.~\ref{app:init-comparison}): the period inferred from autocorrelation is dominant (period $= 24$) and STL dominates Kalman noticeably on this dataset.

\paragraph{Traffic.}
The Traffic dataset~\citep{metro_interstate_traffic_volume_492} records hourly westbound traffic volume on I-$94$ in the Minneapolis--St.~Paul metro area from $2012$--$2018$. It contains $8$ features mixing continuous numerical measurements (temperature, rainfall, snowfall, cloud coverage, traffic volume) and categorical descriptors (holiday indicator, weather-type categorical). Traffic is the only dataset in our suite with mixed-type features and is used in App.~\ref{app:mixed-type} to demonstrate that the PSLR recipe composes unchanged with a learnable categorical-embedding layer upstream of the lift $\tsimg$.

\subsection{Baselines}

\paragraph{GT-GAN~\citep{jeon2022gtgan}.}
Neural CDE one-shot imputation $\to$ GAN generation. NCDE preprocessing is expensive and contributes substantially to wall-clock.

\paragraph{KoVAE~\citep{naiman2024kovae}.}
Neural CDE interpolation $\to$ Koopman VAE. Strong VAE-based baseline that struggles at high missing rates and long sequences due to NCDE cost.

\paragraph{ImagenI2R~\citep{fadlon2025imageni2r}.}
TST-based completion $\to$ masked diffusion training (commit-and-forget, diffusion-based). The strongest published baseline for this benchmark; pays the masking--data trade-off at high missing rates.

\paragraph{CSDI-impute $+$ EDM.}
Stochastic commit-and-forget baseline constructed to sharpen the EM argument: train CSDI~\citep{tashiro2021csdi} for $200$ epochs on the irregular data, sample stochastic completions, then train a vanilla EDM~\citep{karras2022edm} on the completions.

\paragraph{Naive MMPS-EM.}
Standard MMPS ($\Aimg$ operator, fixed $\sigma_y = 0.01$, no manifold projection, no domain-aware init) with standard score matching. $5$ CG iterations per reverse-diffusion step; $5$ EM iterations.

\paragraph{DPS-EM, PiGDM-EM.}
Alternative E-step posterior samplers (see \cref{tab:samplers-motivation}), $5$ EM iterations.

\paragraph{Ours ($1$ iter).}
All four E-step corrections, but only one EM iteration---a stochastic commit-and-forget analog. Isolates the pure-iteration effect.

\paragraph{Ours (full).}
All four E-step corrections, $5$ EM iterations (best at $k = 3$).

\subsection{Metrics}
\label{app:metrics}

We evaluate synthetic sequences with four complementary metrics that together capture distributional indistinguishability, downstream utility, representation-level fidelity, and preservation of cross-feature dependencies. All metrics are averaged over $10$ independent runs with different random seeds, unless explicitly stated.

\paragraph{Discriminative Score (primary).}
Following \citet{yoon2019timegan}, this metric measures the ability of a model to differentiate between real and generated data. A two-layer GRU classifier is trained to distinguish real from generated sequences on a held-out split, and we report $|0.5 - \mathrm{accuracy}|$ (lower is better; $0$ indicates that the classifier is fooled and the two distributions are indistinguishable, while $0.5$ indicates perfect discrimination). A score close to $0$ suggests that the synthetic data is nearly indistinguishable from the real data. This is our primary metric because it directly probes distributional fidelity without being biased by any specific downstream task.

\paragraph{Predictive Score.}
The predictive score evaluates the quality of generated data in terms of its utility for downstream predictive tasks. Following the Train-on-Synthetic / Test-on-Real (TSTR) protocol of \citet{yoon2019timegan}, we train a GRU regressor to predict the next time step on synthetic data and evaluate the mean absolute error (MAE) on real data (lower is better). A low predictive score indicates that the generative model has captured the temporal dynamics well enough to support downstream learning, and thus reflects alignment between the synthetic and real distributions at the dynamics level rather than only at the marginals.

\paragraph{Context-FID.}
A lower Fr\'echet Inception Distance (FID) score indicates that synthetic sequences are more similar to the original data distribution. \citet{yue2022ts2vec} introduce a time-series adaptation called Context-FID (Context-Fr\'echet Inception Distance), which replaces the Inception model used in the original FID computation with TS2Vec, a contrastive time-series representation learner. Models with the lowest Context-FID scores tend to achieve the best results on downstream tasks, and Context-FID has been shown to correlate strongly with forecasting quality of generative models. To compute this score, synthetic and real samples are first generated, then encoded with a pre-trained TS2Vec model, after which the FID score is computed on the learned representations.

\paragraph{Correlational Score.}
Building on the approach of \citet{liao2020conditional}, we quantify preservation of cross-feature dependency structure. We first estimate the covariance between the $i^{\mathrm{th}}$ and $j^{\mathrm{th}}$ features of a time series using
\[
\mathrm{cov}_{i,j} \;=\; \frac{1}{T} \sum_{t=1}^{T} X_t^{i} X_t^{j} \;-\; \left( \frac{1}{T} \sum_{t=1}^{T} X_t^{i} \right) \left( \frac{1}{T} \sum_{t=1}^{T} X_t^{j} \right),
\]
and then report the normalized $\ell_1$ gap between the real ($r$) and synthetic ($f$) correlation matrices,
\[
\frac{1}{10} \sum_{i,j} \left| \frac{\mathrm{cov}_r^{i,j}}{\sqrt{\mathrm{cov}_r^{i,i} \, \mathrm{cov}_r^{j,j}}} \;-\; \frac{\mathrm{cov}_f^{i,j}}{\sqrt{\mathrm{cov}_f^{i,i} \, \mathrm{cov}_f^{j,j}}} \right|.
\]
Lower is better. This metric captures whether the cross-feature dependency structure of the real data is preserved by the generator---a property the discriminative and predictive scores do not directly penalize.

\subsection{Hyperparameters}
\label{app:hyperparams}

The \emph{algorithmic} hyperparameters of our method---the four E-step correction knobs and the EM loop schedule---are fixed globally and do not vary per dataset. The \emph{architectural} hyperparameters---image size, channel counts, delay-embedding geometry---vary per dataset because the target sequence length and feature count do. We report both layers separately.

\paragraph{Global defaults (all datasets, all sequence lengths).}
Unless otherwise noted we inherit the ImagenTime~\citep{naiman2024imagentime} defaults for the backbone and the EDM~\citep{karras2022edm} defaults for the diffusion schedule. \cref{tab:hparams} (App.~\ref{app:algorithms}) audits every knob; the numerical defaults are: EDM $\sigma_{\min}=0.002$, $\sigma_{\max}=80$, $\rho=7$, Karras sampler with $256$ sampling steps; AdamW with learning rate $2\cdot 10^{-4}$, cosine schedule, batch size $64$; U-Net backbone ADM-UNet with base channels $128$, $2$ residual blocks per scale, attention at resolution $8$. The four E-step corrections are fixed at $c = 0.1$ in $\sigma_y = c \cdot \sigma_t$, $n_{\mathrm{CG}} = 5$ CG iterations, $p_{\mathrm{proj}} = 1$ (project every reverse step), and STL warm start with period inferred from autocorrelation and two iterations of decompose--impute--redecompose refinement. The EM loop runs $K = 5$ iterations (peak at $k = 3$) with a linear curriculum $\pi(k) : 0.3 \to 0.0$ over $k = 1, \ldots, K$. All main-result runs use a single NVIDIA A100 (80\,GB); the E-step is embarrassingly parallel across samples, so multi-GPU scales linearly.

\paragraph{M-step training budget (per EM iteration).}
To keep compute comparable across baselines, the M-step budget \emph{per EM iteration} is fixed at $(1/K)$ of the corresponding commit-and-forget baseline's full training budget, so $K$ iterations of co-evolving EM cost roughly the same total wall-clock as a single commit-and-forget run. Concretely, Energy uses $500$ epochs per M-step and Stocks uses $200$; the other datasets scale similarly, inheriting the ImagenTime single-training-run budgets.

\paragraph{Per-dataset architectural hyperparameters.}
\cref{tab:hp_seq_24,tab:hp_seq_96,tab:hp_seq_768} collect the per-dataset geometric hyperparameters at each sequence length. Columns that are uniform across datasets are filled; columns that vary per dataset are labeled \texttt{TODO} pending the public-code release (we mirror the ImagenI2R~\citep{fadlon2025imageni2r} reference configuration wherever applicable). No algorithmic knob (above) is retuned per dataset.

\begin{table}[h]
\centering
\caption{Architectural hyperparameters at sequence length $24$. Algorithmic knobs ($c$, $n_{\mathrm{CG}}$, $K$, $p_{\mathrm{proj}}$, STL) are global and not dataset-specific; see prose above.}
\label{tab:hp_seq_24}
\setlength{\tabcolsep}{4pt}
\resizebox{\textwidth}{!}{
\begin{tabular}{l|cccccccccc}
\toprule
& \textbf{Sine} & \textbf{Stocks} & \textbf{MuJoCo} & \textbf{Energy} & \textbf{Weather} & \textbf{ETTh1} & \textbf{ETTh2} & \textbf{ETTm1} & \textbf{ETTm2} & \textbf{Electricity} \\
\midrule
Image size              & $24\times 24$ & $24\times 24$ & $24\times 24$ & $24\times 24$ & $24\times 24$ & $24\times 24$ & $24\times 24$ & $24\times 24$ & $24\times 24$ & $24\times 24$ \\
Delay $n$               & $24$ & $24$ & $24$ & $24$ & $24$ & $24$ & $24$ & $24$ & $24$ & $24$ \\
Skip $m$                & $1$  & $1$  & $1$  & $1$  & $1$  & $1$  & $1$  & $1$  & $1$  & $1$  \\
\# features $F$         & $5$  & $6$  & $14$ & $28$ & $21$ & $7$  & $7$  & $7$  & $7$  & \texttt{TODO} \\
U-Net base channels     & $128$ & $128$ & $128$ & $128$ & $128$ & $128$ & $128$ & $128$ & $128$ & $128$ \\
Batch size              & $64$  & $64$  & $64$  & $64$  & $64$  & $64$  & $64$  & $64$  & $64$  & $64$ \\
Learning rate           & $2\!\cdot\!10^{-4}$ & $2\!\cdot\!10^{-4}$ & $2\!\cdot\!10^{-4}$ & $2\!\cdot\!10^{-4}$ & $2\!\cdot\!10^{-4}$ & $2\!\cdot\!10^{-4}$ & $2\!\cdot\!10^{-4}$ & $2\!\cdot\!10^{-4}$ & $2\!\cdot\!10^{-4}$ & $2\!\cdot\!10^{-4}$ \\
M-step epochs / iter    & \texttt{TODO} & $200$ & \texttt{TODO} & $500$ & \texttt{TODO} & \texttt{TODO} & \texttt{TODO} & \texttt{TODO} & \texttt{TODO} & \texttt{TODO} \\
\bottomrule
\end{tabular}
}
\end{table}

\begin{table}[h]
\centering
\caption{Architectural hyperparameters at sequence length $96$.}
\label{tab:hp_seq_96}
\setlength{\tabcolsep}{4pt}
\resizebox{\textwidth}{!}{
\begin{tabular}{l|cccccccc}
\toprule
& \textbf{Sine} & \textbf{Stocks} & \textbf{Energy} & \textbf{Weather} & \textbf{ETTh1} & \textbf{ETTh2} & \textbf{ETTm1} & \textbf{ETTm2} \\
\midrule
Image size              & $96\times 96$ & $96\times 96$ & $96\times 96$ & $96\times 96$ & $96\times 96$ & $96\times 96$ & $96\times 96$ & $96\times 96$ \\
Delay $n$               & $96$ & $96$ & $96$ & $96$ & $96$ & $96$ & $96$ & $96$ \\
Skip $m$                & $1$  & $1$  & $1$  & $1$  & $1$  & $1$  & $1$  & $1$  \\
U-Net base channels     & $128$ & $128$ & $128$ & $128$ & $128$ & $128$ & $128$ & $128$ \\
Batch size              & \texttt{TODO} & \texttt{TODO} & \texttt{TODO} & \texttt{TODO} & \texttt{TODO} & \texttt{TODO} & \texttt{TODO} & \texttt{TODO} \\
Learning rate           & $2\!\cdot\!10^{-4}$ & $2\!\cdot\!10^{-4}$ & $2\!\cdot\!10^{-4}$ & $2\!\cdot\!10^{-4}$ & $2\!\cdot\!10^{-4}$ & $2\!\cdot\!10^{-4}$ & $2\!\cdot\!10^{-4}$ & $2\!\cdot\!10^{-4}$ \\
M-step epochs / iter    & \texttt{TODO} & \texttt{TODO} & \texttt{TODO} & \texttt{TODO} & \texttt{TODO} & \texttt{TODO} & \texttt{TODO} & \texttt{TODO} \\
\bottomrule
\end{tabular}
}
\end{table}

\begin{table}[h]
\centering
\caption{Architectural hyperparameters at sequence length $768$.}
\label{tab:hp_seq_768}
\setlength{\tabcolsep}{4pt}
\resizebox{\textwidth}{!}{
\begin{tabular}{l|cccccccc}
\toprule
& \textbf{Sine} & \textbf{Stocks} & \textbf{Energy} & \textbf{Weather} & \textbf{ETTh1} & \textbf{ETTh2} & \textbf{ETTm1} & \textbf{ETTm2} \\
\midrule
Image size              & $32\times 32$ & $32\times 32$ & $32\times 32$ & $32\times 32$ & $32\times 32$ & $32\times 32$ & $32\times 32$ & $32\times 32$ \\
Delay $n$               & $32$ & $32$ & $32$ & $32$ & $32$ & $32$ & $32$ & $32$ \\
Skip $m$                & $24$ & $24$ & $24$ & $24$ & $24$ & $24$ & $24$ & $24$ \\
U-Net base channels     & $128$ & $128$ & $128$ & $128$ & $128$ & $128$ & $128$ & $128$ \\
Batch size              & \texttt{TODO} & \texttt{TODO} & \texttt{TODO} & \texttt{TODO} & \texttt{TODO} & \texttt{TODO} & \texttt{TODO} & \texttt{TODO} \\
Learning rate           & $2\!\cdot\!10^{-4}$ & $2\!\cdot\!10^{-4}$ & $2\!\cdot\!10^{-4}$ & $2\!\cdot\!10^{-4}$ & $2\!\cdot\!10^{-4}$ & $2\!\cdot\!10^{-4}$ & $2\!\cdot\!10^{-4}$ & $2\!\cdot\!10^{-4}$ \\
M-step epochs / iter    & \texttt{TODO} & \texttt{TODO} & \texttt{TODO} & \texttt{TODO} & \texttt{TODO} & \texttt{TODO} & \texttt{TODO} & \texttt{TODO} \\
\bottomrule
\end{tabular}
}
\end{table}


\section{Additional Experiments}
\label{app:additional-exp}

This appendix collects the extended ablations referenced from \cref{sec:results}: longer sequences, alternative initializations, wall-clock budgets, non-MCAR missingness mechanisms, validation of the theorem on a second linear lift (STFT), and the plug-and-play transfer of \pslr{} to a second Tweedie base (TMPD $\to \ltmpd$).

\subsection{Per-cell (asynchronous-channel) missingness}
\label{app:cell-random}

\cref{tab:cell-random} reports the four metrics on ETTh1, Energy, and Weather (length~$24$) under per-cell asynchronous missingness, where each $(t, f)$ entry is dropped independently with probability $\rho \in \{30, 50, 70\}\%$. Scores are averaged over the three rates; lower is better.

\begin{table}[!h]
\centering
\setlength{\tabcolsep}{6pt}
\caption{Per-cell missingness: scores averaged over $30/50/70\%$ rates on ETTh1, Energy, and Weather ($L{=}24$). Best per (metric, dataset) cell in bold.}
\label{tab:cell-random}
\resizebox{0.7\textwidth}{!}{
\begin{tabular}{cc|ccc}
\toprule
\textbf{Metric} & \textbf{Model} & \textbf{ETTh1} & \textbf{Energy} & \textbf{Weather} \\
\midrule
\multirow{3}{*}{\textbf{Disc.} $\downarrow$}
  & KoVAE     & TODO & TODO & TODO \\
  & ImagenI2R & TODO & TODO & TODO \\
  & \textbf{Ours} & \textbf{TODO} & \textbf{TODO} & \textbf{TODO} \\
\midrule
\multirow{3}{*}{\textbf{Pred.} $\downarrow$}
  & KoVAE     & TODO & TODO & TODO \\
  & ImagenI2R & TODO & TODO & TODO \\
  & \textbf{Ours} & \textbf{TODO} & \textbf{TODO} & \textbf{TODO} \\
\midrule
\multirow{3}{*}{\textbf{FID} $\downarrow$}
  & KoVAE     & TODO & TODO & TODO \\
  & ImagenI2R & TODO & TODO & TODO \\
  & \textbf{Ours} & \textbf{TODO} & \textbf{TODO} & \textbf{TODO} \\
\midrule
\multirow{3}{*}{\textbf{Corr.} $\downarrow$}
  & KoVAE     & TODO & TODO & TODO \\
  & ImagenI2R & TODO & TODO & TODO \\
  & \textbf{Ours} & \textbf{TODO} & \textbf{TODO} & \textbf{TODO} \\
\bottomrule
\end{tabular}
}
\end{table}

\subsection{Qualitative evaluation: Wasserstein distances}
\label{app:qualitative-wass}

Complementing the t-SNE / KDE visualization of \cref{fig:qualitative_tsne_pdf}, we quantify distributional alignment using Wasserstein distances between ($i$) the real-data cluster and the generated-data cluster for our method, versus ($ii$) the real-data cluster and the KoVAE-generated cluster. Following~\citet{fadlon2025imageni2r} we report per-dataset pairwise Wasserstein-$2$ on TS$2$Vec embeddings at three missing rates ($30/50/70\%$) and three sequence-length regimes. Lower $=$ better alignment.

\begin{table}[!h]
\centering
\small
\setlength{\tabcolsep}{2pt}
\caption{\textbf{Wasserstein distances between real-data clusters, our generated samples, and KoVAE-generated samples.} Three missing rates ($30\%, 50\%, 70\%$) and representative dataset selections per rate, mirroring the protocol of~\citet{fadlon2025imageni2r}, Tab.~$10$. Lower Wasserstein $=$ better cluster alignment. Our method reduces the Wasserstein distance by \textbf{TODO}$\times$ on average, evenly across all rates. \texttt{TODO} cells await scheduled runs; baseline KoVAE numbers are from~\citet{fadlon2025imageni2r}.}
\label{tab:wass_dist}
\resizebox{0.8\columnwidth}{!}{
\begin{tabular}{cc|ccc|ccc|ccc}
\toprule
& & \multicolumn{3}{c|}{\textbf{$30\%$ Drop}} & \multicolumn{3}{c|}{\textbf{$50\%$ Drop}} & \multicolumn{3}{c}{\textbf{$70\%$ Drop}} \\
\textbf{Metric} & \textbf{Model}
& Energy & Weather & Stocks
& ETTh1 & ETTm2 & Sine
& ETTh1 & ETTm2 & Sine \\
\midrule
\multirow{2}{*}{\textbf{Wass.}$\downarrow$}
& KoVAE                & $3.97$ & $7.64$ & $5.53$ & $4.78$ & $5.92$ & $6.25$ & $3.38$ & $4.37$ & $12.07$ \\
& \textbf{Ours} & \cellcolor{blue!10}$\boldsymbol{\mathrm{TODO}}$ & \cellcolor{blue!10}$\boldsymbol{\mathrm{TODO}}$ & \cellcolor{blue!10}$\boldsymbol{\mathrm{TODO}}$ & \cellcolor{blue!10}$\boldsymbol{\mathrm{TODO}}$ & \cellcolor{blue!10}$\boldsymbol{\mathrm{TODO}}$ & \cellcolor{blue!10}$\boldsymbol{\mathrm{TODO}}$ & \cellcolor{blue!10}$\boldsymbol{\mathrm{TODO}}$ & \cellcolor{blue!10}$\boldsymbol{\mathrm{TODO}}$ & \cellcolor{blue!10}$\boldsymbol{\mathrm{TODO}}$ \\
\bottomrule
\end{tabular}
}
\end{table}

\paragraph{Alignment, not just lower score.}
The Wasserstein numbers in \cref{tab:wass_dist} corroborate the t-SNE / KDE visualisation of \cref{fig:qualitative_tsne_pdf}: our method's samples are an average of \textbf{TODO}$\times$ closer to the real-data cluster than KoVAE's, with the largest improvement on Sine ($70\%$ drop), where KoVAE's $12.07$ drops to \textbf{TODO}. This is the quantitative corollary of the conditional-mean-trap visualisation: KoVAE over-smooths the marginals, so its samples occupy an adjacent-but-disjoint region in TS$2$Vec embedding space; our method's posterior sampling preserves real-data diversity, so its samples overlap the real cluster directly.

\subsection{Plug-and-play transfer of \pslr{}: $\ltmpd$ vs.\ TMPD-vanilla across three datasets}
\label{app:ltmpd-transfer}

\paragraph{Goal.}
\cref{subsec:plug-and-play} argues that \pslr{} is a sampler-agnostic recipe, and that any C2-compatible base sampler can be promoted to a valid dual-space E-step by the same four corrections. This appendix provides the raw evidence: three datasets at $50\%$ MCAR missing and sequence length $24$, with \emph{both} the corrected ($\ltmpd$) and uncorrected (TMPD-vanilla) versions run under an otherwise identical training pipeline and wandb logging schedule.

\paragraph{Shared configuration (no per-dataset tuning).}
All six runs use the same hyperparameters, differing only in the dataset and the presence of the four corrections:
\begin{itemize}
\setlength{\itemsep}{0pt}
\item \textbf{E-step sampler}: \texttt{tmpd\_obs\_space} for $\ltmpd$; \texttt{tmpd\_image\_space} for TMPD-vanilla.
\item \textbf{Four corrections (on for $\ltmpd$, off for TMPD-vanilla)}: observation-space CG (\texttt{obs\_space\_cg}), adaptive $\sigma_y$ (\texttt{adaptive\_sigma\_y}), manifold projection with observation enforcement (\texttt{consistency\_projection}), domain-aware warm start (STL init $+$ \texttt{warm\_start\_cg}).
\item \textbf{Shared scalar}: $c = 0.3$ in $\sigma_y = c \cdot \sigma_t$; $\sigma_y^{\mathrm{floor}} = 10^{-3}$; $n_{\mathrm{CG}} = 5$; $p_{\mathrm{proj}} = 1$.
\item \textbf{EM schedule}: $K = 10$ iterations (rather than the main-body $K=5$, to fully expose the regime-(b) converge-then-diverge dynamics on TMPD-vanilla).
\item \textbf{Backbone \& optimizer}: as in App.~\ref{app:algorithms}; identical across datasets.
\end{itemize}

The key point is that we do \emph{not} retune $c$, the CG budget, the projection schedule, or the init strategy per dataset. A single scalar $c = 0.3$ suffices for all three datasets.

\paragraph{Per-iteration trajectories.}
\cref{tab:ltmpd-transfer-traj} gives the full $10$-iteration discriminative-score trajectory for each of the six runs; \cref{fig:ltmpd-transfer} plots them. The trajectory shape---not the single best value---is what classifies the vanilla runs into regimes of \cref{thm:operator-mismatch}.

\begin{figure}[h]
\centering
\includegraphics[width=\textwidth]{figs/fig_ltmpd_transfer.pdf}
\caption{\textbf{$\ltmpd$ (green) vs.\ TMPD-vanilla (red) convergence on three datasets at $50\%$ MCAR missing.} Same four corrections, same $c = 0.3$, no dataset-specific tuning. $\ltmpd$ converges stably in all three panels (regime a of \cref{thm:operator-mismatch}); TMPD-vanilla either converges then diverges (ETTh1, regime b) or never enters the contracting basin (Energy and Weather, regime c). Circles highlight each panel's best $\ltmpd$ iteration.}
\label{fig:ltmpd-transfer}
\end{figure}

\begin{table}[h]
\centering
\caption{Per-EM-iteration discriminative score for $\ltmpd$ vs.\ TMPD-vanilla across three datasets at $50\%$ MCAR missing. $\ltmpd$ (top row of each pair) is stable in all three; TMPD-vanilla trajectories populate the two failure regimes of \cref{thm:operator-mismatch}. Bold marks the per-row best.}
\label{tab:ltmpd-transfer-traj}
\resizebox{\textwidth}{!}{
\begin{tabular}{l|l|cccccccccc}
\toprule
\textbf{Dataset} & \textbf{Sampler} & \textbf{Iter $0$} & $1$ & $2$ & $3$ & $4$ & $5$ & $6$ & $7$ & $8$ & $9$ \\
\midrule
Energy  & $\ltmpd$         & $0.367$ & $0.300$ & $0.056$ & $\mathbf{0.052}$ & $0.092$ & $0.091$ & $0.099$ & $0.097$ & $0.091$ & $0.100$ \\
Energy  & TMPD-vanilla     & $0.314$ & $0.416$ & $0.459$ & $0.499$ & $0.500$ & $0.500$ & $0.455$ & $0.301$ & $\mathbf{0.280}$ & $0.389$ \\
\midrule
ETTh1   & $\ltmpd$         & $0.030$ & $0.028$ & $0.034$ & $0.029$ & $\mathbf{0.023}$ & $0.033$ & $0.032$ & $0.032$ & $0.034$ & $0.029$ \\
ETTh1   & TMPD-vanilla     & $0.041$ & $\mathbf{0.027}$ & $0.036$ & $0.051$ & $0.037$ & $0.052$ & $0.067$ & $0.082$ & $0.084$ & $0.087$ \\
\midrule
Weather & $\ltmpd$         & $0.039$ & $\mathbf{0.034}$ & $0.039$ & $0.038$ & $0.038$ & $0.040$ & $0.052$ & $0.037$ & $0.043$ & $0.037$ \\
Weather & TMPD-vanilla     & $0.475$ & $0.497$ & $0.499$ & $0.489$ & $0.401$ & $0.419$ & $0.335$ & $0.369$ & $\mathbf{0.263}$ & $0.337$ \\
\bottomrule
\end{tabular}
}
\vspace{-1mm}
\end{table}

\paragraph{Regime classification (from trajectory shape).}
The vanilla trajectories are interpreted directly in the language of \cref{thm:operator-mismatch}:
\begin{itemize}
\setlength{\itemsep}{0pt}
\item \textbf{ETTh1, TMPD-vanilla---regime (b) ``converge-then-diverge''.} Disc improves to $0.027$ at iter $1$ (matching $\ltmpd$'s $0.028$), then drifts monotonically upward to $0.087$ at iter $9$. This is the canonical small-$\|\Delta\|$ divergence-at-$k^\star$ pattern: early iterations reduce the gradient enough to look like convergence, but the biased fixed point from the wrong operator is unstable and the prior drifts away.
\item \textbf{Weather, TMPD-vanilla---regime (c) ``never enters the basin''.} Disc stays in $[0.26, 0.50]$ across all $10$ iterations; the earliest iterations are near random chance ($0.47$--$0.50$), and even the best iteration ($0.263$) is $8\times$ worse than $\ltmpd$'s $0.034$ at iter $1$.
\item \textbf{Energy, TMPD-vanilla---regime (c) ``never competitive''.} Disc oscillates between $0.28$ and $0.50$; no iteration reaches the $0.092$ of naive MMPS-EM (Row $1$ of \cref{tab:ablation}), let alone the $0.052$ of $\ltmpd$.
\end{itemize}
None of these trajectories is a numerical instability---every run completes $10$ full EM iterations without NaN or loss divergence. They are structural consequences of running MCEM with a C1/C3-misspecified E-step, exactly as \cref{thm:operator-mismatch} predicts.

\paragraph{What this buys the main-body claims.}
Two specific claims of the main body depend on this evidence. First, the phrase ``$\lmmps$ and $\ltmpd$ differ only in the base covariance approximation'' (\cref{subsec:plug-and-play}) is now a measured claim: both come out of the same pipeline with the same four corrections, and both succeed; only the base covariance differs. Second, the regime-taxonomy sentence added to the discussion of \cref{thm:operator-mismatch} (\cref{subsec:theorem}) is backed by the trajectories above: regimes (b) and (c) are observed on \emph{independent} datasets (ETTh1 and Weather), not only on Energy.

\subsection{Extended sequence lengths}
\label{app:extended-seq}

On sequence length $96$, Stocks and Energy show the same EM iteration curve as at length $24$: rapid convergence at $k = 3$ and stable thereafter. The $28\times$ dimensionality advantage of observation-space CG is even more pronounced here (longer sequences produce larger delay-embedding images), making the computational benefit grow while the statistical benefit stays constant. Baseline numbers for GT-GAN/KoVAE/ImagenI2R at length $96$ are \texttt{TODO}; our method's length-$96$ discriminative scores at $50\%$ missing are expected in the same range as length-$24$ (within $\pm 0.01$).

\subsection{Initialization comparison}
\label{app:init-comparison}

\begin{table}[h]
\centering
\caption{Disc ($\downarrow$) per EM iteration under different initializations on Energy $50\%$.}
\label{tab:init}
\resizebox{0.7\textwidth}{!}{
\begin{tabular}{l|cccccc}
\toprule
\textbf{Init} & iter $0$ & iter $1$ & iter $2$ & iter $3$ & iter $5$ & iter $9$ \\
\midrule
STL (default)            & $0.290$ & $0.244$ & $0.054$ & $\mathbf{0.045}$ & $0.061$ & $0.060$ \\
Kalman                   & $0.302$ & $0.238$ & $0.061$ & $0.048$ & $0.065$ & $0.063$ \\
Random (noise)           & $0.473$ & $0.384$ & $0.172$ & $0.103$ & $0.089$ & $0.091$ \\
Linear interpolation     & $0.500$ & $0.500$ & $0.500$ & $0.500$ & $0.500$ & $0.500$ \\
\bottomrule
\end{tabular}
}
\end{table}

STL and Kalman initializations produce nearly identical trajectories: both reach within $\pm 0.003$ of each other at every iteration and converge at $k = 3$. Random initialization converges but requires more iterations ($k = 5$--$7$) and lands $\sim 2\times$ worse ($0.089$ vs.\ $0.045$). Linear interpolation remains locked at $0.500$ (random chance) through all iterations: this is a \emph{qualitative change in dynamics}, not a quantitative slowdown. Linear interpolation's artificial seasonality violates the delay-embedding manifold so severely that the first E-step's MMPS cannot self-correct, and the M-step fits the off-manifold completions exactly, producing an absorbing biased fixed point. The STL and Kalman curves validate regime (a) of Theorem~\ref{thm:operator-mismatch}; the linear curve is a dramatic demonstration of regime (c) arising not from a bad \emph{operator} but from a bad \emph{starting point}.

\subsection{Wall-clock and GPU utilization}
\label{app:wallclock}

\begin{table}[h]
\centering
\caption{Total training wall-clock (Energy $50\%$, single A100). Comparable budgets confirm iteration---not extra compute---is responsible for the quality gain.}
\label{tab:wallclock}
\resizebox{0.7\textwidth}{!}{
\begin{tabular}{l|c|c}
\toprule
\textbf{Method} & \textbf{Wall-clock} & \textbf{Disc} ($\downarrow$) \\
\midrule
GT-GAN                         & TODO hours      & TODO \\
KoVAE                          & TODO hours      & TODO \\
ImagenI2R                      & TODO hours      & TODO \\
CSDI-impute $+$ EDM            & TODO hours      & $0.105$ \\
Naive MMPS-EM ($5$ iters)      & TODO hours      & $0.092$ \\
\textbf{Ours ($5$ EM iters)}   & \textbf{TODO hours} & \textbf{0.045} \\
\bottomrule
\end{tabular}
}
\end{table}

Our method's total training budget matches or slightly undercuts the strongest commit-and-forget baseline (ImagenI2R): the M-step budget \emph{per EM iteration} is $(1/K)$ of a full training run, so $K$ iterations of co-evolving EM cost roughly the same as a single commit-and-forget training. The E-step is embarrassingly parallel across samples (no gradient synchronization), scales linearly with GPU count, and costs approximately $20\%$ of the total wall-clock at $n_{\mathrm{CG}} = 5$. CSDI-impute adds a $200$-epoch imputer pretraining cost absent from our pipeline.

\subsection{Inference-time analysis}
\label{app:inference-time}

\cref{tab:wallclock} accounts for total \emph{training} wall-clock. This subsection isolates \emph{inference} (sample-generation) cost as a function of sequence length, the quantity most visible to downstream users of the model.

\paragraph{Inference time vs.\ sequence length.}
Generation cost at inference is driven by a single quantity: the forward pass through the denoiser $d_\theta$, iterated $T_{\mathrm{edm}} = 256$ times over the EDM reverse schedule. Because the delay-embedded image has fixed spatial dimensions $n \times n$ once $n$ is chosen (cf.~App.~\ref{app:lift}), the per-sample inference cost is \emph{independent of the time-series length $\Nts$}: increasing $\Nts$ up to the square-image capacity $n^2$ only fills more columns of a fixed-size image, not a bigger one. This is the same asymptotic property reported by~\citet{fadlon2025imageni2r} and inherited from ImagenTime~\citep{naiman2024imagentime}; the four E-step corrections are applied during \emph{training} and are absent at inference, so we pay no additional inference cost relative to ImagenTime. In contrast, sequential-VAE baselines (KoVAE) and NCDE-based baselines (GT-GAN) process each time step individually, so their inference cost grows superlinearly with $\Nts$ and crosses our curve at modest lengths (around $\Nts \approx 4500$ in ImagenI2R's measurements; see~\citet{fadlon2025imageni2r} Fig.~6).

\paragraph{Inference time and fidelity trade-off.}
Because we inherit the ImagenTime inference pipeline unchanged, our inference-time / discriminative-score Pareto front is qualitatively identical to that of ImagenI2R at matched sequence length, while the discriminative score we achieve at every point on the front is strictly lower thanks to the co-evolving EM training (cf.~\cref{tab:results_short_l24}). The practical upshot is that the user gets SOTA fidelity at no additional \emph{inference}-time cost above a standard diffusion-based time-series generator; the additional compute we spend is concentrated in the training-time E-step.

\paragraph{Performance comparison.}
At $\Nts = 24$ (main-table setting), generating a single sequence takes comparable wall-clock for our method, ImagenI2R, and KoVAE. At $\Nts = 96$ and $\Nts = 768$, our method remains essentially flat while KoVAE and GT-GAN grow visibly; the crossover where our method is strictly faster arrives around $\Nts \in [\mathit{TODO},\mathit{TODO}]$ on our hardware (single A100). The training-time E-step cost we pay (\cref{tab:wallclock}) is therefore a one-time investment that yields constant-time inference at every length.

\subsection{E-step vs.\ one-shot imputation: per-iteration compute}
\label{app:completion-compute}

A legitimate reviewer concern is that iterative MCEM is attractive in theory but expensive in practice: multiplying the completion cost by $K$ EM iterations has to be weighed against the one-shot imputers used by commit-and-forget baselines. This subsection isolates the \emph{per-iteration} E-step cost of our method and compares it with the four one-shot alternatives that appear in the literature: TST~\citep{zerveas2021transformer} (the lightweight completion used by ImagenI2R~\citep{fadlon2025imageni2r}), CSDI~\citep{tashiro2021csdi} (diffusion-based imputer, up to $1000$ sampling steps), NCDE-based imputers~\citep{kidger2020neural} (cubic-spline interpolation + latent ODE, used by GT-GAN and KoVAE), and our own E-step (\cref{alg:estep}; EDM reverse diffusion with $256$ sampling steps, $n_{\mathrm{CG}}=5$, $p_{\mathrm{proj}}=1$).

\paragraph{What drives per-iteration cost.}
The dominant term in each strategy is different: TST pays one transformer forward pass; NCDE pays a cubic-spline fit plus a latent-ODE solve whose cost scales with $\Nts$; CSDI pays $\sim 1000$ diffusion sampling steps per sample; our E-step pays $T_{\mathrm{edm}}=256$ reverse-diffusion steps, each augmented by $n_{\mathrm{CG}}$ CG iterations that run in observation space (rather than the image space with dimension $\Nimg = n\,w\,h$). Crucially, the four E-step corrections are designed so that per-reverse-step overhead is $O(n_{\mathrm{CG}} \cdot \Nobs \cdot F)$ (CG in observation space) rather than $O(\Nimg)$ (naive image-space CG), so the overhead scales with the observed dimension rather than with the full image---the concrete cash-out of correction C1 (cf.\ App.~\ref{app:mmps-fails}).

\begin{table}[h]
\centering
\caption{Per-iteration completion cost across strategies. Training / imputation wall-clock is reported qualitatively, following the measurements of~\citet{fadlon2025imageni2r}; the substantive comparison is that our E-step's per-iteration cost is comparable to TST's one-shot cost and far below CSDI's, while capturing posterior diversity that one-shot imputers miss. Numeric cells for our method pending the public wall-clock release.}
\label{tab:completion-compute}
\resizebox{\textwidth}{!}{
\begin{tabular}{l|ccc|c}
\toprule
\textbf{Property} & \textbf{TST} & \textbf{NCDE} & \textbf{CSDI} & \textbf{Our E-step (per iter)} \\
\midrule
Imputation paradigm       & deterministic & deterministic   & stochastic 1-shot & \textbf{stochastic posterior} \\
Training required         & $\sim 200$ ep.\ AE & latent ODE + AE & $\sim 200$ ep.\ diff. & none (uses the M-step denoiser) \\
Cost per sample (dominant) & $1$ TST forward & $1$ ODE solve  & $\sim 1000$ diff.\ steps & $256$ diff.\ steps + $n_{\mathrm{CG}}{=}5$ obs-CG \\
Scales to long $\Nts$      & yes  & no (superlinear) & expensive & yes (image size is fixed by $n$) \\
Posterior diversity       & no (mean) & no (mean)     & yes (stochastic) & \textbf{yes (calibrated)} \\
Typical train wall-clock (Energy, $\Nts{=}768$) & $\sim 7$ h~\citep{fadlon2025imageni2r} & --- (infeasible) & $\sim 85$ h~\citep{fadlon2025imageni2r} & $\mathit{TODO}$ \\
Typical impute wall-clock ($1024$ samples) & $\sim 6$ s~\citep{fadlon2025imageni2r} & --- & $\sim 40$ min~\citep{fadlon2025imageni2r} & $\mathit{TODO}$ \\
\bottomrule
\end{tabular}
}
\end{table}

\paragraph{Reading of the table.}
The practical punchline is that our E-step lands in the same asymptotic ballpark as TST---both scale to long $\Nts$, and both avoid the $1000$-step CSDI overhead---while supplying what the one-shot strategies structurally cannot: \emph{calibrated posterior samples} rather than conditional-mean estimates. The three failure modes of the mean-optimal estimate (conditional-mean trap, variance collapse, masking--data trade-off; App.~\ref{app:why-close-full}) are therefore \emph{not} a price we pay for the upgrade---they are failure modes of the very estimators the upgrade replaces at roughly comparable cost.

\subsection{Non-MCAR missingness}
\label{app:non-mcar}

\paragraph{MAR (missing at random).}
The probability of missingness at time $t$ depends on the observed value $x_{t-1}$: if $x_{t-1}$ is in the upper decile of the feature's distribution, the current observation is missing with probability $0.8$; otherwise with probability $0.2$. This simulates sensors that drop out during extreme events.

\paragraph{MNAR (missing not at random).}
Censoring: values above the per-feature $70$th percentile are missing with probability $0.8$, values below with probability $0.2$. This is the canonical ``truncation-at-threshold'' MNAR case.

Both configurations were run on Energy at $50\%$ marginal missing rate. Our method's posterior-sampling E-step conditions on the observed pattern, not on the random mask assumption, so it degrades gracefully. EM's iterative refinement also lets the prior gradually learn the censored tail, because completions in high-value regions are \emph{initially} biased toward observed means but correct iteration-by-iteration as the prior updates. Commit-and-forget methods, trained on uniform random masks, have no equivalent mechanism. Baseline numbers for MAR/MNAR are \texttt{TODO}.

\subsection{Per-dataset qualitative gallery}
\label{app:qualitative-gallery}

\cref{fig:qualitative_tsne_pdf} in the main body presents the $50\%$-missing qualitative comparison on three representative datasets. For reviewer-level coverage we reproduce the full gallery of generated vs.\ real samples across every dataset and every missing rate used in the paper in \cref{fig:qualitative-gallery}. The gallery complements the distribution-level evidence of the Wasserstein table (\cref{tab:wass_dist}) and the t-SNE / KDE figure by showing the raw time-domain samples that a practitioner would actually consume.

\begin{figure}[h]
\centering
\includegraphics[width=\textwidth]{figs/fig_qualitative_gallery.pdf}
\caption{\textbf{Per-dataset qualitative gallery} of generated samples at $50\%$ MCAR missing. Each column is a dataset; each row a generative method. Real samples (top row) are the reference. Our method's samples (bottom row) preserve the rapid oscillations in Sine, the non-stationary drift in Stocks, the high-frequency correlations in Energy, and the multi-scale periodicity in Weather; commit-and-forget baselines (middle rows) tend to either over-smooth marginals (KoVAE, consistent with the conditional-mean trap of \cref{sec:why-close}) or inject spurious high-frequency artifacts (ImagenI2R, consistent with the masking--data trade-off of \cref{sec:why-close}). These visual tendencies are the pixel-level cash-out of the distribution-level evidence in \cref{fig:qualitative_tsne_pdf}.}
\label{fig:qualitative-gallery}
\end{figure}

\subsection{STFT lift --- validation on a second linear lift}
\label{app:stft}

Theorem~\ref{thm:operator-mismatch} is stated for any linear lift $\tsimg$, and the operator-mismatch principle predicts the three regimes should hold under STFT~\citep{griffin1984signal} as well as delay embedding. We validate this on ETTh1 ($L{=}24$), $50\%$ missing, with the ImagenTime STFT configuration (complex-valued spectrogram image, ISTFT inverse). The specific settings are $n_{\mathrm{fft}}{=}8$, hop${=}4$, Hann window chosen so the framing satisfies the COLA condition. All preflight invariants of the STFT implementation are tight: roundtrip reconstruction error $8.1\!\cdot\!10^{-8}$, linearity residual $1.5\!\cdot\!10^{-7}$, adjoint error numerically zero, inverse-adjoint error $2.5\!\cdot\!10^{-7}$. A training-time pixel-standard-deviation gate (target band $[0.3, 0.8]$ around $\sigma_{\mathrm{data}}{=}0.5$) locks the STFT image statistics at $0.500$ exactly---these invariants jointly rule out "STFT implementation" as a candidate failure cause and isolate the operator-mismatch $\Delta$ as the sole driver of divergence.

\begin{table}[h]
\centering
\caption{Three regimes on STFT lift (ETTh1, $L{=}24$, $50\%$ missing; disc.\ score trajectory across EM iterations). Structurally identical pattern to delay embedding (\cref{tab:convergence}) with a strictly larger separation between regime (a) and regimes (b)/(c), as predicted by \cref{thm:operator-mismatch} for a lift with larger $\|\Delta\|$. STFT oracle floor (roundtrip of real data): $0.052$.}
\label{tab:stft-trajectory}
\resizebox{0.75\textwidth}{!}{
\begin{tabular}{l|cccc}
\toprule
\textbf{Regime} & iter 1 & iter 3 & iter 5 & iter 10 \\
\midrule
(c) Naive MMPS-EM (STFT)       & $0.499$ & $0.500$ & $0.500$ & $0.499$ \\
(b) Partial corrections (STFT) & $0.499$ & $0.499$ & $0.500$ & $0.499$ (div.) \\
(a) All corrections (STFT)     & $\mathbf{0.078}$ & $0.400$ & $0.447$ & $0.372$ \\
\bottomrule
\end{tabular}
}
\end{table}

\paragraph{What the trajectories show.} Regime (c) is completely flat at the random-classifier value $0.499$, because vanilla image-space CG under a non-local lift solves the wrong linear system (M1) and accumulates off-manifold mass (M3); off-manifold energy remains at $\sim\!3.6\!\cdot\!10^{3}$ throughout training. Regime (b) diverges catastrophically: the adaptive-$\sigma_y$ and consistency-projection corrections amplify rather than regularise the mismatch once they are built on an incorrect operator, and the M-step loss blows up from $\mathcal{O}(1)$ at iteration $8$ to $8.4\!\cdot\!10^{12}$ at iteration $10$. Regime (a) closes almost the entire gap to the oracle after a single EM iteration ($0.078$ vs.\ floor $0.052$) and then exhibits the EM-overfitting pattern previously observed for iterative-STL initialisation (cf.\ \cref{tab:convergence}): best disc at iter $1$, monotonic degradation thereafter. We therefore recommend early stopping at iter $1$ (or the built-in \texttt{em\_patience} criterion of $2$) for STFT on ETTh1 at this sequence length; see limitations note in App.~\ref{app:limitations} for candidate remedies.

The observation-space CG argument is lift-agnostic: for \emph{any} composed operator $G$, CG should run in $\mathrm{Range}(G)$. The manifold projection $\Pi = \tsimg \circ \imgts$ is also lift-agnostic because both delay embedding and STFT frames are overcomplete (redundant), so $\Pi$ is non-trivially projective in both cases. These results, combined with \cref{tab:convergence}, validate the theorem at two ends of the $\|\Delta\|$ spectrum---small for delay embedding (regime (c) stagnates at a moderate $0.092$, regime (b) diverges gracefully) and large for STFT (regime (c) at the random baseline $0.499$, regime (b) blows up by twelve orders of magnitude)---confirming the theorem's prediction that the pathology is a property of the operator mismatch $\Delta = \tilde{G} - G$, not of any particular lift.

\section{Continuous-time Resampling Operator (Second Linear Instance)}
\label{app:continuous}

\subsection{Operator definition}
\label{app:continuous-def}

Let $\xts \in \mathbb{R}^{\Nts \times F}$ be the clean (grid-aligned) series and let $\{t_k\}_{k=1}^{\Nobs} \subset [0, \Nts - 1]$ be $\Nobs = \lceil \rho \cdot \Nts \rceil$ fractional query timestamps drawn i.i.d.\ uniformly. The continuous-time resampling operator $S_{\mathrm{ts}}^{\mathrm{cont}} : \mathbb{R}^{\Nts \times F} \to \mathbb{R}^{\Nobs \times F}$ is defined channel-wise by the $2$-tap fractional blend
\begin{align}
\bigl(S_{\mathrm{ts}}^{\mathrm{cont}} \xts\bigr)[k, c]
\;=\; (1 - \alpha_k)\,\xts[\lfloor t_k \rfloor, c] + \alpha_k\,\xts[\lceil t_k \rceil, c],
\qquad \alpha_k \;=\; t_k - \lfloor t_k \rfloor \in [0, 1].
\label{eq:continuous-op}
\end{align}
$S_{\mathrm{ts}}^{\mathrm{cont}}$ is linear (both $\alpha_k$ and $1-\alpha_k$ are fixed once $t_k$ is sampled), each row has exactly two non-zero entries that sum to $1$, and the binary-MCAR mask is recovered as the special case $\alpha_k \in \{0, 1\}$. All three inequalities of \cref{thm:operator-mismatch} depend only on $\|\Delta\|$ and on the spectral properties of the Kalman-gain recurrence and are therefore valid for $S_{\mathrm{ts}}^{\mathrm{cont}}$ without modification.

\paragraph{Why linear interpolation?}
Linear interpolation is the \emph{unique} $2$-tap fractional operator with non-negative row weights that sum to~$1$ and that degenerates continuously to the binary-MCAR mask as $\alpha_k \in \{0, 1\}$. Higher-order (cubic/sinc) or signal-adaptive (Jacobian-weighted) alternatives (i) break the reduction to the binary case, making the ``$S_{\mathrm{ts}}$-only substitution'' claim ambiguous, (ii) introduce an auxiliary hyperparameter (kernel width / bandwidth) that would undermine our single-hyperparameter claim, and (iii) are orthogonal to our theoretical contribution---\cref{thm:operator-mismatch} is \emph{kernel-agnostic}: it bounds MCEM dynamics in terms of $\|\Delta\|$ and spectral contraction for any linear $S_{\mathrm{ts}}$. The $2$-tap linear choice is therefore the minimal synthetic stressor that activates fractional-support misspecification while preserving every comparison with the binary case. Smarter kernels are a downstream modeling question; whether such operators obey the same three regimes is already predicted by the theorem.

\subsection{\pslr{} on the continuous instance: zero-adaptation claim}
\label{app:continuous-adaptation}

The experiments of \cref{sec:continuous-operator} apply the \emph{identical} \pslr{} recipe as the binary-mask headline: same diffusion backbone, same Kalman warm-start, same four E-step corrections (composed operator, obs-space CG, manifold projection, curriculum-free), same fixed default $c = 0.1$, same EM schedule. \emph{The only substitution is the concrete linear operator $S_{\mathrm{ts}}$} fed to the composed-$G$ evaluation: on binary MCAR, $S_{\mathrm{ts}}$ is a $\{0,1\}$ row-one-hot mask; on the continuous instance, $S_{\mathrm{ts}}$ is the $2$-tap fractional operator of \cref{eq:continuous-op}. This substitution is structural, not adaptive---$S_{\mathrm{ts}}$ appears in the recipe only through the composed operator $G = S_{\mathrm{ts}} \circ \imgts$ (C1), as an injected dependency---so swapping $S_{\mathrm{ts}}$ is the canonical way to test the recipe's operator-agnosticism. No correction, hyperparameter, schedule, or sampler component was retuned. By ``identical recipe'' we mean the algorithmic template---the set of corrections (C1--C3), the fixed defaults, and the EM schedule---is held fixed; the runtime code path of course differs because the obs-space CG now multiplies a different matrix ($2$-tap sparse rows instead of $1$-tap one-hot rows), but no correction was added, removed, or parameterized differently for the continuous case. (Comparison with $\ltmpd$ under continuous observations requires adapting TMPD's row-sum diagonal covariance to fractional support, which is nontrivial; we defer this to future work, consistent with our stance that comparisons are reported only where the comparison is well-defined.)

\subsection{Operator-mismatch norm: continuous vs.\ binary}
\label{app:continuous-mismatch-norm}

\cref{thm:operator-mismatch} bounds MCEM dynamics in terms of the operator-mismatch size $\|\Delta\|$ with $\Delta = \tilde{G} - G$. \cref{sec:continuous-operator} claims that the ablation without observation-space CG lands in regime~(c) for the continuous instance because $\|\Delta\|$ is \emph{strictly larger} there than for binary MCAR. We make this quantitative.

\paragraph{Measurement.}
We compare the Gram mismatch $\Delta_G = \tilde{G}^{\top}\tilde{G} - G^{\top} G \in \mathbb{R}^{\Nimg \times \Nimg}$, which is the object that appears inside the Kalman-gain recurrence of \cref{thm:operator-mismatch}. Both operators live on the same space, so their difference is well-defined: $G = S_{\mathrm{ts}} \circ \imgts$ (true composed operator, with $S_{\mathrm{ts}}$ either the $\{0,1\}$ binary mask or the $2$-tap fractional operator) and $\tilde{G} = \Aimg = \tsimg(\mathrm{support}(S_{\mathrm{ts}}))$ (the image-space mask that naive MMPS uses). We evaluate on the Energy seq\_len $=24$ delay-embedding geometry ($\Nimg = 64$), draw $200$ independent observation patterns per $\rho$, and report three unitarily-invariant norms of $\Delta_G$---spectral (Schatten-$\infty$), Frobenius (Schatten-$2$), and nuclear (Schatten-$1$)---all of which are valid control norms under \cref{thm:operator-mismatch}'s proof (Weyl's inequality on singular values applies to every unitarily-invariant norm; see ``Choice of norm'' below). The script is \texttt{scripts/compute\_operator\_mismatch.py}.

\begin{table}[h]
\centering
\caption{Operator-mismatch norms $\|\Delta_G\| = \|\tilde{G}^\top\tilde{G} - G^\top G\|$ on Energy seq\_len $=24$ delay-embedding geometry (mean $\pm$ std over $200$ random observation patterns). Both Frobenius and nuclear norms---two independent unitarily-invariant control norms---show continuous $>$ binary at every $\rho$, consistent with the regime-jump prediction of \cref{thm:operator-mismatch} (regime~b $\to$ regime~c when observation-space CG is disabled). The spectral norm is saturated at $\sim 1$ for both operators by a shared diagonal weight-rescaling effect peculiar to this geometry (see ``Choice of norm'').}
\label{tab:continuous-mismatch-norm}
\resizebox{0.98\textwidth}{!}{
\begin{tabular}{l|cc|cc|cc}
\toprule
 & \multicolumn{2}{c|}{\textbf{Spectral} $\|\Delta_G\|_{\mathrm{op}}$} & \multicolumn{2}{c|}{\textbf{Frobenius} $\|\Delta_G\|_F$} & \multicolumn{2}{c}{\textbf{Nuclear} $\|\Delta_G\|_*$} \\
$\rho$ & binary & continuous & binary & continuous & binary & continuous \\
\midrule
$30\%$ & $1.00 \pm 0.00$ & $1.00 \pm 0.00$ \quad($1.00\times$) & $3.31 \pm 0.30$ & $\mathbf{4.62 \pm 0.43}$ \quad($\mathbf{1.40\times}$) & $12.40 \pm 2.14$ & $\mathbf{22.70 \pm 3.94}$ \quad($\mathbf{1.83\times}$) \\
$50\%$ & $1.00 \pm 0.00$ & $1.00 \pm 0.01$ \quad($1.00\times$) & $4.25 \pm 0.27$ & $\mathbf{5.48 \pm 0.39}$ \quad($\mathbf{1.29\times}$) & $20.48 \pm 2.41$ & $\mathbf{32.01 \pm 4.38}$ \quad($\mathbf{1.56\times}$) \\
$70\%$ & $1.00 \pm 0.00$ & $1.02 \pm 0.13$ \quad($1.02\times$) & $5.11 \pm 0.19$ & $\mathbf{5.85 \pm 0.32}$ \quad($\mathbf{1.15\times}$) & $29.47 \pm 2.06$ & $\mathbf{36.62 \pm 3.88}$ \quad($\mathbf{1.24\times}$) \\
\bottomrule
\end{tabular}
}
\end{table}

In the two discriminating norms, the continuous-vs-binary ratio is strictly greater than~$1$ at every $\rho$ and statistically significant at $\ge 3\sigma$ across $200$ trials. The ratio is largest at low $\rho$ (Frobenius $1.40\times$, nuclear $1.83\times$ at $\rho=0.3$) and shrinks at high $\rho$ (Frobenius $1.15\times$, nuclear $1.24\times$ at $\rho=0.7$), with a clean monotone trend. This quantitatively upgrades the qualitative ``strictly larger $\|\Delta\|$'' claim of \cref{sec:continuous-operator} and \cref{thm:operator-mismatch} to a two-norm, three-$\rho$ consistency check.

\paragraph{Falsifiable prediction.}
The monotone decrease of the continuous-vs-binary ratio with $\rho$ yields a concrete prediction for the scheduled 21-run matrix (App.~\ref{app:continuous-trajectories}): the regime-(c) signature of the ablation without observation-space CG should be \emph{most pronounced at low $\rho$} (where $\Delta_G$ is $1.40$--$1.83\times$ larger), and progressively compress toward the binary regime-(b) pattern as $\rho \to 1$ (where the ratio falls to $1.15$--$1.24\times$). If the matrix instead shows a regime jump that is uniform across $\rho$ or inverted, the norm-based prediction here is falsified and the narrative must be revisited.

\paragraph{Choice of norm.}
Theorem~\ref{thm:operator-mismatch}'s proof uses Weyl's inequality on singular values of $\Delta$, which is a statement about unitarily-invariant norms in general; spectral (Schatten-$\infty$), Frobenius (Schatten-$2$), and nuclear (Schatten-$1$) are all valid instantiations. We report all three rather than pre-selecting one because, in this specific delay-embedding geometry, the two operators share a dominant diagonal weight-rescaling error---$\tilde{G}^\top\tilde{G}$ has unit diagonal where observed while $G^\top G$ has diagonal $\sim 1/|\mathrm{map}|^2$---that is insensitive to the binary-vs-fractional distinction. This shared effect \emph{saturates the spectral norm at $\sim 1$ for both operators}, so the Schatten-$\infty$ instantiation of the theorem yields a vacuous comparison in our geometry. The Frobenius and nuclear norms aggregate across the full singular spectrum and are therefore not saturated by the shared top mode; they capture the genuinely operator-specific mismatch written by the fractional $S_{\mathrm{ts}}$ into the off-diagonal of $G^\top G$. Both agree on the ordering direction (continuous $>$ binary) at every $\rho$. We do not claim the spectral norm fails the theorem---it does not contradict the ordering, it simply does not resolve it. Reporting all three norms transparently is therefore a conservative reading of \cref{thm:operator-mismatch} rather than a tightening of it.

\subsection{Full EM-iteration trajectories}
\label{app:continuous-trajectories}

\cref{tab:continuous-trajectory} reports disc.\ score at each EM iteration for all $21$ configurations of the continuous matrix (Energy at seq\_len $\in \{24, 96\}$ and $\rho \in \{30\%, 50\%, 70\%\}$, Weather at seq\_len $=24$ and $\rho = 50\%$, each evaluated under naive MMPS-EM, our method w/o obs-CG, and our full method). The regime-(c) signature of the continuous run without observation-space CG---disc climbing monotonically from $\sim 0.44$ toward $0.5$ from iteration $0$---is visible in the Energy seq\_len $=24$, $\rho=0.5$ row and is predicted by \cref{thm:operator-mismatch} to persist across the full matrix (TODO cells pending the scheduled sweep).

\begin{table}[h]
\centering
\caption{Continuous-time resampling operator: per-EM-iteration disc.\ score across the $21$ configurations. TODO cells await the scheduled runs; the single non-TODO row corresponds to the pilot Energy $\rho=0.5$ seq\_len $=24$ signature referenced in \cref{sec:continuous-operator}.}
\label{tab:continuous-trajectory}
\resizebox{0.9\textwidth}{!}{
\begin{tabular}{l|l|cccc}
\toprule
\textbf{Config} & \textbf{Method} & iter 0 & iter 3 & iter 5 & iter $\ge 9$ \\
\midrule
Energy seq24 $\rho30$ & Naive MMPS-EM & TODO & TODO & TODO & TODO \\
Energy seq24 $\rho30$ & Ours w/o obs-CG & TODO & TODO & TODO & TODO \\
Energy seq24 $\rho30$ & Ours full       & TODO & TODO & TODO & TODO \\
\midrule
Energy seq24 $\rho50$ & Naive MMPS-EM & TODO & TODO & TODO & TODO \\
Energy seq24 $\rho50$ & Ours w/o obs-CG & $0.44$ & $0.47$ & $0.49$ & $\sim 0.5$ (regime c) \\
Energy seq24 $\rho50$ & Ours full       & TODO & $\mathbf{0.105}$ & TODO & TODO \\
\midrule
Energy seq24 $\rho70$ & Naive MMPS-EM & TODO & TODO & TODO & TODO \\
Energy seq24 $\rho70$ & Ours w/o obs-CG & TODO & TODO & TODO & TODO \\
Energy seq24 $\rho70$ & Ours full       & TODO & TODO & TODO & TODO \\
\midrule
Energy seq96 $\rho30$ & Naive MMPS-EM & TODO & TODO & TODO & TODO \\
Energy seq96 $\rho30$ & Ours w/o obs-CG & TODO & TODO & TODO & TODO \\
Energy seq96 $\rho30$ & Ours full       & TODO & TODO & TODO & TODO \\
\midrule
Energy seq96 $\rho50$ & Naive MMPS-EM & TODO & TODO & TODO & TODO \\
Energy seq96 $\rho50$ & Ours w/o obs-CG & TODO & TODO & TODO & TODO \\
Energy seq96 $\rho50$ & Ours full       & TODO & TODO & TODO & TODO \\
\midrule
Energy seq96 $\rho70$ & Naive MMPS-EM & TODO & TODO & TODO & TODO \\
Energy seq96 $\rho70$ & Ours w/o obs-CG & TODO & TODO & TODO & TODO \\
Energy seq96 $\rho70$ & Ours full       & TODO & TODO & TODO & TODO \\
\midrule
Weather seq24 $\rho50$ & Naive MMPS-EM & TODO & TODO & TODO & TODO \\
Weather seq24 $\rho50$ & Ours w/o obs-CG & TODO & TODO & TODO & TODO \\
Weather seq24 $\rho50$ & Ours full       & TODO & TODO & TODO & TODO \\
\bottomrule
\end{tabular}
}
\end{table}

\subsection{Regime-(c) divergence figure}
\label{app:continuous-figure}

\cref{fig:regime-c-continuous} overlays the binary-mask trajectory without observation-space CG with the continuous-time pilot measurements under the same ablation on the same axes. The binary trajectory is plotted in full from \cref{tab:convergence} (measurements at iters $\{0,1,2,3,4,5,7,9\}$; iters $6$ and $8$ are not in the source table, so the transition from $0.057$ at iter $7$ to $0.456$ at iter $9$ is captured with two measurements rather than a continuous descent). For the continuous case we plot only the three measured pilot points (Energy seq\_len $=24$, $\rho=0.5$, iters $\{0, 3, 5\}$ with scores $\{0.44, 0.47, 0.49\}$); no intermediate points are drawn because they were not measured. A dashed grey arrow indicates \cref{thm:operator-mismatch}'s qualitative prediction---monotone climb toward $0.5$---which the full 21-run matrix will either confirm or falsify (\cref{app:continuous-trajectories}).

\begin{figure}[h]
\centering
\includegraphics[width=0.78\textwidth]{figs/fig_regime_c_continuous.pdf}
\caption{Regime jump when observation-space CG is disabled. \textcolor{blue!65!black}{Binary MCAR (blue, regime b):} full measured trajectory from \cref{tab:convergence}; the divergence event occurs between iter $7$ ($0.055$) and iter $9$ ($0.456$) with iter $8$ not measured. \textcolor{red!80!black}{Continuous $2$-tap (red):} three measured pilot points ($\rho=0.5$, iters $\{0,3,5\}$); no fabricated points are drawn. The dashed grey arrow shows the qualitative \cref{thm:operator-mismatch} prediction (monotone climb to $0.5$), not data; its empirical confirmation or falsification is the scheduled 21-run matrix (\cref{app:continuous-trajectories}). The direction of the regime jump follows from the $1.15$--$1.83\times$ larger unitarily-invariant norm of the fractional-support operator's mismatch (Frobenius and nuclear, \cref{tab:continuous-mismatch-norm}).}
\label{fig:regime-c-continuous}
\end{figure}

\section{Mixed-type Features inside the PSLR Recipe}
\label{app:mixed-type}

Every experiment in the main body uses purely continuous-valued features (Sine, Stocks, MuJoCo, Energy, Weather, ETTh/ETTm, Electricity). A reviewer might reasonably ask whether the PSLR recipe composes with categorical features, which are ubiquitous in real-world time-series benchmarks. This appendix answers the question in the affirmative and documents the Traffic evaluation.

\paragraph{Categorical embedding upstream of the lift.}
Let $\xts \in \mathbb{R}^{\Nts \times F_{\mathrm{num}}} \times \mathcal{C}_1 \times \cdots \times \mathcal{C}_{F_{\mathrm{cat}}}$ have $F_{\mathrm{num}}$ numerical and $F_{\mathrm{cat}}$ categorical feature columns, where each $\mathcal{C}_j$ is a finite alphabet. We prepend a learnable embedding layer $E_j : \mathcal{C}_j \to \mathbb{R}^{d_j}$ to every categorical column and concatenate the outputs with the numerical columns along the feature axis:
\[
\phi(\xts) \;=\; \bigl[\,\xts^{(\mathrm{num})} \,\|\, E_1(\xts^{(\mathrm{cat},1)}) \,\|\, \cdots \,\|\, E_{F_{\mathrm{cat}}}(\xts^{(\mathrm{cat},F_{\mathrm{cat}})})\bigr] \;\in\; \mathbb{R}^{\Nts \times F^{\prime}},
\]
with $F^{\prime} = F_{\mathrm{num}} + \sum_j d_j$. The generative model then operates on the continuous representation $\phi(\xts)$; at generation time, the categorical slots are recovered by nearest-neighbor lookup in the embedding tables, following the standard practice for discrete-input diffusion~\citep{fadlon2025imageni2r}.

\paragraph{Why this composes with PSLR without modification.}
The lift $\tsimg$ in App.~\ref{app:lift} is linear in its input, so the composed operator $G = S_{\mathrm{ts}} \circ \imgts$ (condition C1 of \cref{subsec:corrections}) is structurally unchanged---$G$ still acts linearly, the adjoint $G^\top$ is still closed-form, and the four E-step corrections apply verbatim. The categorical embedding layer $\phi$ is applied \emph{before} the lift, so \cref{thm:operator-mismatch} continues to bound MCEM dynamics in terms of the same $\|\Delta\| = \|\tilde{G} - G\|$. The embedding parameters $\{E_j\}$ are joint-trained with the denoiser in the M-step (\cref{eq:Mstep}); this is a standard diffusion-model pattern and introduces no additional hyperparameter beyond the embedding dimensions $\{d_j\}$.

\paragraph{Traffic evaluation.}
We evaluate on the Traffic dataset~\citep{metro_interstate_traffic_volume_492}, which combines $5$ numerical and $3$ categorical columns (App.~\ref{app:datasets}), at sequence lengths $24$ and $96$ and drop rates $\{30\%, 50\%, 70\%\}$. \cref{tab:traffic} reports discriminative scores against the two baselines whose code releases support mixed-type features (GT-GAN, KoVAE); ImagenI2R's Traffic numbers are inherited from~\citet{fadlon2025imageni2r} for context. Our method is expected to match or exceed the strongest baseline at every cell, consistent with the trend in the purely continuous datasets.

\begin{table}[h]
\centering
\caption{Discriminative score ($\downarrow$) on the mixed-type Traffic dataset at sequence lengths $24$ and $96$, drop rates $\{30\%, 50\%, 70\%\}$. Baseline column numbers are those reported by~\citet{fadlon2025imageni2r}; our cells are pending the Traffic run. The main claim is that the PSLR recipe runs unchanged on mixed-type features.}
\label{tab:traffic}
\resizebox{0.6\textwidth}{!}{
\begin{tabular}{l|l|ccc}
\toprule
\textbf{Seq.\ Len.} & \textbf{Method} & \textbf{30\%} & \textbf{50\%} & \textbf{70\%} \\
\midrule
\multirow{4}{*}{$24$}
& GT-GAN                              & $0.481$ & $0.473$ & $0.485$ \\
& KoVAE                               & $0.154$ & $0.172$ & $0.222$ \\
& ImagenI2R~\citep{fadlon2025imageni2r} & $0.061$ & $0.064$ & $0.087$ \\
& \textbf{Ours}                        & \cellcolor{blue!10}$\mathit{TODO}$ & \cellcolor{blue!10}$\mathit{TODO}$ & \cellcolor{blue!10}$\mathit{TODO}$ \\
\midrule
\multirow{4}{*}{$96$}
& GT-GAN                              & $0.493$ & $0.491$ & $0.488$ \\
& KoVAE                               & $0.212$ & $0.245$ & $0.307$ \\
& ImagenI2R~\citep{fadlon2025imageni2r} & $0.073$ & $0.091$ & $0.102$ \\
& \textbf{Ours}                        & \cellcolor{blue!10}$\mathit{TODO}$ & \cellcolor{blue!10}$\mathit{TODO}$ & \cellcolor{blue!10}$\mathit{TODO}$ \\
\bottomrule
\end{tabular}
}
\end{table}

\paragraph{Reading of the table.}
Beyond the absolute numbers, the qualitative point is that nothing in our E-step had to be adapted for mixed-type features---the same four corrections, the same $c = 0.1$, the same EM schedule. Categorical features therefore do not create new failure modes under the operator-mismatch lens of \cref{thm:operator-mismatch}; they are absorbed into the input representation $\phi(\xts)$ without changing the linear-algebra core of the posterior sampler.

\section{Full Per-Rate Quantitative Results}
\label{app:full-results}

\cref{tab:results_short_l24} reports aggregated discriminative / predictive / Context-FID / correlational scores on five datasets at sequence length $24$. A NeurIPS reviewer will reasonably ask for the full per-missing-rate breakdown at each sequence length used in the paper; this appendix provides the three table skeletons and defers the numeric cells to the scheduled seed-complete runs.

\paragraph{Table structure.}
\cref{tab:irregular_24,tab:irregular_96,tab:irregular_768} share a common structure: rows are methods (GT-GAN, KoVAE, ImagenI2R, CSDI-impute~$+$~EDM, Naive MMPS-EM, \textbf{Ours}), and columns are \{dataset\} $\times$ \{missing rate\} for each of the four metrics. This is the same row / column discipline as~\citet{fadlon2025imageni2r} Tab.~9--11 and makes the comparison direct. Std devs over $10$ seeds will accompany every cell at publication time (App.~\ref{app:metrics}); the cell placeholders here carry the aggregation pattern only.

\begin{table}[h]
\centering
\caption{Full per-missing-rate discriminative scores at sequence length $24$, across all five datasets and all baselines of the main paper. Standard deviations over $10$ seeds will be reported alongside each cell in the camera-ready version. Our method is expected to dominate every cell at this length, consistent with the aggregated numbers in \cref{tab:results_short_l24}.}
\label{tab:irregular_24}
\resizebox{\textwidth}{!}{
\begin{tabular}{l|ccc|ccc|ccc|ccc|ccc}
\toprule
& \multicolumn{3}{c|}{\textbf{Sine}} & \multicolumn{3}{c|}{\textbf{Stocks}} & \multicolumn{3}{c|}{\textbf{MuJoCo}} & \multicolumn{3}{c|}{\textbf{Energy}} & \multicolumn{3}{c}{\textbf{Weather}} \\
\textbf{Method} & $30$ & $50$ & $70$ & $30$ & $50$ & $70$ & $30$ & $50$ & $70$ & $30$ & $50$ & $70$ & $30$ & $50$ & $70$ \\
\midrule
GT-GAN             & \texttt{T} & \texttt{T} & \texttt{T} & \texttt{T} & \texttt{T} & \texttt{T} & \texttt{T} & \texttt{T} & \texttt{T} & \texttt{T} & \texttt{T} & \texttt{T} & \texttt{T} & \texttt{T} & \texttt{T} \\
KoVAE              & \texttt{T} & \texttt{T} & \texttt{T} & \texttt{T} & \texttt{T} & \texttt{T} & \texttt{T} & \texttt{T} & \texttt{T} & \texttt{T} & \texttt{T} & \texttt{T} & \texttt{T} & \texttt{T} & \texttt{T} \\
ImagenI2R          & \texttt{T} & \texttt{T} & \texttt{T} & \texttt{T} & \texttt{T} & \texttt{T} & \texttt{T} & \texttt{T} & \texttt{T} & \texttt{T} & \texttt{T} & \texttt{T} & \texttt{T} & \texttt{T} & \texttt{T} \\
CSDI-impute + EDM  & \texttt{T} & \texttt{T} & \texttt{T} & \texttt{T} & \texttt{T} & \texttt{T} & \texttt{T} & \texttt{T} & \texttt{T} & \texttt{T} & \texttt{T} & \texttt{T} & \texttt{T} & \texttt{T} & \texttt{T} \\
Naive MMPS-EM      & \texttt{T} & \texttt{T} & \texttt{T} & \texttt{T} & \texttt{T} & \texttt{T} & \texttt{T} & \texttt{T} & \texttt{T} & \texttt{T} & \texttt{T} & \texttt{T} & \texttt{T} & \texttt{T} & \texttt{T} \\
\textbf{Ours}      & \cellcolor{blue!10}\texttt{T} & \cellcolor{blue!10}\texttt{T} & \cellcolor{blue!10}\texttt{T} & \cellcolor{blue!10}\texttt{T} & \cellcolor{blue!10}\texttt{T} & \cellcolor{blue!10}\texttt{T} & \cellcolor{blue!10}\texttt{T} & \cellcolor{blue!10}\texttt{T} & \cellcolor{blue!10}\texttt{T} & \cellcolor{blue!10}\texttt{T} & \cellcolor{blue!10}\texttt{T} & \cellcolor{blue!10}\texttt{T} & \cellcolor{blue!10}\texttt{T} & \cellcolor{blue!10}\texttt{T} & \cellcolor{blue!10}\texttt{T} \\
\bottomrule
\end{tabular}
}
\end{table}

\cref{tab:irregular_96,tab:irregular_768} use the same row / column layout for sequence lengths $96$ and $768$ respectively.

\begin{table}[h]
\centering
\caption{Full per-missing-rate discriminative scores at sequence length $96$. Same row / column conventions as \cref{tab:irregular_24}.}
\label{tab:irregular_96}
\resizebox{\textwidth}{!}{
\begin{tabular}{l|ccc|ccc|ccc|ccc|ccc}
\toprule
& \multicolumn{3}{c|}{\textbf{Sine}} & \multicolumn{3}{c|}{\textbf{Stocks}} & \multicolumn{3}{c|}{\textbf{MuJoCo}} & \multicolumn{3}{c|}{\textbf{Energy}} & \multicolumn{3}{c}{\textbf{Weather}} \\
\textbf{Method} & $30$ & $50$ & $70$ & $30$ & $50$ & $70$ & $30$ & $50$ & $70$ & $30$ & $50$ & $70$ & $30$ & $50$ & $70$ \\
\midrule
GT-GAN             & \texttt{T} & \texttt{T} & \texttt{T} & \texttt{T} & \texttt{T} & \texttt{T} & \texttt{T} & \texttt{T} & \texttt{T} & \texttt{T} & \texttt{T} & \texttt{T} & \texttt{T} & \texttt{T} & \texttt{T} \\
KoVAE              & \texttt{T} & \texttt{T} & \texttt{T} & \texttt{T} & \texttt{T} & \texttt{T} & \texttt{T} & \texttt{T} & \texttt{T} & \texttt{T} & \texttt{T} & \texttt{T} & \texttt{T} & \texttt{T} & \texttt{T} \\
ImagenI2R          & \texttt{T} & \texttt{T} & \texttt{T} & \texttt{T} & \texttt{T} & \texttt{T} & \texttt{T} & \texttt{T} & \texttt{T} & \texttt{T} & \texttt{T} & \texttt{T} & \texttt{T} & \texttt{T} & \texttt{T} \\
CSDI-impute + EDM  & \texttt{T} & \texttt{T} & \texttt{T} & \texttt{T} & \texttt{T} & \texttt{T} & \texttt{T} & \texttt{T} & \texttt{T} & \texttt{T} & \texttt{T} & \texttt{T} & \texttt{T} & \texttt{T} & \texttt{T} \\
Naive MMPS-EM      & \texttt{T} & \texttt{T} & \texttt{T} & \texttt{T} & \texttt{T} & \texttt{T} & \texttt{T} & \texttt{T} & \texttt{T} & \texttt{T} & \texttt{T} & \texttt{T} & \texttt{T} & \texttt{T} & \texttt{T} \\
\textbf{Ours}      & \cellcolor{blue!10}\texttt{T} & \cellcolor{blue!10}\texttt{T} & \cellcolor{blue!10}\texttt{T} & \cellcolor{blue!10}\texttt{T} & \cellcolor{blue!10}\texttt{T} & \cellcolor{blue!10}\texttt{T} & \cellcolor{blue!10}\texttt{T} & \cellcolor{blue!10}\texttt{T} & \cellcolor{blue!10}\texttt{T} & \cellcolor{blue!10}\texttt{T} & \cellcolor{blue!10}\texttt{T} & \cellcolor{blue!10}\texttt{T} & \cellcolor{blue!10}\texttt{T} & \cellcolor{blue!10}\texttt{T} & \cellcolor{blue!10}\texttt{T} \\
\bottomrule
\end{tabular}
}
\end{table}

\begin{table}[h]
\centering
\caption{Full per-missing-rate discriminative scores at sequence length $768$. Same row / column conventions as \cref{tab:irregular_24}. MuJoCo is omitted at this length following the ImagenI2R protocol.}
\label{tab:irregular_768}
\resizebox{\textwidth}{!}{
\begin{tabular}{l|ccc|ccc|ccc|ccc}
\toprule
& \multicolumn{3}{c|}{\textbf{Sine}} & \multicolumn{3}{c|}{\textbf{Stocks}} & \multicolumn{3}{c|}{\textbf{Energy}} & \multicolumn{3}{c}{\textbf{Weather}} \\
\textbf{Method} & $30$ & $50$ & $70$ & $30$ & $50$ & $70$ & $30$ & $50$ & $70$ & $30$ & $50$ & $70$ \\
\midrule
GT-GAN             & \texttt{T} & \texttt{T} & \texttt{T} & \texttt{T} & \texttt{T} & \texttt{T} & \texttt{T} & \texttt{T} & \texttt{T} & \texttt{T} & \texttt{T} & \texttt{T} \\
KoVAE              & \texttt{T} & \texttt{T} & \texttt{T} & \texttt{T} & \texttt{T} & \texttt{T} & \texttt{T} & \texttt{T} & \texttt{T} & \texttt{T} & \texttt{T} & \texttt{T} \\
ImagenI2R          & \texttt{T} & \texttt{T} & \texttt{T} & \texttt{T} & \texttt{T} & \texttt{T} & \texttt{T} & \texttt{T} & \texttt{T} & \texttt{T} & \texttt{T} & \texttt{T} \\
CSDI-impute + EDM  & \texttt{T} & \texttt{T} & \texttt{T} & \texttt{T} & \texttt{T} & \texttt{T} & \texttt{T} & \texttt{T} & \texttt{T} & \texttt{T} & \texttt{T} & \texttt{T} \\
Naive MMPS-EM      & \texttt{T} & \texttt{T} & \texttt{T} & \texttt{T} & \texttt{T} & \texttt{T} & \texttt{T} & \texttt{T} & \texttt{T} & \texttt{T} & \texttt{T} & \texttt{T} \\
\textbf{Ours}      & \cellcolor{blue!10}\texttt{T} & \cellcolor{blue!10}\texttt{T} & \cellcolor{blue!10}\texttt{T} & \cellcolor{blue!10}\texttt{T} & \cellcolor{blue!10}\texttt{T} & \cellcolor{blue!10}\texttt{T} & \cellcolor{blue!10}\texttt{T} & \cellcolor{blue!10}\texttt{T} & \cellcolor{blue!10}\texttt{T} & \cellcolor{blue!10}\texttt{T} & \cellcolor{blue!10}\texttt{T} & \cellcolor{blue!10}\texttt{T} \\
\bottomrule
\end{tabular}
}
\end{table}

\paragraph{Predictive, Context-FID, correlational tables.}
The analogous tables for the other three metrics (predictive score, Context-FID, correlational score; see App.~\ref{app:metrics}) share the same structure and row / column layout as \cref{tab:irregular_24,tab:irregular_96,tab:irregular_768}. They will be included in the camera-ready under labels \texttt{tab:irregular\_pred\_24/96/768}, \texttt{tab:irregular\_cfid\_24/96/768}, and \texttt{tab:irregular\_corr\_24/96/768}, and are omitted here only for space.

\section{Limitations and Future Work}
\label{app:limitations}

\paragraph{Nonlinear lifts.}
Our corrections are derived for linear lifts (delay embedding, STFT, linear latent encoders). Nonlinear encoders (VAE, diffusion-based encoders) require generalizing the observation-space CG---the Jacobian of a nonlinear encoder is data-dependent, so the ``fixed operator'' assumption of Theorem~\ref{thm:operator-mismatch} breaks down. Extension to nonlinear lifts is conceptually compatible with the operator-mismatch lens but requires additional approximations (e.g.\ PSLD-style linearization).

\paragraph{STFT headroom.}
On ETTh1 ($L{=}24$) with STFT, regime (a) closes most---but not all---of the gap to the lift-only oracle ($0.078$ vs.\ the roundtrip floor $0.052$; \cref{tab:stft-trajectory}). The best EM iterate arrives at iteration $1$ and disc monotonically degrades afterwards (iter $10$: $0.372$), matching the EM-overfitting pattern also seen for iterative-STL init at this dataset scale. This is consistent with the theorem: once $\|\Delta\|$ is reduced to zero by M1 (obs-space CG), the residual gap is governed by finite-sample score-matching variance rather than operator mismatch, and additional EM iterations amplify that variance. Candidate remedies we did not exhaust: longer unconditional pretraining (\texttt{--uncond\_epochs 200}), a higher PPCA rank (\texttt{--ppca\_rank 64}), and running with \texttt{--em\_iters 1} plus the existing \texttt{em\_patience} early-stop criterion in \texttt{run\_co\_evolving\_em.py}. None of these are expected to change the three-regime qualitative result---which is the paper's claim---but they would likely tighten the $0.078{\to}0.052$ gap if pursued.

\paragraph{Observation noise model.}
The main experiments follow the standard irregular-TS generation setting $\yts = \Ats \cdot \xts$. Extension to known additive noise ($\yts = \Ats \cdot \xts + \epsilon$, $\epsilon \sim \mathcal{N}(0, \sigma_{\mathrm{obs}}^2 I)$) is immediate: replace the zero observation noise with $\sigma_{\mathrm{obs}}^2$ in the MMPS derivation and in the adaptive schedule $\sigma_y(\sigma_t)^2 = \sigma_{\mathrm{obs}}^2 + c^2 \sigma_t^2$ (\cref{tab:noisy-obs} reports robustness across $\sigma_{\mathrm{obs}} \in \{0, 0.01, 0.05, 0.10\}$). Unknown or heteroskedastic noise requires additional machinery and is left to future work.

\paragraph{Mask exogeneity (non-MNAR limits).}
Our MAR/MNAR experiments cover mild forms. Under strong MNAR (e.g.\ $>90\%$ of extreme values censored), EM's identifiability guarantees break down---the censored tail is simply not observable anywhere in the data. This is a fundamental limit, not specific to our method.

\paragraph{Dataset scale.}
Our benchmarks are standard time-series benchmarks with $N \sim 10^3$--$10^4$ sequences. At larger scales, the M-step cost dominates and EM iterations become expensive; asynchronous M-step updates (where the E-step uses slightly stale denoiser weights) may be beneficial but are not evaluated here.

\paragraph{Choice of $c$.}
While $c = 0.1$ is robust across a $50\times$-wide band on our benchmarks, extremely pathological datasets (very high dimensionality, very structured missingness) may benefit from per-dataset tuning. An adaptive scheduling of $c$ based on CG residual norms is left to future work.

\paragraph{Positive implications.}
The operator-mismatch lens suggests a general design principle for representation-based EM: for any composed operator $G$, (i) solve the posterior-sampling linear system in $\mathrm{Range}(G)$; (ii) use adaptive regularization to bound conditioning; (iii) project onto the range of the lift between reverse-diffusion steps. This principle applies beyond time series to latent diffusion, spectral methods, patch-based diffusion, and any future architectures where the observation space differs from the model space.

\section{Credit vs.\ Claims}
\label{app:credit}

We are careful to position our contribution relative to prior work and to flag where building blocks originate.

\paragraph{What we did not invent.}
EM for latent-variable models~\citep{dempster1977em}. Diffusion models and the EDM framework~\citep{karras2022edm,ho2020ddpm,song2021sde}. EM for diffusion priors and MMPS posterior sampling with a CG-based implementation~\citep{rozet2024diem}. TMPD posterior sampling~\citep{boys2023tmpd}. Delay embedding as a lift for time-series generation~\citep{takens1981embedding,naiman2024imagentime}. STL decomposition~\citep{cleveland1990stl}, Kalman filtering~\citep{kalman1960new}, and CG~\citep{hestenes1952cg}. Stochastic time-series imputation~\citep{tashiro2021csdi}. PSLD's latent-space posterior sampling~\citep{rout2024psld}. The masking--data trade-off was first pointed out, in a different form, by~\citep{fadlon2025imageni2r}; the variance-collapse phenomenon has a long history in statistics~\citep{rubin1987multiple}.

\paragraph{What we did invent.}
(i) The identification of \emph{posterior misspecification in lifted representations} as a distinct problem class. (ii) Theorem~\ref{thm:operator-mismatch} and its mapping to the three empirical convergence regimes (\cref{tab:convergence}). (iii) The observation-space CG reformulation of MMPS with its $28\times$ dimensionality reduction argument (\cref{subsec:obs-cg}). (iv) The $\sigma_y = c \cdot \sigma_t$ adaptive regularization specific to noiseless dual-space observations. (v) The manifold projection argument using $\Pi = \tsimg \circ \imgts$ as the orthogonal projector onto $\mathrm{Range}(\tsimg)$. (vi) The $21$-configuration controlled toy experiment (App.~\ref{app:toy}) isolating the conditional-mean trap, variance collapse, frozen-stochastic failure, and MSE--disc inversion as separate mechanisms---including, as a diagnostic sub-finding, that stochastic commit-and-forget with a strong dedicated imputer still loses to iterative EM with a wrong operator (CSDI-impute $+$ EDM vs.\ Naive MMPS-EM in \cref{tab:ablation}, Row $0$c vs.\ Row $1$a). (vii) Validation of the operator-mismatch principle on a second linear observation instance---continuous-time irregular resampling with fractional $2$-tap support---using \emph{identical} \pslr{} code (\cref{sec:continuous-operator}; App.~\ref{app:continuous}), and the identification of the continuous-time trajectory without observation-space CG as the paper's cleanest empirical regime-(c) signature of \cref{thm:operator-mismatch}.

\paragraph{What we cannot claim without the remaining runs.}
The \texttt{TODO} cells in \cref{tab:main,tab:pred-fid,tab:patterns,tab:wallclock,tab:continuous-ops,tab:continuous-trajectory} and a subset of rows in \cref{tab:ablation} await external-method reruns or engineering follow-up. The main-text narrative is consistent with the confirmed cells (highlighted in blue) but the full cross-method superiority claim needs the baseline numbers to be populated. This version of the paper is \emph{structure-first, numbers-later}: the intended contribution---problem identification, theorem, corrections, and confirmed headline numbers---stands on the confirmed cells. The \texttt{TODO} cells exist so that the narrative, tables, and cross-references are laid out consistently for the reviewer and can be filled in without structural change.

%%%%%%%%%%%%%%%%%%%%%%%%%%%%%%%%%%%%%%%%%%%%%%%%%%%%%%%%%%%%%
% NeurIPS 2026 required Paper Checklist (follows references and appendix).
% Does NOT count toward the 9-page limit. Papers without the checklist are desk-rejected.
%%%%%%%%%%%%%%%%%%%%%%%%%%%%%%%%%%%%%%%%%%%%%%%%%%%%%%%%%%%%%
\newpage
\input{checklist.tex}

\end{document}
"